\documentclass[11pt,a4paper]{article}
\usepackage[T1]{fontenc}
\usepackage[utf8]{inputenc}
\usepackage{lmodern,amsmath,amssymb,amsthm,mathtools,microtype,mathrsfs}
\usepackage[margin=25mm,top=23mm,bottom=25mm]{geometry}
\usepackage{booktabs,array,longtable,enumitem,xcolor,fancyhdr}
\usepackage[colorlinks=true,linkcolor=blue!45!black,urlcolor=blue!45!black,citecolor=blue!45!black]{hyperref}
\hypersetup{pdftitle={Reading the Alpoge--Claude RH and Dirichlet-GRH proof: working version 0.45},pdfauthor={ChatGPT},pdfsubject={Partial human-readable reconstruction with explicit arithmetic inputs},pdfkeywords={Riemann hypothesis, Lean, theta lifts, prime sums, Arakelov slopes}}
\setlength{\parindent}{0pt}
\setlength{\parskip}{5pt plus 1pt}
\setlist{topsep=3pt,itemsep=3pt,leftmargin=*}
\setlength{\headheight}{14pt}
\pagestyle{fancy}\fancyhf{}
\fancyhead[L]{\small RH and Dirichlet GRH: a reconstruction}
\fancyhead[R]{\small Working version 0.45}
\fancyfoot[C]{\thepage}
\renewcommand{\headrulewidth}{0.3pt}
\numberwithin{equation}{section}
\newtheorem{theorem}{Theorem}[section]
\newtheorem{proposition}[theorem]{Proposition}
\newtheorem{lemma}[theorem]{Lemma}
\newtheorem{corollary}[theorem]{Corollary}
\theoremstyle{remark}\newtheorem{remark}[theorem]{Remark}
\newcommand{\Rea}{\operatorname{Re}}
\newcommand{\HS}{\mathrm{HS}}
\newcommand{\tr}{\operatorname{tr}}
\newcommand{\codim}{\operatorname{codim}}
\newcommand{\covol}{\operatorname{covol}}
\newcommand{\OK}{\mathcal O_K}
\newcommand{\src}[2]{\href{https://github.com/CocoSgt/riemann-hypothesis-lean/blob/7cbf7a0c32fca1badf05931879a37c3f8a2b03b4/src/#1}{\texttt{#2}}}
\newcommand{\status}[1]{\textbf{Status:} #1}
\begin{document}
\begin{center}
{\LARGE\bfseries RH and Dirichlet GRH: new reconstruction chapters}\\[7pt]
{\large Version 0.45 --- 23 November 2026}\\[5pt]
{\large The six Hida fields and both support installations}
\end{center}
\textbf{Scope.} This edition contains the current audit and the seven new
chapters. It remains a partial human reconstruction. Earlier named
claims are in the cumulative manuscript and preserved component files.
The underlying formal proof is accepted as correct; the exact remaining
human reconstruction boundaries are stated below.
\setcounter{tocdepth}{1}
\tableofcontents
\clearpage
\begingroup\emergencystretch=3em
\section{Current dependency audit: the six Hida fields and both support installations}
\label{sec:dependency-audit-v45}

\subsection{Current status and reading order}

Version 0.45 completes the six fields of the original
\texttt{Hida\allowbreak Props\allowbreak Alt} on one ideal $J$, constructs the actual
all-level \texttt{Level\allowbreak Alt} and its \texttt{Has\allowbreak Aut}, and installs
the selected rank-four Hida/$P_1$ support result on both original
Links above one common coefficient threshold.  The entire selected
\texttt{COUNT\_\allowbreak II\_\allowbreak SUPPORT} obligation is consequently discharged.
Its original field, embedding, residual, determinant and
pro-$\ell$ data are retained, including $\ell=3$.

The pinned Alpoge--Claude formal proof is accepted as fully checked
and correct.  This remains a \textbf{partial human reconstruction}
of RH and unrestricted Dirichlet GRH.  The common geometric and
arithmetic families, global estimates and several case producers
still have open human reconstruction contracts.  The all-rank
\texttt{Hida\allowbreak UAlt} conclusion is distinguished from full
\texttt{Mod\allowbreak UAlt}: the latter still requires the original
\texttt{Pfrom\allowbreak Aut} supplier, including arbitrary-chain $P_3$.
No Lean build or kernel replay is claimed.

There are 117 new and 2677 cumulative claims.  All 2560
predecessor claim records, source arrays and dependency edges are
preserved.  The selected sufficient-route graph still has 121 nodes
and 161 edges.  The whole-obligation register decreases from 31 to
30: 21 class-C and 9 class-B tasks.  Of these, 23 are shared by RH
and Dirichlet GRH, and 7 are additional for Dirichlet GRH.  The
claim inventory is not a percentage of proof completion.

HED and HAW prove arithmetic-point decomposition and automorphy
away from the exceptional set.  HPT returns an arbitrary supplied
automorphic point to the actual Hecke support; HTS proves twisting
on that same support.  HLM constructs the complete family of
auxiliary-level modules, and HTW constructs its integral
Taylor--Wiles data.  HIA assembles all six Hida fields and performs
the two selected support installations.  The separate reader
contains this audit and these seven chapters.  The cumulative
edition preserves the previous proof body verbatim.  CMP1--CMP6
remain the complete conditional terminal deduction on their
original hypotheses.

\subsection{The same ideal and all six Hida properties}

The empty auxiliary level gives an actual ordinary Hecke module
and its original \texttt{Is\allowbreak J} ideal.  Every subsequent field refers
to that ideal, the same coefficient ring and the same universal
determinant.  The nonzero module branch uses the kernel of the
actual Hecke-algebra map; the zero branch retains the top ideal.

\begin{center}
\begin{tabular}{@{}p{0.13\textwidth}p{0.53\textwidth}p{0.22\textwidth}@{}}
\toprule
Original field & Human-readable output & Completed supplier\\
\midrule
\texttt{ring} & All original quotient-ring and determinant properties & JRP15\\
\texttt{dense} & Bare arithmetic-point density for this same $J$ & UPA18\\
\texttt{pt} & The original point theorem, with its rank-$2$--$4$ antecedent & HPT18\\
\texttt{end1} & Full arithmetic-point decomposition & HED16\\
\texttt{end2} & Full irreducible arithmetic-point automorphy away & HAW19\\
\texttt{twist} & The original all-positive-rank twisting property & HTS8\\
\bottomrule
\end{tabular}
\end{center}

HED starts with the given arithmetic point and its given continuous
semisimple representation.  Surjective quotient factorization
preserves the point, weight and full polynomial law after every
scalar extension.  The base-change decomposition supplies one
initial Galois representation for every cusp and every positive
Speh length.  Exact inverse cyclotomic strings retain the
Frobenius normalization.  Degree-one Chebotarev turns the cofinite
Frobenius product into an identity of characteristic polynomials
at every group element.

When all string lengths are one, the already matched
generic-unitary classification gives the strict spherical interval
\[
                   q^{-1/2}<|\alpha|<q^{1/2}.
\]
It yields the original cofinite strict ratio
$|z|<q|z'|$ for every supplied Frobenius multiset at every required
place.  This branch does not discard the other positive lengths
from the full decomposition theorem.

HAW proves the original ordinary principal-series description at
every arithmetic point at its literal scope.  A nonzero Jacquet
character block has a character quotient by finite-length
Nakayama; a functional need not detect a preselected eigenvector.
The ordinary projector selects coefficient-field units, including
the exclusion of every nonunit.  The ordered flag and conjugate
place comparison retain the polarization multiplier
$\epsilon^{1-n}$ and the reversal of dual characters.

One global real norm twist unitarizes the central character before
any finite place is selected.  It therefore gives one common center
for all the required local exponent windows.  The exact local
character calculation then shows that a forbidden positive-power
cyclotomic ratio would force a factor $q$ between two character
norms.  The strict window excludes this.  The resulting
\texttt{Aut\allowbreak Off\allowbreak N} keeps the source's conditional continuity in
its original location.

HPT begins with the supplied automorphic representation and
constructs a nonzero vector in the actual ordinary space, with
the integral type, infinity weight, depth and residual block all
retained.  Its Hecke point equals the supplied point because the
full determinant and ordered characters agree in a finite
coefficient ring, where the universal-ring uniqueness theorem
applies.  Thus $J\subseteq\ker x$ is the conclusion.  The
all-positive-rank return theorem is separate from the last point
field, whose original antecedent includes $2\le n\le4$ and
whose use of \texttt{Desc\allowbreak Weak} remains explicit.

HTS realizes a permitted odd-order conjugate-dual twist by an
actual adelic unitary determinant character.  If its order is
$m=2N+1$, the power $\mu^{N+1}$ is a square root because
$2(N+1)\equiv1\pmod m$.  At a split place the resulting
determinant character has precisely the required square.
Multiplication of the actual automorphic functions preserves the
infinity weight and the required local vectors.  The guarded
local comparisons and full polynomial-law twist supply every
input of the all-rank point-return theorem.  This proves
\texttt{Twist\allowbreak Ok} without importing the low-rank descent premise.

\subsection{One family over every auxiliary level}

HLM constructs the actual directed union of smooth automorphic
modules and identifies its fixed modules at every permitted
level.  Integral trace projectivity supplies a group-algebra
splitting without dividing by the quotient-group order.  The
ordinary and residual projectors preserve the required module
properties, with the residual decomposition justified by
Artin--Rees.

At every level the determinant law is constructed before using
any point theorem which assumes it.  Compact interpolation is
performed over the coefficient ring; integral word identities
are tested in finite-degree truncated scalar algebras.  Polynomial
coefficient separation establishes the complete law after
arbitrary scalar extension.  The full ordinary pair then gives
the surjective universal-ring map.  Universal uniqueness glues
these maps into one coherent action on the directed union.

Classical points separate actual module vectors.  The type and
local correspondence comparisons retain all of the original
quantifiers.  Dense Frobenius comparison identifies the old-level
residual component with the initial Hecke module, preserving
$J$.  One base is chosen before all auxiliary level data,
including the empty auxiliary set.  This proves the literal
\texttt{Level\allowbreak Alt} in every positive rank.

\subsection{Integral Taylor--Wiles comparison and completion}

HTW retains the diagonal-unit terms in the tame Hecke relation.
For a simple reflection it has the form
\[
 T_s^2=Q+T_s\sum_c D_c,\qquad T_sD_c=D_c^sT_s.
\]
The corresponding left and right inverses have their respective
unit terms in the correct order.  The usual Iwahori relation is
the specialization in which those diagonal terms act trivially.
The resulting torus operators and an explicit Jacquet dimension
bound give one exponent $n!$, independent of the later auxiliary
choices and of isotypic multiplicity.

The comparison of spherical and Iwahori modules is integral.
For ordered variables put
\[
 d=\prod_{i<k}(X_k-X_i),\qquad
 N_Q=\prod_{i<k}(X_k-QX_i).
\]
An explicit finite family satisfies
\[
              \sum_T a_T\,w(b_T)=\delta_{w,1}\,d.
\]
With the divided difference $\partial=\mathrm{Alt}/d$, these
identities give both dual expansions over the symmetric ring.
The actual Hecke trace operator is
$E(f)=\partial(N_Qf)$; its integral identity is proved through
the faithful polynomial representation.  These formulas construct
both comparison maps.  Multiplying by the remaining prescribed
factors makes both compositions exactly the original ordered-pair
\texttt{f\allowbreak El}.  No inversion of $n!$ or replacement by an
unspecified annihilator is made.

The desired topologically nilpotent part is a canonical summand
over the finite coefficient algebra.  Its projectivity makes it
finite free over the local diamond group algebra.  Power actions
extend coherently through the inverse limit.  Reduction by the
diamond ideal gives the correct ordered-root summand, and
completion preserves the two comparison maps.  The zero-module
case is handled separately when using the finite image algebra.

The coefficient map is also explicit: the root variables map
to the selected roots minus their prechosen residue lifts, and
the additional variables map to zero.  Its composition with the
coefficient embedding equals the original prescribed map, not
merely a map agreeing on closed points.  Every field of
\texttt{TWRing} and \texttt{TWMaps}, both full trace relations,
and \texttt{Has\allowbreak Aut} is obtained.  One base and one lift of each
residual root are fixed before all later $N$.

\subsection{Both rank-four support installations}

HIA uses the empty-level \texttt{Is\allowbreak J} once, attaches the six
fields listed above to it, and attaches the same $J$'s
\texttt{Has\allowbreak Aut}.  The inherited selected patching, normalization,
pro-$\ell$ and $P_1$ consumers then yield the required rank-four
support condition.  Their original constants and choice order
are preserved; no arbitrary-choice or arbitrary-chain
strengthening is used.

The original two QX Links are installed with fixed $\iota$ above
one common coefficient threshold, before the later coefficient
DVR, universal ring and weight algebra are quantified.  Full
\texttt{Coh\allowbreak P}, \texttt{Loc\allowbreak CFT}, the original residual
cyclotomic condition, determinant relation and pro-$\ell$
conditions remain visible.  This is the precise selected closure
recorded by \texttt{COUNT\_\allowbreak II\_\allowbreak SUPPORT}.

The general six-field \texttt{Hida\allowbreak UAlt} is also obtained for
all positive ranks: outside ranks two through four the point
field has an impossible antecedent.  Full \texttt{Modular\allowbreak Alt}
and \texttt{Mod\allowbreak UAlt} retain the original \texttt{Pfrom\allowbreak Aut}
premise.  The selected rank-four $P_1$ conclusion is not used
as a substitute for arbitrary-chain $P_3$.

\subsection{Every remaining selected obligation}

Each contract below is an entire selected output.  The other 30
contracts retain their original scope and dependency edges.
The next principal work is the prescribed common geometric and
arithmetic family, its three sufficient cylinder counts, the
global quantitative estimates on that family, and the remaining
one-character, congruence, rigidity and case producers.  The named
base-change, attachment, local, descent, cohomological and period
foundations remain explicit at their recorded scopes.  Current
frontier descriptions supersede historical notes in the preserved
chapters.
\paragraph{1. GLOBAL\_\allowbreak DENOMINATOR — C — Actual LinkN local table and binding producer}

Closed(Des $\mathcal C$ 3) LinkN: the actual etaSum is bounded with exact dZ0 and same-host vvTerm, subject only to the literal ShaFinOdd and summability guards.

\emph{Shared objects and choice order.} one FamP; one class $\mathcal A$ and its target-compatible subclass; one coefficient field per family datum; one x and i at each application; original selected rows, words and levels; thresholds chosen before the common family; actual D.nd; actual etaJ and Aj; sha\_\allowbreak eq

\emph{Source.} \href{https://github.com/CocoSgt/riemann-hypothesis-lean/blob/7cbf7a0c32fca1badf05931879a37c3f8a2b03b4/src/Formalism/SharpRestriction/GlobalTheoremInputs.lean#L45-L48}{LinkN}.

\emph{Source.} \href{https://github.com/CocoSgt/riemann-hypothesis-lean/blob/7cbf7a0c32fca1badf05931879a37c3f8a2b03b4/src/Formalism/SharpRestriction/GlobalTheoremInputs.lean#L40-L44}{LinkN\_\allowbreak at}.

\emph{Source.} \href{https://github.com/CocoSgt/riemann-hypothesis-lean/blob/7cbf7a0c32fca1badf05931879a37c3f8a2b03b4/src/Development/AutomorphicArithmetic/CharacterBoxes/DivinNOfHypV.lean#L117-L134}{divinN\_\allowbreak of\_\allowbreak hypV}.



\paragraph{2. GLOBAL\_\allowbreak RETAINED\_\allowbreak DEFECT — C — Original retained-row defect bound LinkIRR}

Rows.DelpNonneg and delp$\le$exp(-x\textasciicircum{}(ez/7)) uniformly for every late actual row datum.

\emph{Shared objects and choice order.} one FamP; one class $\mathcal A$ and its target-compatible subclass; one coefficient field per family datum; one x and i at each application; original selected rows, words and levels; thresholds chosen before the common family; AllRows A; LabelCountGoodInjective; AllJoinA; CountAllRows

\emph{Source.} \href{https://github.com/CocoSgt/riemann-hypothesis-lean/blob/7cbf7a0c32fca1badf05931879a37c3f8a2b03b4/src/Formalism/SharpRestriction/GlobalTheoremInputs.lean#L55-L58}{LinkIRR}.

\emph{Source.} \href{https://github.com/CocoSgt/riemann-hypothesis-lean/blob/7cbf7a0c32fca1badf05931879a37c3f8a2b03b4/src/Development/AbelianVarietyTheory/Proofs/GlobalInterpolation/LinkIRROfJoinEx.lean#L17-L38}{linkIRR\_\allowbreak of\_\allowbreak join\_\allowbreak ex}.

\emph{Source.} \href{https://github.com/CocoSgt/riemann-hypothesis-lean/blob/7cbf7a0c32fca1badf05931879a37c3f8a2b03b4/src/Development/AbelianVarietyTheory/Proofs/GlobalInterpolation/LinkIRROfJoinEx.lean#L143-L166}{link\_\allowbreak irr}.



\paragraph{3. GLOBAL\_\allowbreak PERIOD\_\allowbreak HEIGHT — C — Full actual period and height comparison LinkDIVt}

Closed conditional TuSize -> DGamSize -> (LinkEta -> DIVinN -> LinkDIV), with every original LinkDIV\_\allowbreak at clause.

\emph{Shared objects and choice order.} one FamP; one class $\mathcal A$ and its target-compatible subclass; one coefficient field per family datum; one x and i at each application; original selected rows, words and levels; thresholds chosen before the common family; same RR target; same Bu/besS; same Nu and Bk

\emph{Source.} \href{https://github.com/CocoSgt/riemann-hypothesis-lean/blob/7cbf7a0c32fca1badf05931879a37c3f8a2b03b4/src/Formalism/SharpRestriction/GlobalTheoremInputs.lean#L77-L90}{LinkDIV\_\allowbreak at}.

\emph{Source.} \href{https://github.com/CocoSgt/riemann-hypothesis-lean/blob/7cbf7a0c32fca1badf05931879a37c3f8a2b03b4/src/Formalism/SharpRestriction/GlobalTheoremInputs.lean#L94-L96}{LinkDIVt}.

\emph{Source.} \href{https://github.com/CocoSgt/riemann-hypothesis-lean/blob/7cbf7a0c32fca1badf05931879a37c3f8a2b03b4/src/Development/AbelianVarietyTheory/Proofs/GlobalInterpolation/LinkBrest.lean#L7-L25}{GlobalTheory.Top.link\_\allowbreak brest}.

\emph{Source.} \href{https://github.com/CocoSgt/riemann-hypothesis-lean/blob/7cbf7a0c32fca1badf05931879a37c3f8a2b03b4/src/Development/AbelianVarietyTheory/Proofs/GlobalInterpolation/WireSel.lean#L11-L41}{GlobalTheory.Top.wire\_\allowbreak sel}.



\paragraph{4. GLOBAL\_\allowbreak LOCAL\_\allowbreak GAIN — C — Original global LinkG installation}

One uniform LinkG with actual local lattice lengths, served-orbit coverage and MeanKappaBound; delta(x)->0 before varying actual indices.

\emph{Shared objects and choice order.} one FamP; one class $\mathcal A$ and its target-compatible subclass; one coefficient field per family datum; one x and i at each application; original selected rows, words and levels; thresholds chosen before the common family; actual G.Inst; actual GK; actual lenMN; same orbit representatives

\emph{Source.} \href{https://github.com/CocoSgt/riemann-hypothesis-lean/blob/7cbf7a0c32fca1badf05931879a37c3f8a2b03b4/src/Formalism/SharpRestriction/GlobalTheoremInputs.lean#L122-L137}{LinkG\_\allowbreak at}.

\emph{Source.} \href{https://github.com/CocoSgt/riemann-hypothesis-lean/blob/7cbf7a0c32fca1badf05931879a37c3f8a2b03b4/src/Formalism/SharpRestriction/GlobalTheoremInputs.lean#L138-L141}{LinkG}.

\emph{Source.} \href{https://github.com/CocoSgt/riemann-hypothesis-lean/blob/7cbf7a0c32fca1badf05931879a37c3f8a2b03b4/src/Development/AbelianVarietyTheory/Proofs/GlobalInterpolation/LinkG.lean#L14-L27}{GlobalTheory.Top.link\_\allowbreak g}.



\paragraph{5. GLOBAL\_\allowbreak SPECTRAL — C — Original normalized kernel/Gram estimates LinkS}

Closed TuSize -> LinkS with exact Gamma normalization, off-diagonal Hilbert-Schmidt bound, diagonal error and Dall/Dstar estimate, errors exp(-x\textasciicircum{}(epsilon/8)).

\emph{Shared objects and choice order.} one FamP; one class $\mathcal A$ and its target-compatible subclass; one coefficient field per family datum; one x and i at each application; original selected rows, words and levels; thresholds chosen before the common family; actual KD; same S.Cu spaces; GamOf; Dstar; Dall

\emph{Source.} \href{https://github.com/CocoSgt/riemann-hypothesis-lean/blob/7cbf7a0c32fca1badf05931879a37c3f8a2b03b4/src/Formalism/SharpRestriction/GlobalTheoremInputs.lean#L145-L149}{LinkS\_\allowbreak at}.

\emph{Source.} \href{https://github.com/CocoSgt/riemann-hypothesis-lean/blob/7cbf7a0c32fca1badf05931879a37c3f8a2b03b4/src/Formalism/SharpRestriction/GlobalTheoremInputs.lean#L152-L157}{LinkS}.

\emph{Source.} \href{https://github.com/CocoSgt/riemann-hypothesis-lean/blob/7cbf7a0c32fca1badf05931879a37c3f8a2b03b4/src/Development/AbelianVarietyTheory/Proofs/GlobalInterpolation/LinkAlt.lean#L16-L50}{GlobalTheory.Top.link\_\allowbreak alt}.



\paragraph{6. TASK\_\allowbreak PROP\_\allowbreak GK\_\allowbreak INSTALL — B — Install OneChar and the congruence integer into PropGK}

Complete the finite orbit/Frobenius/valuation argument producing one nonzero integer before every later coefficient prime, with the original e,N,kernel,p and fixed-K constants.

\emph{Source.} \href{https://github.com/CocoSgt/riemann-hypothesis-lean/blob/7cbf7a0c32fca1badf05931879a37c3f8a2b03b4/src/Development/AutomorphicArithmetic/AutomorphicVarieties/UnitaryThreeGroups/KatoUse.lean#L36-L42}{PropGK}.

\emph{Source.} \href{https://github.com/CocoSgt/riemann-hypothesis-lean/blob/7cbf7a0c32fca1badf05931879a37c3f8a2b03b4/src/Development/AutomorphicArithmetic/AutomorphicVarieties/UnitaryThreeGroups/OrbEquiv.lean#L111-L173}{UpperTheory.V.U3.propGK}.



\paragraph{7. TASK\_\allowbreak ONE\_\allowbreak CHAR\_\allowbreak PROOF — C — Reconstruct the original one-character argument}

From exact InpZ and AutomorphyUnramifiedCoefficients, prove the original OneChar trace/length estimate on the supplied coefficient datum and GD input.

\emph{Source.} \href{https://github.com/CocoSgt/riemann-hypothesis-lean/blob/7cbf7a0c32fca1badf05931879a37c3f8a2b03b4/src/Development/AutomorphicArithmetic/AutomorphicVarieties/UnitaryThreeGroups/KatoUse.lean#L54-L64}{OneChar}.

\emph{Source.} \href{https://github.com/CocoSgt/riemann-hypothesis-lean/blob/7cbf7a0c32fca1badf05931879a37c3f8a2b03b4/src/Development/AutomorphicArithmetic/AutomorphicVarieties/ModularityArgument/KatoQuadraticCharacters/OneChar.lean#L158-L178}{UpperTheory.V.Z.Kato.oneChar}.



\paragraph{8. TASK\_\allowbreak INP\_\allowbreak Z\_\allowbreak ASSEMBLY — B — Install the exact Kato input record}

Build the original CohInp/InpZ record using exactly lemM,lemMprime,lemP,RigII,AdVanishK and number on the same arithmetic objects. No rank-four replacement is permitted.

\emph{Source.} \href{https://github.com/CocoSgt/riemann-hypothesis-lean/blob/7cbf7a0c32fca1badf05931879a37c3f8a2b03b4/src/Development/AutomorphicArithmetic/AutomorphicVarieties/ModularityArgument/KatoQuadraticCharacters/KatoTwoOfUseN.lean#L55-L61}{InpZ}.



\paragraph{9. RIG\_\allowbreak II — C — Type-II polarized adjoint rigidity}

Marked original datum, actual type-II point and modularity of all source extensions imply vanishing for its exact adRep.

\emph{Source.} \href{https://github.com/CocoSgt/riemann-hypothesis-lean/blob/7cbf7a0c32fca1badf05931879a37c3f8a2b03b4/src/Development/AutomorphicArithmetic/AutomorphicVarieties/ModularityArgument/KatoQuadraticCharacters/KatoTwoOfUseN.lean#L49-L53}{RigII}.

\emph{Source.} \href{https://github.com/CocoSgt/riemann-hypothesis-lean/blob/7cbf7a0c32fca1badf05931879a37c3f8a2b03b4/src/Development/AutomorphicArithmetic/AutomorphicVarieties/ModularityArgument/KatoQuadraticCharacters/KatoTwoOfUseN.lean#L43-L47}{KatoTwo}.

\emph{Source.} \href{https://github.com/CocoSgt/riemann-hypothesis-lean/blob/7cbf7a0c32fca1badf05931879a37c3f8a2b03b4/src/Development/AutomorphicArithmetic/AutomorphicVarieties/ModularityArgument/KatoQuadraticCharacters/KatoTwoOfUseN.lean#L11-L41}{Used}.



\paragraph{10. AUT\_\allowbreak UNRAM\_\allowbreak COEF — C — Actual unramified-coefficient automorphy}

For each actual Dat with coefficient ramification one and fixed iota, choose one marked g1 supplying FinA, AutB and LemW together.

\emph{Source.} \href{https://github.com/CocoSgt/riemann-hypothesis-lean/blob/7cbf7a0c32fca1badf05931879a37c3f8a2b03b4/src/Development/AutomorphicArithmetic/AutomorphicVarieties/AutomorphicCases/AutomorphyUnramifiedCoefficients.lean#L24-L27}{AutomorphyUnramifiedCoefficients}.

\emph{Source.} \href{https://github.com/CocoSgt/riemann-hypothesis-lean/blob/7cbf7a0c32fca1badf05931879a37c3f8a2b03b4/src/Development/AutomorphicArithmetic/AutomorphicVarieties/AutomorphicCases/AutomorphyUnramifiedCoefficients.lean#L12-L13}{FinA}.

\emph{Source.} \href{https://github.com/CocoSgt/riemann-hypothesis-lean/blob/7cbf7a0c32fca1badf05931879a37c3f8a2b03b4/src/Development/AutomorphicArithmetic/AutomorphicVarieties/AutomorphicCases/AutomorphyUnramifiedCoefficients.lean#L15-L17}{AutB}.

\emph{Source.} \href{https://github.com/CocoSgt/riemann-hypothesis-lean/blob/7cbf7a0c32fca1badf05931879a37c3f8a2b03b4/src/Development/AutomorphicArithmetic/AutomorphicVarieties/AutomorphicCases/AutomorphyUnramifiedCoefficients.lean#L19-L22}{LemW}.



\paragraph{11. CNR3 — C — Uniform rank-three and rank-two congruence integer}

Fixed quadratic K -> c,A -> e,N,p,w,finite Galois-stable roots-of-unity orbit -> nonzero integer NN -> every later coefficient prime, closure, embedding, Frobenius and finite admissible collection.

\emph{Source.} \href{https://github.com/CocoSgt/riemann-hypothesis-lean/blob/7cbf7a0c32fca1badf05931879a37c3f8a2b03b4/src/Development/AutomorphicArithmetic/AutomorphicVarieties/AutomorphicForms/CongruenceNumberRankThree.lean#L125-L140}{CongruenceNumberRankThree}.



\paragraph{12. TASK\_\allowbreak CASE3\_\allowbreak AUT\_\allowbreak JOIN — B — Complete the original Case-III automorphy/image join}

Install the actual Case-III ProModAlt, SeedIVAlt and shared PMExAlt with the inherited CohomologyHyp/Ferrero-Washington/HFit/SeedII facts into the original whole Case-III automorphy/image conclusion.

\emph{Source.} \href{https://github.com/CocoSgt/riemann-hypothesis-lean/blob/7cbf7a0c32fca1badf05931879a37c3f8a2b03b4/src/Development/AutomorphicArithmetic/AutomorphicVarieties/ModularityArgument/ThresholdArguments/DenseThresholdCases/PromodularityCaseThree.lean#L43-L44}{AutomorphyRankFourDense}.

\emph{Source.} \href{https://github.com/CocoSgt/riemann-hypothesis-lean/blob/7cbf7a0c32fca1badf05931879a37c3f8a2b03b4/src/Development/AutomorphicArithmetic/AutomorphicVarieties/ModularityArgument/ThresholdArguments/DenseThresholdCases/CuFinal.lean#L44-L44}{main4\_\allowbreak alt}.



\paragraph{13. TASK\_\allowbreak CASE3\_\allowbreak PROMOD\_\allowbreak JOIN — C — Reconstruct the actual Case-III chain-to-promodularity proof}

Use the original tensor-factor tower, ChainRankFour, ChainDesc4, seeded rank-four chain reduction, actual PMExAlt and SeedIVAlt to construct ProModAlt on every original Ctx3 star point at the same fixed iota, preserving all finite-coefficient, inertia and support maps.

\emph{Source.} \href{https://github.com/CocoSgt/riemann-hypothesis-lean/blob/7cbf7a0c32fca1badf05931879a37c3f8a2b03b4/src/Development/AutomorphicArithmetic/AutomorphicVarieties/ModularityArgument/ThresholdArguments/DenseThresholdCases/PromodularityCaseThreeOfLemmas.lean#L131-L153}{UpperTheory.V.Cu.promodularity\_\allowbreak case\_\allowbreak three\_\allowbreak of\_\allowbreak chain\_\allowbreak reduction}.

\emph{Source.} \href{https://github.com/CocoSgt/riemann-hypothesis-lean/blob/7cbf7a0c32fca1badf05931879a37c3f8a2b03b4/src/Development/AutomorphicArithmetic/AutomorphicVarieties/ModularityArgument/ThresholdArguments/DenseThresholdCases/PromodularityCaseThree.lean#L40-L42}{PromodularityCaseThreeDense}.



\paragraph{14. SEED\_\allowbreak IV\_\allowbreak ALT — C — Actual Case-III seed with prechosen residual datum}

Base Ctx3/Four tower, iota and Amb -> w,s before every later Dies tower -> compatible finite-coefficient seed with exact residual support/root entries.

\emph{Source.} \href{https://github.com/CocoSgt/riemann-hypothesis-lean/blob/7cbf7a0c32fca1badf05931879a37c3f8a2b03b4/src/Development/AutomorphicArithmetic/AutomorphicVarieties/ModularityArgument/ThresholdArguments/DenseThresholdCases/PromodularityCaseThree.lean#L19-L25}{SeedIVAlt}.



\paragraph{15. TENSOR\_\allowbreak FACTOR\_\allowbreak TOWER — B — Actual tensor-factor tower and characters}

Original finite-field, tame, factor and character constraints; source-level construction is sizeable and has not been reconstructed here.

\emph{Source.} \href{https://github.com/CocoSgt/riemann-hypothesis-lean/blob/7cbf7a0c32fca1badf05931879a37c3f8a2b03b4/src/Development/AutomorphicArithmetic/AutomorphicVarieties/ModularityArgument/ThresholdArguments/DenseThresholdCases/PromodularityCaseThreeOfLemmas.lean#L131-L153}{UpperTheory.V.Cu.promodularity\_\allowbreak case\_\allowbreak three\_\allowbreak of\_\allowbreak chain\_\allowbreak reduction}.



\paragraph{16. CHAIN\_\allowbreak RANK\_\allowbreak FOUR — C — Actual SomeIn-to-AllIn propagation}

Same source Space and PM ideal, residual point Good4, original bound 30(C+1)\textasciicircum{}2 <= d; geometry must preserve all identities.

\emph{Source.} \href{https://github.com/CocoSgt/riemann-hypothesis-lean/blob/7cbf7a0c32fca1badf05931879a37c3f8a2b03b4/src/Development/AutomorphicArithmetic/AutomorphicVarieties/ModularityArgument/ThresholdArguments/DenseThresholdCases/PromodularityCaseThree.lean#L26-L32}{ChainRankFour}.



\paragraph{17. CHAIN\_\allowbreak DESC4 — B — Original residual chain descent}

Finite residual descent retaining source Chain4 compatibility.

\emph{Source.} \href{https://github.com/CocoSgt/riemann-hypothesis-lean/blob/7cbf7a0c32fca1badf05931879a37c3f8a2b03b4/src/Development/AutomorphicArithmetic/AutomorphicVarieties/ModularityArgument/ThresholdArguments/DenseThresholdCases/PromodularityCaseThreeOfLemmas.lean#L131-L153}{UpperTheory.V.Cu.promodularity\_\allowbreak case\_\allowbreak three\_\allowbreak of\_\allowbreak chain\_\allowbreak reduction}.



\paragraph{18. RANK4\_\allowbreak CHAIN\_\allowbreak REDUCTION — C — Seeded reduction to the original residual chain}

Select T1 before every later Deep tower; retain one compatible seed and Chain4 for every original residual lift over Q.

\emph{Source.} \href{https://github.com/CocoSgt/riemann-hypothesis-lean/blob/7cbf7a0c32fca1badf05931879a37c3f8a2b03b4/src/Development/AutomorphicArithmetic/AutomorphicVarieties/ModularityArgument/ThresholdArguments/DenseThresholdCases/PromodularityCaseThree.lean#L33-L39}{RankFourChainReductionDense}.

\emph{Source.} \href{https://github.com/CocoSgt/riemann-hypothesis-lean/blob/7cbf7a0c32fca1badf05931879a37c3f8a2b03b4/src/Development/AutomorphicArithmetic/AutomorphicVarieties/ModularityArgument/ThresholdArguments/DenseThresholdCases/CuPChainIV.lean#L16-L82}{UpperTheory.V.Cu.chainIV\_\allowbreak alt}.



\paragraph{19. TASK\_\allowbreak COUNT\_\allowbreak V\_\allowbreak INSTALL — B — Complete the original CountV finite assembly}

From exactly the listed source host/cohomology inputs, complete countV\_\allowbreak of\_\allowbreak finiteness\_\allowbreak duality\_\allowbreak euler\_\allowbreak host, including original host/frame formation, normalized length, the CaseV guard and R1<=2\textasciicircum{}c2.

\emph{Source.} \href{https://github.com/CocoSgt/riemann-hypothesis-lean/blob/7cbf7a0c32fca1badf05931879a37c3f8a2b03b4/src/Development/AutomorphicArithmetic/AutomorphicVarieties/HeightAndVanishing/Counts.lean#L69-L69}{CountV}.

\emph{Source.} \href{https://github.com/CocoSgt/riemann-hypothesis-lean/blob/7cbf7a0c32fca1badf05931879a37c3f8a2b03b4/src/Development/AutomorphicArithmetic/AutomorphicVarieties/ModularityArgument/ThresholdArguments/DenseThresholdCases/CountVOf.lean#L33-L60}{UpperTheory.V.UnitaryThreeInertia.countV\_\allowbreak of\_\allowbreak finiteness\_\allowbreak duality\_\allowbreak euler\_\allowbreak host}.



\paragraph{20. THM\_\allowbreak S6\_\allowbreak S — C — Original unitary-three ordinary support and automorphy}

Original host/frame and fixed iota -> one marked g1 -> every universal R,u -> weight-zero finiteness and every Type-I BPrime\_\allowbreak s point.

\emph{Source.} \href{https://github.com/CocoSgt/riemann-hypothesis-lean/blob/7cbf7a0c32fca1badf05931879a37c3f8a2b03b4/src/Development/AutomorphicArithmetic/AutomorphicVarieties/ModularityArgument/ThresholdArguments/DenseThresholdCases/ThmS6CoreS.lean#L17-L26}{ThmS6\_\allowbreak s}.

\emph{Source.} \href{https://github.com/CocoSgt/riemann-hypothesis-lean/blob/7cbf7a0c32fca1badf05931879a37c3f8a2b03b4/src/Development/AutomorphicArithmetic/AutomorphicVarieties/ModularityArgument/ThresholdArguments/DenseThresholdCases/ThmS6CoreS.lean#L12-L16}{BPrime\_\allowbreak s}.



\paragraph{21. HOST\_\allowbreak FRAME\_\allowbreak LENGTH — B — Original Case-V host/frame and normalized-length installation}

Actual HostData, frame, local-condition equality, abstract count and normalized original Selmer length must all refer to the same objects.

\emph{Source.} \href{https://github.com/CocoSgt/riemann-hypothesis-lean/blob/7cbf7a0c32fca1badf05931879a37c3f8a2b03b4/src/Development/AutomorphicArithmetic/AutomorphicVarieties/ModularityArgument/ThresholdArguments/DenseThresholdCases/CountVOf.lean#L33-L60}{UpperTheory.V.UnitaryThreeInertia.countV\_\allowbreak of\_\allowbreak finiteness\_\allowbreak duality\_\allowbreak euler\_\allowbreak host}.



\paragraph{22. HOST\_\allowbreak NUMBER\_\allowbreak S — C — Uniform original-host rank-three catalogue}

One exponent before all closures, coefficient rings, levels, characters, frames, embeddings and marked points; one separating trace function controls all later finite pairwise nonisomorphic collections.

\emph{Source.} \href{https://github.com/CocoSgt/riemann-hypothesis-lean/blob/7cbf7a0c32fca1badf05931879a37c3f8a2b03b4/src/Development/AutomorphicArithmetic/AutomorphicVarieties/ModularityArgument/ThresholdArguments/DenseThresholdCases/HostNumberS.lean#L18-L28}{HostNumber\_\allowbreak s}.



\paragraph{23. TASK\_\allowbreak HOST\_\allowbreak VAN\_\allowbreak E\_\allowbreak INSTALL — B — Transport K/duality to literal HostVanE}

Use the already supplied LemmaK/local-global duality and finiteness on the actual CountV host/frame and exact original coefficient modules.

\emph{Source.} \href{https://github.com/CocoSgt/riemann-hypothesis-lean/blob/7cbf7a0c32fca1badf05931879a37c3f8a2b03b4/src/Development/AutomorphicArithmetic/AutomorphicVarieties/UnitaryThreeDeformations/HostNumber.lean#L9-L12}{HostVanE}.

\emph{Source.} \href{https://github.com/CocoSgt/riemann-hypothesis-lean/blob/7cbf7a0c32fca1badf05931879a37c3f8a2b03b4/src/Development/AutomorphicArithmetic/AutomorphicVarieties/LevelCohomology/SoHostVanE.lean#L11-L18}{UpperTheory.V.X.So.hostVanE}.



\paragraph{24. TASK\_\allowbreak HOST\_\allowbreak XPP\_\allowbreak INSTALL — B — Transport CorX to literal HostXpp}

Preserve Psi\textasciicircum{}-2 eta, the prime-to-ell quotient, actual local conditions and the original CountV host while installing BR7.

\emph{Source.} \href{https://github.com/CocoSgt/riemann-hypothesis-lean/blob/7cbf7a0c32fca1badf05931879a37c3f8a2b03b4/src/Development/AutomorphicArithmetic/AutomorphicVarieties/UnitaryThreeDeformations/HostNumber.lean#L14-L17}{HostXpp}.

\emph{Source.} \href{https://github.com/CocoSgt/riemann-hypothesis-lean/blob/7cbf7a0c32fca1badf05931879a37c3f8a2b03b4/src/Development/AutomorphicArithmetic/AutomorphicVarieties/LevelCohomology/Xpp.lean#L11-L84}{UpperTheory.V.X.xpp}.



\paragraph{25. DESIGN\_\allowbreak CHOICES\_\allowbreak COMMON — C — Construct the original common design family and bundle}

For the one target HostExW, construct the original compatible DesignChoices on Cdes=WithTw(WithSm Cprime), where Cprime=WithChN(A intersect target class), A is the one Built construction class, and cT(eps,y)=cTmin(y). Retain every original field, tuple, word, metric and common bundle duty. DSZ10 and RMP11 give size/rank properties of class members, not member existence. The original RecExP/RecAtG and common object/row route remains coupled. Construct ONE original common design family R and its recursive/geometric half with original Inputs/Pins/Holds/ExAt/InGeo and required CoreP. PKT now reconstructs full exists\_\allowbreak pkT, rowPk\_\allowbreak nonempty and exists\_\allowbreak moved with Copy/Coh on those original given inputs and a supplied fitting later stage: one catalogue, both ChT, complete LfT, actual integral lattices and all metric/measure/period constants are retained; one numerical record precedes every eta and V/S/N/G. RDR supplies one Side0 before all readers, ReadPk/pins and full ReadL with exact o.eta=q.eta. The initial actual geometry, Fits/division/central conditions, primitive rational tests and eta agreement must still be produced. Inherited CSR supplies finite columns, one ReadC before later twists, period/evaluation and scalar/error transports, RowSideDetermined and same-R Bundle/BTw/BWit/BunU consumers. RPI supplies exact RowPairInt for actual pairs and finite translates with original mu\_\allowbreak b/fd0; first literal DivDef3 at the same updated R,F,Y,Z,C,sEE follows from its retained original readings, Pins/Holds, stored nonzero IsItW chain, five real-group predicates and group tie. Remaining common fields include host/frame, MultEq nonzero-part support, that original nonzero chain/group tie, eigensystem/carrier/parts and all-embedding logLam, HD/d3b, NEb, actual rational tests and base rows. Retain every simultaneous family hypothesis and eventual threshold. Stronger second DivRows3 remains separate; universal quotient-height CuspRapidDecay and actual common member existence are not inferred from the comparison constructors. Version 0.41 supplies the finite phase-correct rational theta route: exact coefficient extension, split representatives, square-law/PhiIs, Pins from CoversAt, class-sum/point-identity transfer to OneMapOverL, actual-good-projector RatG and both-level LiftTestB. Its original higher covers, all-tuple class-sum uniqueness including phase rescaling, one PointIdentCore before O, base ClassSumsSpan and geometric LatFull remain explicit producers. MLat is available from the original LatFull input. This sharpens the rational-test work without claiming actual common geometry, Fits or eta agreement.



\paragraph{26. LOCAL\_\allowbreak NBUILT\_\allowbreak COMMON — C — Actual local specialization on the common selected family}

Construct, with the required closed-class property on the same Des(C,3), one eventual family of actual NBuilt certificates for the original Fm.D(x,i), retaining the certificate before forgetting to TSel. HLP now reconstructs full HeckeLift0, all-index commutation, HeckeStd0 and HeckeIntTw from the original MtInput and complex Hecke facts, in the full ambient GenAt range k<=nV. Every maximal ideal is an actual supported residual character class; literal MtSub is the O-span of projections of the original MtLat generators, is finite free, and localizes onto the exact NCov.part. Literal Psi has its exact range and off-space projection error, both original span assertions hold, and ActEx/ActHK are constructed on the same lattice under original semisimplicity. Original MjLat relation quotients, universal coefficient-map relation fields, selected tame generators, coefficient embeddings and all identity transports give the complete nonempty\_\allowbreak spec constructor on its full printed input list. A single compositum combines the four original coefficient requirements with big fixed beforehand. Whole ndSpec\_\allowbreak exists still needs actual common geometric/choice MtInput, complex HeckeStdPlain and the no-Jordan producer, original F3 bridges with valuation/GoodCoef/iq, original F4 common-overfield reading of the entire coefficient field, and exact arithmetic character/reciprocity/admissibility suppliers. Original AutSpec, Galois specialization and the remaining NBuilt fields on that same installation also remain. Do not replace the actual projected lattice or maps by abstract objects with the same dimensions.



\paragraph{27. SECOND\_\allowbreak LAYER\_\allowbreak COMPATIBLE\_\allowbreak SELECTION — C — Original second-layer selector and class duties}

For the same A, supply Duty on A and a Closed common original TSel family whose selected records are the original true-column rec7 records (or have the entire dependent tuple transported), with their real witness and original coefficient-field root data. Existence/compatibility remains arithmetic; merely abstractly isomorphic fields or independently selected columns do not suffice. CSR21–22 now construct the same-R Bundle/BTw/BWit/BunU outputs from the exact selected realization/row-side packet and NEb, with the original EvP/ParEv quantifier order. Actual common selection, the five stated realization packets and NEb remain upstream; none is replaced by an independently chosen compatible record. Version 0.38 closes the selected analytic extension and reconstructs the first updated DivDef3 consumer; it leaves the original common selection, geometric readings, rational tests, base rows, host/parts/NEb and exact retained DI premises upstream, with unchanged EvP/ParEv thresholds. Version 0.39 supplies the original given-data reader constructors and exact RowPairInt/first DivDef3 on original premises. It leaves one actual common geometric/comparison/test member, eta agreement, compatible selection, rational tests, base rows, host/parts/NEb and original EvP/ParEv thresholds upstream. Version 0.40 also supplies original comparison/RowPk/fitting-stage Moved algorithms and exact common data transport. The actual initial geometric member, fitting tests, eta agreement and compatible second-layer selector with original eventual thresholds remain upstream. Version 0.41 supplies the finite phase-correct rational theta route: exact coefficient extension, split representatives, square-law/PhiIs, Pins from CoversAt, class-sum/point-identity transfer to OneMapOverL, actual-good-projector RatG and both-level LiftTestB. Its original higher covers, all-tuple class-sum uniqueness including phase rescaling, one PointIdentCore before O, base ClassSumsSpan and geometric LatFull remain explicit producers. MLat is available from the original LatFull input. This sharpens the rational-test work without claiming actual common geometry, Fits or eta agreement.



\paragraph{28. BASE\_\allowbreak THETA\_\allowbreak READINGS — C — Actual base reading on the common bundle}

For every family in the original class, eventually on the SAME selected Bundle and for each i<m, supply the actual R,F,Read,ReadL, weak models/Weil data and phase-one pins. Identify their LiftTie.theta, mtRead, mjRead and ThDef3 basis with that Bundle.P.A3, P.Zs, P.S0 and its stored maps. Exact field, embeddings, metrics, measures, row family and weak choices are retained. This remains one part of the coupled BTH12 producer. PKT now supplies the complete original comparison packet, RowPk and fitting-stage Moved with Copy/Coh on original given geometry, Inputs/Pins/Holds/ExAt/InGeo and Fits. Its ChT/LfT/Reps, actual lattice carriers, original Haar and norm constants, period maps, representatives and all-place RrT are constructed together; its one numerical record precedes every eta/V/S/N/G. RDR gives the original level-zero side, moved-Read and ReadPk/pin constructors on those objects, with one Side0 before every reader and all twelve Read ties. Full ReadL retains the exact o.eta(j,v)=q.eta(j,v) guard for every 1<=j<=m and place v; arbitrary Over does not imply it. Inherited CSR fixes one ReadC before all TwArgs, transports both metrics and all thirteen reading maps, and gives actual period/evaluation with kappa=(c*nP)\textasciicircum{}(-1/2). RPI supplies exact row-pair integrability, so first updated DivDef3 follows from the retained original Pins/Holds, nonzero chain, five real-group predicates and group tie. Initial theta extension remains complete. Remaining work is actual initial common geometry and its original predicates, Fits/primitive rational tests, eta agreement and outer family/eventual/same-Bundle duties, reFace compatibility, MultEq support, original chain/group data, actual base rows and any stronger DivRows3 theta square. No new compatible reading, family threshold or Haar normalization is substituted. Version 0.41 supplies the finite phase-correct rational theta route: exact coefficient extension, split representatives, square-law/PhiIs, Pins from CoversAt, class-sum/point-identity transfer to OneMapOverL, actual-good-projector RatG and both-level LiftTestB. Its original higher covers, all-tuple class-sum uniqueness including phase rescaling, one PointIdentCore before O, base ClassSumsSpan and geometric LatFull remain explicit producers. MLat is available from the original LatFull input. This sharpens the rational-test work without claiming actual common geometry, Fits or eta agreement.



\paragraph{29. BASE\_\allowbreak THETA\_\allowbreak RATIONAL\_\allowbreak TESTS — C — Simultaneous original base rational tests}

On the SAME original Bundle/reading/theta map, for every i<m supply one finite separating complex-linear test family whose values lie in the fixed embedded coefficient field on theta of EVERY original rational source vector and on EVERY stored target-coordinate basis vector. BTH8/BTH11 use this same LiftTestB for rational adjoint image and matrix coordinates. RTG now reconstructs prescribed-field coefficient extension, original split class representatives/RepsOK, the self-dual split-lattice square law and finite residue PhiIs; every Pins field follows from the actual CoversAt. With the base coefficient record literally fixed, exact class-sum existence/uniqueness including phase-rescaled O, PointIdentCore with common lambda,z before every O, and base ClassSumsSpan imply all-tuple finite lift span and OneMapOverL. The phase is exactly cj(twAt embLL A3.A.v0)(R.phs i). RatG uses the actual good projector eOp(y), and full LiftTestB follows at BOTH original levels. These finite implication proofs are complete. Remaining original producers are the prescribed higher covers, those actual class sums/uniqueness, PointIdentCore, base ClassSumsSpan and the base/common geometric datum; exact MLat is supplied only from original LatFull, whose geometric production remains separate. Finite dimension or an unrelated rational theta model cannot replace these same-object identities. Original outer eventual-family/Bundle quantifiers remain.



\paragraph{30. SLR\_\allowbreak BASE\_\allowbreak ROWS\_\allowbreak PRODUCER — C — Actual base row identities on that same bundle}

For the SAME selected Bundle, common original fields and readings used by the other three BTH12 inputs, supply actual RowsHold0\_\allowbreak 3 uniformly eventually for every family. Its original primitive row realization, spectrum/type and basis identifications must be retained. BTH supplies a theta consumer only; it does not construct the base row identity. No separate witness may be chosen merely to satisfy this row requirement.



\clearpage
\section{Arithmetic points, cyclotomic strings, and the strict cofinite bound}
\label{sec:hida-decomposition-v45}

This chapter proves the original arithmetic-point extraction theorem and
the complete \texttt{Arithmetic\allowbreak Point\allowbreak Decomposition} field, also called
\texttt{end1}.  The same point and the same supplied representation are
retained throughout.  The strict local estimate needed at the last step
follows from the generic-unitary classification already matched in BCS8
and the spherical calculation in GUA12.  It requires no additional
Ramanujan statement or stronger uniform window theorem.

We retain the original \texttt{Hypotheses\allowbreak H'}, the coefficient DVR
\(\mathcal O\), the fixed field isomorphism
\(\iota:\overline{\mathbb Q}_\ell\simeq\mathbb C\), the Hida setting
\(s\), level \(V\), auxiliary datum \(Q\), residual representation,
universal determinant ring \(B\), weight object \(\Lambda\), and
\(R_H:\mathrm{UnitaryHeckeData}(Q,\iota,\Lambda)\).
The original base-change and polarized cuspidal attachment foundations,
and their local and period foundations, have exactly the scope stated
in BCS and GUA.  We use CAP6 for the already matched degree-one
Chebotarev statement.  In particular \(\iota\) is used algebraically;
no continuity of \(\iota\) is assumed.

For a number field \(E\), write \(G_E=\operatorname{Gal}(\Omega/E)\),
\(q_w=\#(\mathcal O_E/w)\), and \(\varepsilon=\mathrm{cycBar}\).
The source convention is
\(\varepsilon(\operatorname{Frob}^{\mathrm{ar}}_w)=q_w\).
The predicate \(\mathrm{FrobEig}(\iota,r,w,S)\) includes triviality
on every inertia group above \(w\) and, for every place prolongation
and every arithmetic Frobenius lift \(g\), the identity
\[
 \iota\det(X-r(g^{-1}))=\prod_{z\in S}(X-z).
 \tag{HED.1}
\]
All products over multisets include multiplicities.

\subsection{Extracting the actual representation at a point of the ideal}

% BEGIN HED1
\paragraph{HED1. Factoring a point through the actual Hecke quotient.}
\label{claim:HED1}
Let \(f:A\twoheadrightarrow T\) be a surjective \(\mathcal O\)-algebra
map and let \(x:A\to k\) be another such map, where \(k\) is a field.
If \(\ker f\subseteq\ker x\), there is a unique
\(\mathcal O\)-algebra map \(y:T\to k\) with \(y f=x\).
\begin{proof}
For \(t=f(a)\), set \(y(t)=x(a)\).  If \(f(a)=f(a')\), then
\(a-a'\in\ker f\subseteq\ker x\), proving well-definedness.
Choose preimages of two elements to verify addition and multiplication;
the identity and the \(\mathcal O\)-algebra relation follow from those
for \(f,x\).  Surjectivity proves uniqueness.  Equivalently, this is
the quotient lift \(A/\ker f\to k\) transported through the actual
isomorphism \(A/\ker f\simeq T\).  No section of \(f\) as an algebra
map is required.
\end{proof}
% END HED1

% BEGIN HED2
\paragraph{HED2. Naturality preserves the full determinant law.}
\label{claim:HED2}
Let \(D\) be the original homogeneous multiplicative determinant
polynomial law over \(A\).  For the maps of HED1,
\[
                  (D\otimes_{A,f}T)\otimes_{T,y}k
                              =D\otimes_{A,x}k.
 \tag{HED.2}
\]
Moreover, for a group map \(p:\Gamma'\to\Gamma\), pullback of the
determinant law of a representation \(r\) is the determinant law of
\(r\circ p\).
\begin{proof}
Evaluate both sides on an arbitrary commutative \(k\)-algebra \(C\)
and an arbitrary finite sum in its group algebra.  Both coefficient
structures from \(A\) to \(C\) are identical, because \(yf=x\).
The base-changed polynomial laws therefore have the same value by
their defining naturality.  This proves equality over every scalar
algebra, including those with nilpotents.  Pullback sends a finite sum
\(\sum a_g g\) to \(\sum a_g p(g)\); its matrix under \(r\) is
\(\sum a_g r(p(g))\), exactly the matrix for \(r\circ p\).
Taking determinants proves the second assertion at the same full-law
strength.
\end{proof}
% END HED2

% BEGIN HED3
\paragraph{HED3. Transferring the entire arithmetic-point output.}
\label{claim:HED3}
Suppose the original Hecke algebra \(T=H.TK(\theta,a,K)\) has its
given \(H.\mathrm{Out}\) record for a determinant \(D_0\otimes f\)
and property \(P(\lambda,r)\).  Let \(j:A\to B\) and
\(f:B\twoheadrightarrow T\) satisfy
\(f(j(b))=\operatorname{algebraMap}_{A,T}(b)\).
If \(x:B\to\overline{\mathbb Q}_\ell\) kills \(\ker f\), and
\(xj\) is arithmetic at an admitted \(\lambda\), then one obtains
\(r:H.Gq\to\operatorname{GL}_n(\overline{\mathbb Q}_\ell)\) with
\[
 \mathrm{Continuous}(r),\quad \mathrm{IsSSRep}(r),\quad
 P(\lambda,r),\quad D_0\otimes_x\overline{\mathbb Q}_\ell=\det r.
 \tag{HED.3}
\]
\begin{proof}
Factor \(x=yf\) by HED1.  For every \(b\in A\),
\[
 y(\operatorname{algebraMap}_{A,T}(b))=y(f(j(b)))=x(j(b)).
\]
Thus the weight character induced by \(y\) is exactly the given
character induced by \(xj\).  Apply the literal \(\mathrm{Out.pt}\)
field to this \(y,\lambda\).  It supplies one representation with
continuity, semisimplicity, \(P\), and the evaluated determinant law.
HED2 identifies that law with the last law in (HED.3).
No good-point condition is needed for these four conclusions.
\end{proof}
% END HED3

% BEGIN HED4
\paragraph{HED4. The literal original \texttt{hecke\_\allowbreak ideal\_\allowbreak end}.}
\label{claim:HED4}
Retain all original data with \(n\geq1\), any allowed \(K\), and
\(R_H.\mathrm{IsJ}(B,K,J)\).  For every \(x:B.\mathrm{Pt}\) and
weight \(h\) satisfying \(J\subseteq\ker x\) and
\(x.\mathrm{Arith}(h)\), there are an admitted \(\lambda\) and one
frame representation \(r:G_0\to\operatorname{GL}_n
(\overline{\mathbb Q}_\ell)\) such that
\[
 \mathrm{Continuous}(r),\quad \mathrm{IsSSRep}(r),\quad
 \mathrm{UnitaryPointProperties}(\bar r,\iota,\lambda,r),\quad
                     x.\mathrm{IsRep}(r\circ V.\mathrm{pr}).
 \tag{HED.4}
\]
\begin{proof}
In the top-ideal branch of \(\mathrm{IsJ}\), the premise sends
\(1\in J\) to zero under \(x\), contradicting \(x(1)=1\) in the
nonzero field.  Hence that branch has no such point.

In the other branch, the actual \(\mathrm{IsJ}\) witness supplies
the original Hecke algebra \(T\), its commutativity and coefficient
conditions, a surjection \(f:B\twoheadrightarrow T\), and
\(J=\ker f\).  It also supplies the original determinant and
ordinary-character comparison and the weight-map relation
\[
 f(\Lambda.\mathrm{toPair}(B.u)(b))
                         =\operatorname{algebraMap}_{\Lambda,T}(b).
\]
The increasing-weight part of \(x.\mathrm{Arith}(h)\), together
with AFS8, gives \(\lambda\) and exactly the arithmetic Hecke
weight required by HED3.  All dominant weights of this original
\(hset\) are admitted.

HEP12 supplies the complete \(\mathrm{Out}\) record on precisely
these data; its automorphic premise is the actual one installed in
UPA17.  Apply HED3 with \(j=\Lambda.\mathrm{toPair}(B.u)\).
Its property \(P\) is the full \(\mathrm{UnitaryPointProperties}\),
so the representation carries all four fields, including all
Teichm\"uller families and the complete automorphic witness.
Its full frame determinant equals \(x.D\).  Pulling this equality
back along \(V.\mathrm{pr}\) by HED2 gives exactly the source
predicate \(x.\mathrm{IsRep}(r\circ V.\mathrm{pr})\).
The same \(r\) occurs in every field of (HED.4).
\end{proof}
% END HED4

% BEGIN HED5
\paragraph{HED5. Pullback to the original absolute Galois group.}
\label{claim:HED5}
The original map \(p=V.\mathrm{pr}:G_L\to G_0\) is surjective.
Consequently a continuous semisimple frame representation \(r\)
pulls back to a continuous semisimple representation \(r\circ p\).
\begin{proof}
Given \(g\in G_0\), lift its underlying element through the
original quotient map \(\mathrm{toGS}\).  The condition defining
\(G_0\) says that the lift has trivial \(\eta\)-character.
The original \(V.h\eta\) identifies this kernel with the image of
\(G_L\).  Choose the corresponding element of \(G_L\); its image
under \(p\) is \(g\), since both have the chosen quotient image.
This proves surjectivity on the actual source groups.

The source map \(p\) is continuous, so the composition is continuous.
A vector subspace is invariant under \(r\circ p\) if and only if
it is invariant under \(r\), by surjectivity.  An invariant
complement furnished by semisimplicity of \(r\) is therefore an
invariant complement for the pullback.  No local semisimplification
or new representation is involved.
\end{proof}
% END HED5

\subsection{Frobenius comparison with all original normalizations}

% BEGIN HED6
\paragraph{HED6. Degree-one places split in the unramified CM extension.}
\label{claim:HED6}
Let \(E/E^+\) be the original CM quadratic extension, unramified at
every finite place.  If a finite place \(w\) of \(E\) has absolute
residue degree one, then \(w\) is split over \(E^+\).
\begin{proof}
Put \(v=w|_{E^+}\).  Multiplicativity of residue degrees gives
\(f(w/v)f(v/p)=1\), hence \(f(w/v)=1\).  Unramifiedness gives
\(e(w/v)=1\).  If there were only one prime above \(v\), the
degree formula \(\sum_{w'\mid v}e(w'/v)f(w'/v)=2\) would give
\(1=2\).  Thus there are two primes above \(v\), which is the
literal splitting condition.  The argument applies to every such
\(w\); its only ramification hypothesis is the original one.
\end{proof}
% END HED6

% BEGIN HED7
\paragraph{HED7. A cofinite degree-one Frobenius comparison is global.}
\label{claim:HED7}
Let \(R:G_E\to\operatorname{GL}_n(\overline{\mathbb Q}_\ell)\)
and a finite family \(r_i:G_E\to\operatorname{GL}_{a_i}
(\overline{\mathbb Q}_\ell)\) be continuous.  Suppose that outside
one finite set, every degree-one place admits multisets \(B_{i,w}\)
such that
\[
 \mathrm{FrobEig}(\iota,R,w,\coprod_i B_{i,w}),\qquad
              \mathrm{FrobEig}(\iota,r_i,w,B_{i,w})\quad\text{for all }i.
\]
Then for every \(g\in G_E\),
\[
                  \det(X-R(g))=\prod_i\det(X-r_i(g)).
 \tag{HED.5}
\]
\begin{proof}
CAP6 puts an arithmetic Frobenius at a degree-one place outside
the finite exceptional set into every prescribed open coset of
\(G_E\).  These Frobenius elements are therefore dense.  At each
one, (HED.1) and the multiset union give (HED.5) at its inverse,
first after applying \(\iota\), then over
\(\overline{\mathbb Q}_\ell\) by injectivity of \(\iota\).

For each coefficient, the characteristic polynomial is a polynomial
in the matrix entries and is continuous.  Coefficients of the finite
product on the right are finite sums of products of those continuous
functions.  Inversion on \(G_E\) is continuous.  Since the coefficient
field is Hausdorff, equality extends from the dense set to every
\(g^{-1}\), and hence to every \(g\).  All continuity is used over
\(\overline{\mathbb Q}_\ell\), before applying the merely algebraic
\(\iota\).  Semisimplicity is not needed for this comparison lemma.
\end{proof}
% END HED7

% BEGIN HED8
\paragraph{HED8. The exact spherical Frobenius law for any real shift.}
\label{claim:HED8}
Let \(w\nmid\ell\), let \(\pi_w\) be an actual smooth irreducible
spherical representation with normalized Satake multiset \(S_w\),
and let \(t\in\mathbb R\).  If the original
\(\mathrm{ParAtT}(\iota,r,w,\mathrm{rc}(\pi_w)(t))\) holds,
with its full quantification over \(\mathrm{IsRecT}\) maps, then
\[
 \mathrm{FrobEig}\left(\iota,r,w,
              \{\alpha q_w^{-t}:\alpha\in S_w\}\right).
 \tag{HED.6}
\]
This is the original \texttt{frob\allowbreak Tw\allowbreak T} conclusion, for arbitrary real
\(t\), not only an exponent arising from a particular Speh length.
\begin{proof}
Fix any prolongation \(A\) over \(w\).  The original existence
theorem supplies a permissible local correspondence there.  GUA12
computes its spherical parameter: inertia is trivial, monodromy
is zero, and arithmetic Frobenius has eigenvalues \(\alpha^{-1}\).
The real twist multiplies these by \(q_w^t\).  Its parameter is
unramified and hence good by GUA1, so the conditional equality in
\(\mathrm{ParAtT}\) applies to the actual Weil--Deligne parameter
of \(r\).

For clarity, this equality implies triviality of the original
inertia action, as in GUA14.  Zero monodromy makes that action
trivial on an open subgroup; its image is finite.  Every inertia
matrix is thus semisimple in characteristic zero.  The original
Frobenius-semisimplification witness commutes with it and differs
by a nilpotent matrix.  Simultaneous diagonalization makes this
difference zero, giving the trivial original inertia action.
For general Weil elements the same commuting correction preserves
the characteristic polynomial.  Inverting the arithmetic
eigenvalues \(\alpha^{-1}q_w^t\) gives
\(\alpha q_w^{-t}\).  Trivial inertia makes the identity hold for
every lift, and \(A\) was arbitrary.  This proves the entire
\(\mathrm{FrobEig}\) predicate.
\end{proof}
% END HED8

% BEGIN HED9
\paragraph{HED9. Cyclotomic strings have the source's sign and multiplicities.}
\label{claim:HED9}
At \(w\nmid\ell\), the character \(\varepsilon^{-j}\),
\(j\geq0\), is unramified and has geometric Frobenius value
\(q_w^j\).  Therefore
\[
 \mathrm{FrobEig}(\iota,r,w,B)
 \ \Longrightarrow\
 \mathrm{FrobEig}(\iota,r\otimes\varepsilon^{-j},w,
                                     \{zq_w^j:z\in B\}).
 \tag{HED.7}
\]
Every twist of a continuous representation occurring here is continuous.
\begin{proof}
Inertia away from \(\ell\) fixes every \(\ell\)-power root of
unity; arithmetic Frobenius raises them to their \(q_w\)-th powers.
Passing to the defining inverse limit gives
\(\varepsilon(I_w)=1\) and
\(\varepsilon(\operatorname{Frob}^{\mathrm{ar}}_w)=q_w\).
Thus \(\varepsilon^{-j}(g^{-1})=q_w^j\) for an arithmetic
Frobenius \(g\).

For any matrix \(M\) and nonzero scalar \(u\),
\(\det(X-uM)=u^{\dim M}\det((X/u)-M)\).  If the characteristic
polynomial of \(M\) is \(\prod_{z\in B}(X-z)\), this identity
is \(\det(X-uM)=\prod_{z\in B}(X-uz)\), with every multiplicity
unchanged.  Apply it to the actual Frobenius matrix, and use the
triviality of both inertia actions.  The original cyclotomic
character is continuous by its inverse-limit construction; powers,
inversion and scalar matrix multiplication preserve continuity.
\end{proof}
% END HED9

\subsection{The complete decomposition of an arbitrary supplied representation}

% BEGIN HED10
\paragraph{HED10. One initial Galois representation for each cusp.}
\label{claim:HED10}
Let \(E/E^+\), \(\ell\ne2\), \(n\geq1\), and the actual
automorphic unitary subspace \(\sigma\) be as in the original
\texttt{end\allowbreak Shape} theorem.  Choose its BCS18 witnesses: cusps
\(P_i\), ranks \(a_i\geq1\), lengths \(d_i\geq1\), actual local
components \(\pi_{i,w}\), and infinity multisets \(m_{i,\tau}\),
with \(\sum_i a_i d_i=n\).
For each \(i\) there is a continuous representation
\(r_i:G_E\to\operatorname{GL}_{a_i}(\overline{\mathbb Q}_\ell)\)
whose original compatibility record has local parameter
\[
                  \mathrm{rc}(\pi_{i,w})\bigl((d_i-n)/2\bigr).
 \tag{HED.8}
\]
All local and period quantifiers belong to that one \(r_i\).
\begin{proof}
GUA9 supplies an infinity weight at every embedding and applies
the Speh-string identity to the BCS infinity multisets.  In
particular \(m_{i,\tau}\) has distinct entries, and
\(z+(d_i-n)/2\in\mathbb Z\) for every entry \(z\).
These are exactly the hypotheses of GUA8 with the single integer
\(e_i=d_i-n=\mathrm{ex}(n,d_i,0)\).  Apply GUA8 to \(P_i\),
its entire specified local-component family, and all its infinity
multisets.  It supplies \(r_i\) and the full compatibility record
with (HED.8), including continuity.

Make this finite family choice once.  The remaining string members
will be \(r_i\otimes\varepsilon^{-j}\); they are determined by
this initial representation.  Their construction does not require
identifying separately chosen Galois representations for different
half shifts.
\end{proof}
% END HED10

% BEGIN HED11
\paragraph{HED11. Matching the given representation at degree-one places.}
\label{claim:HED11}
Let \(R:G_E\to\operatorname{GL}_n(\overline{\mathbb Q}_\ell)\)
be any continuous representation satisfying the original
\(\mathrm{LocGlobT}\) comparison with \(\sigma\) for every split
place and every actual local component.  For the \(r_i\) of HED10,
outside one finite set and at every degree-one place \(w\), there
are multisets
\[
 B_{ij,w}=\{\alpha q_w^{(n-d_i)/2+j}:\alpha\in S_{i,w}\},
                       \qquad 0\leq j<d_i,
 \tag{HED.9}
\]
such that \(R\) has Frobenius multiset \(\coprod_{i,j}B_{ij,w}\)
and \(r_i\otimes\varepsilon^{-j}\) has multiset \(B_{ij,w}\).
\begin{proof}
GUA13 makes all the finitely many \(\pi_{i,w}\) spherical away
from one finite set.  Add the finitely many places above \(\ell\).
At any remaining degree-one place, HED6 gives splitting over
\(E^+\).  The BCS18 split-component assertion supplies the actual
irreducible normalized product of the specified Speh components.
Apply the assumed local comparison of \(R\) to that component.
GUA6 identifies its local parameter, for every permissible
correspondence, with
\[
 \bigoplus_i\bigoplus_{0\leq j<d_i}
       \mathrm{rc}(\pi_{i,w})\bigl((d_i-2j-n)/2\bigr).
\]
GUA14, whose conclusion applies to any continuous representation
with this parameter, gives the Frobenius multiset in (HED.9) for
\(R\), including trivial original inertia and all Frobenius lifts.

For each initial \(r_i\), HED8 applied to (HED.8) gives
\(\{\alpha q_w^{(n-d_i)/2}:\alpha\in S_{i,w}\}\).
HED9 multiplies this by \(q_w^j\), giving (HED.9) for the
\(j\)-th cyclotomic twist.  The common exponent identity is
\[
 \frac{n-1}{2}+j-\frac{d_i-1}{2}
                          =\frac{n-d_i}{2}+j.
\]
Thus the comparison uses the given \(R\), not a representation
chosen to replace it, and includes every positive Speh length.
\end{proof}
% END HED11

% BEGIN HED12
\paragraph{HED12. The full cyclotomic-string characteristic polynomial.}
\label{claim:HED12}
Under HED10--HED11, for every \(g\in G_E\),
\[
 \det(X-R(g))=
     \prod_i\prod_{j=0}^{d_i-1}
             \det\bigl(X-(r_i\otimes\varepsilon^{-j})(g)\bigr).
 \tag{HED.10}
\]
The ranks and lengths are positive, and every \(r_i\) is continuous.
\begin{proof}
Index the finite family in HED7 by the pairs \((i,j)\),
\(0\leq j<d_i\).  HED9 supplies continuity of every twist;
HED11 supplies exactly its cofinite degree-one Frobenius premise.
Hence HED7 proves (HED.10) at every group element.  Positivity and
the rank identity come from the same BCS witnesses, and continuity
of the initial representations was proved in HED10.  Neither a
semisimplification nor an isomorphism of the given \(R\) is needed
to establish the stated characteristic-polynomial identity.
\end{proof}
% END HED12

% BEGIN HED13
\paragraph{HED13. The strict spherical bound from the existing local classification.}
\label{claim:HED13}
Let \(F\) be a nonarchimedean characteristic-zero local field with
residue cardinality \(q>1\).  If \(\pi\) is smooth, irreducible,
generic, unitarisable and spherical, then every element \(\alpha\)
of its normalized Satake multiset satisfies
\[
                         q^{-1/2}<|\alpha|<q^{1/2}.
 \tag{HED.11}
\]
In particular this applies to every spherical local component of
the actual conjugate-self-dual cusps \(P_i\) above.
\begin{proof}
BCS8 gives the actual normalized induction
\[
          \pi\simeq\delta_1\nu^{t_1}\times\cdots\times
                             \delta_b\nu^{t_b},\qquad |t_s|<1/2,
\]
with each \(\delta_s\) unitary and square-integrable.  Since the
maximal-compact fixed space of \(\pi\) is nonzero, the tensor
formula for that fixed space in BCS14 makes every
\(\delta_s\nu^{t_s}\) spherical.  The real determinant twists
are trivial on the maximal compact, so every \(\delta_s\) is
spherical as well.

Use the original square-integrable-segment theorem, retained in
BCS9, to write each \(\delta_s\) as a segment representation.
The spherical-classification argument in GUA12 proves that a
spherical segment has length one and unramified rank-one cuspidal
support: spherical supercuspidals are unramified characters, and
a longer segment would have the same spherical Hecke character
as a determinant character, contradicting uniqueness of its
Langlands multisegment.  Thus \(\delta_s=\chi_s\) is a unitary
unramified character, and \(|\chi_s(\varpi)|=1\).

The normalized induction calculation in BCS15 identifies the
Satake multiset, with multiplicities, as
\(\{\chi_s(\varpi)q^{-t_s}\}_s\).  Its moduli are \(q^{-t_s}\).
Since \(q>1\) and \(-1/2<t_s<1/2\), both inequalities in
(HED.11) are strict.  BCS1 and BCS5 supply precisely unitarisability
and genericity for the chosen cuspidal local components.
This derives the source's needed local Satake window from its
already matched generic-unitary input; it asserts no uniform
distance of the exponents from the endpoints as the place varies.
\end{proof}
% END HED13

% BEGIN HED14
\paragraph{HED14. Singleton strings give the exact cofinite ratio at all places.}
\label{claim:HED14}
If every \(d_i=1\) in HED12, then for cofinally many finite places
\(w\), every multiset \(b\) with
\(\mathrm{FrobEig}(\iota,R,w,b)\) and every \(z,z'\in b\) satisfy
\[
                             |z|<q_w|z'|.
 \tag{HED.12}
\]
There is no final splitting condition and no extra nonzero-root
premise in this assertion.
\begin{proof}
Outside the finite set consisting of the places above \(\ell\)
and the nonspherical places of the finitely many cusps, HED8 gives
for each \(r_i\) the multiset
\[
                  C_{i,w}=\{\alpha q_w^{(n-1)/2}:
                                              \alpha\in S_{i,w}\}.
\]
These comparisons hold at every such place, whether split or not.
The identity (HED.10) is already global.  When \(d_i=1\), it says
\(\det(X-R(g))=\prod_i\det(X-r_i(g))\) for every \(g\).

Given any \(b\) in the statement, choose a prolongation and an
arithmetic Frobenius at \(w\), whose existence is part of the
original local-place formalism.  Evaluate that global identity at
its inverse and compare the \(\mathrm{FrobEig}\) equations.  It gives
\[
                  \prod_{z\in b}(X-z)
                         =\prod_i\prod_{u\in C_{i,w}}(X-u).
\]
For completeness, if \(z\in b\), evaluation at \(X=z\) makes the
right side zero.  Since \(\mathbb C\) is a field and the product
is finite, \(z\) equals an element of some \(C_{i,w}\).  Thus
every such \(z\) is \(\alpha q_w^{(n-1)/2}\) for some \(i\).
HED13 consequently gives
\[
                    q_w^{(n-2)/2}<|z|<q_w^{n/2}.
 \tag{HED.13}
\]
The same lower bound holds for \(z'\).  As \(q_w>0\),
\[
             |z|<q_w^{n/2}
                    =q_w q_w^{(n-2)/2}<q_w|z'|.
\]
The finite exceptional set was chosen before \(b,z,z'\), so all
their quantifiers have the required order.  The positive lower
bound also proves nonvanishing rather than requiring it as an
additional hypothesis.
\end{proof}
% END HED14

% BEGIN HED15
\paragraph{HED15. The complete end-shape theorem and its original wrapper.}
\label{claim:HED15}
Under the original \(\mathrm{HypothesesH'}\), let \(E/E^+\) be CM
and unramified at every finite place, \(\ell\ne2\), \(n\geq1\),
and \(\sigma\) an actual \(\mathrm{IsAutRepU}\) subspace.
For every continuous semisimple \(R:G_E\to\operatorname{GL}_n
(\overline{\mathbb Q}_\ell)\) satisfying the original
\(\mathrm{LocGlobT}\) comparison with every split local component
of \(\sigma\), there are finitely many positive ranks \(a_i\),
positive lengths \(d_i\), and continuous representations \(r_i\)
with (HED.10).  If every \(d_i=1\), the exact cofinite conclusion
(HED.12) holds for every \(\mathrm{FrobEig}\) witness of this \(R\).
\begin{proof}
BCS18 supplies one complete Speh family for the given \(\sigma\)
under the original base-change foundation.  HED10 uses the original
polarized attachment foundation through GUA8 to construct the
initial representations.  HED11 and HED12 establish the global
characteristic-polynomial identity for the supplied \(R\).
HED13 and HED14 prove the singleton-string clause.  This gives
every field of the displayed source theorem, including positivity,
continuity, every group element, and its strict cofinite conclusion.
Semisimplicity of \(R\) is retained exactly as in the target.
The \texttt{Gu\allowbreak Wrap} specialization uses these same BCS18 and GUA8
constructions, so it adds no new supplier.  All odd primes,
including \(3\), and all positive ranks and lengths remain allowed.
\end{proof}
% END HED15

% BEGIN HED16
\paragraph{HED16. The full arithmetic-point decomposition on the same ideal.}
\label{claim:HED16}
Retain the original data of HED4, the condition \(V.S'=s.S\), and
the same \(R_H.\mathrm{IsJ}(B,K,J)\).  Then the original field
\[
               \mathrm{ArithmeticPointDecomposition}(B,\iota,J)
\]
holds.
Explicitly, every \(x,h\) with \(J\subseteq\ker x\) and
\(x.\mathrm{Arith}(h)\) has a continuous semisimple
\(\rho:G_L\to\operatorname{GL}_n(\overline{\mathbb Q}_\ell)\)
with \(x.\mathrm{IsRep}(\rho)\), and positive ranks and lengths
with the full identity (HED.10) and the implication (HED.12).
\begin{proof}
HED4 supplies \(\lambda,r\) and all four unitary point properties.
Set \(\rho=r\circ V.\mathrm{pr}\).  HED5 proves its continuity
and semisimplicity, and HED4 already gives the required full
\(x.\mathrm{IsRep}\) equality.  The \(\mathrm{aut}\) field of
those same point properties supplies one actual \(\sigma\), its
\(\mathrm{IsAutRepU}\) property, and its \(\mathrm{LocGlobT}\)
comparison with this very \(\rho\) at every split local component.

The original Hida setting has odd \(\ell\), and UPA1 proves that
\(L/L^+\) is unramified at every finite place.  Hence HED15 applies
to \(\sigma,\rho\).  Use its ranks \(a_i\) as the source's
\(d_i\), and its lengths \(d_i\) as the source's \(m_i\).
The initial representations \(r_i\), their continuity and
positivity, (HED.10), and (HED.12) are exactly the remaining
fields of the original arithmetic-point decomposition.
No good-point assumption or change of ideal is introduced.
The source's original \(V.S'=s.S\) hypothesis is retained;
the extraction and comparison lemmas themselves have the broader
scopes explicitly stated above.
\end{proof}
% END HED16

\subsection{Scope of this completion}

HED4 is the original point-extraction supplier; HED16 completes the
\texttt{end1} field on the chosen original ideal.  The stronger
\texttt{Aut\allowbreak Off\allowbreak N} window, the reverse low-rank point theorem, finite
twisting and all-level modules have separate proofs and exact input
contracts.  The local strict interval in HED13 supplies no place-uniform
margin and is not substituted for that stronger window.  The present
chapter changes neither the original coefficient ring nor the choice
order needed when the Hida fields are installed on both rank-four Links.


\section{The complete ordinary automorphy-away endpoint}
\label{sec:automorphy-away-v45}

We retain the original $H'=\mathrm{HypothesesH'}$ and all its
named literature inputs.  In particular, the definite-inner-form
base-change input and polarized cuspidal attachment remain the
ones used in BCS and GUA.  The local inputs are the original
admissibility and Jacquet finite-length theorem, the geometric
lemma and uniqueness of cuspidal support, the generic unlinked
segment classification, smooth contragredient duality and
Gelfand--Kazhdan, and the full $\mathrm{IsRecT}$ laws and existence
statement.  We use the already matched generic-unitary theorem
BCS8, not a new exponent estimate.  Global strong multiplicity one,
Chebotarev, local/global reciprocity, and the basic central-character
and compact norm-one idele-class facts are retained foundations.
Below we prove their finite consumers, including the single global
unitarizing twist; in
particular the window conclusion is not retained as an opaque
extra assumption.

Fix the original $\ell,L,\Omega,s,n,V,Q,c,r,\iota,B,\Lambda,R,K$,
where $n\ge1$, $R$ is the given unitary Hecke datum, and
$\iota:\overline{\mathbb Q}_\ell\simeq\mathbb C$ is the given
field isomorphism.  The final endpoint assumes $V.S'=s.S$ and
$R.\mathrm{IsJ}(B,K,J)$, with exactly the same $J$.  No
$r.\mathrm{V0ok}$ hypothesis is added.  Write $\varepsilon$ for
the cyclotomic character.  At a finite place $w$, set
$q_w=\#(\mathcal O_L/w)>1$, and choose the source's uniformizer
$\varpi_w=\mathrm{unif}(L_w)$.  The absolute value satisfies
$|\varpi_w|_w=q_w^{-1}$.

Here is the exact local vocabulary.  A cuspidal support multiset
$S$ is \emph{toral} when all its cuspidal entries have rank one.
Its smooth rank-one characters form $\mathrm{chars}(S)$, with
the original multiplicities retained in $S$.  Put
\[
 \mathrm{Unl}(q,S):\quad
 \xi,\xi'\in\mathrm{chars}(S),\quad \xi/\xi'\text{ unramified}
       \ \Longrightarrow\ (\xi/\xi')(\varpi)\ne q .
 \tag{HAW.1}
\]
For an actual cusp $\pi$, $\mathrm{TU}(\pi,w)$ means
$\mathrm{TorAt}(\pi,w)$ and (HAW.1) for every one of its
support witnesses.  The literal strict window is
\[
 \begin{gathered}
 \mathrm{WinAt}(\pi,t,w):\quad \mathrm{TorAt}(\pi,w),\qquad
 q_w^{t-1/2}<|\xi(\varpi_w)|<q_w^{t+1/2},\\
 S\text{ a support at }w,\qquad\xi\in\mathrm{chars}(S).
 \end{gathered}
 \tag{HAW.2}
\]
For a Galois representation $\rho$, $\mathrm{Unl1}(\rho,w)$
requires at \emph{every} prolongation over $w$ an actual full flag
with characters $\kappa_i$ satisfying
$\kappa_i\ne\kappa_j\varepsilon$ for all $i,j$.
$\mathrm{PotUnl}(\rho,w)$ requires instead
\[
       \kappa_i^m\ne(\kappa_j\varepsilon)^m
                         \quad(i,j;\ m\ge1).
 \tag{HAW.3}
\]
All flag orders, principal-series normalizations, and full
determinant polynomial laws are the original ones.

\subsection{The same automorphic part and all places away from the auxiliary prime}

% BEGIN HAW1
\paragraph{HAW1. Comparison with an arbitrary given representation.}
\label{claim:HAW1}
Let $\sigma$ be an original actual $\mathrm{IsAutRepU}$ space
over a CM extension unramified at every finite place, and let
$\rho:G_L\to\mathrm{GL}_n(\overline{\mathbb Q}_\ell)$ be
continuous and irreducible.  Suppose that, outside one finite set,
at every split place every actual local component of $\sigma$
and each of its Satake multisets gives the original
$\mathrm{FrobEig}$ formula for this $\rho$.  Then there is an
actual cusp $\pi$ with
\[
 \mathrm{RegAlg}(\pi),\quad \mathrm{DualIsConj}(\pi),\quad
 \mathrm{AttC}(\iota,\rho,\pi),
 \tag{HAW.4}
\]
Satake parameters at every $\mathrm{IsUnramU}(\sigma,w)$,
and, at every split $w$, the equivalence
\[
 \pi.\mathrm{IsLocComp}(w,\tau)
 \ \Longleftrightarrow\
 \mathrm{IsCompAlong}(\sigma,\mathrm{locU},\mathrm{cLoc},\tau).
 \tag{HAW.5}
\]
Here $\tau$ ranges over every actual smooth irreducible representation.
This is the literal same-given-representation $\mathrm{link2\_\allowbreak alt}$.
\begin{proof}
GUA18 supplies a continuous semisimple representation $r_\sigma$
with the full local comparison and GUA17's irreducible-cusp
implication.  Intersect the given cofinite set with the places
away from $\ell$ and those where $\sigma$ is unramified.
An absolute residue-degree-one place in this set is split:
its relative residue degree is one, and unramifiedness in the
quadratic extension makes the degree formula force two places
above the real place.  At it choose an actual local component
by LTP19.  It is spherical by LTP20, so it has a Satake multiset.
UPA7 and the given hypothesis identify the characteristic
polynomials of $r_\sigma$ and $\rho$ at all geometric Frobenius
lifts there.

CAP6's degree-one Chebotarev density extends this equality to
every element of $G_L$: the polynomial coefficients are
continuous over $\overline{\mathbb Q}_\ell$, a Hausdorff field.
The field isomorphism $\iota$ is used only after those
equalities; no complex continuity of $\iota$ is asserted.
The two semisimple representations are isomorphic.  The finite
algebra argument also applies later without any topological
hypothesis: take the image of the group algebra in
$\operatorname{End}(W\oplus W')$.  Its radical kills the
semisimple faithful module, hence is zero.  It is a product of
matrix algebras over the algebraically closed coefficient field.
Equality of traces on group elements gives equality on this
algebra; evaluating its central block idempotents determines
the multiplicity of every simple module.  The representations
are therefore isomorphic.  Characteristic zero makes division
by each positive simple dimension legitimate.

Thus $r_\sigma$ is irreducible.  GUA17 supplies $\pi$ and
(HAW.4), transferring the attachment across this isomorphism.
It also supplies the Satake assertion at every unramified place
and the forward implication in (HAW.5).  For its reverse,
choose one actual local component $\tau'$ of $\pi$ at $w$
by BCS4.
Both $\tau'$ and the supplied $\tau$ are components of the same
$\sigma$.  Uniqueness of the split actual component gives
$\tau\simeq\tau'$, so $\tau$ is a component of $\pi$.
This keeps the supplied $\rho$ and proves all the local
quantifiers, not just a cofinite local assertion.
\end{proof}
% END HAW1

% BEGIN HAW2
\paragraph{HAW2. The literal PartLinkAlt and the partwise ordinary cusp predicate.}
\label{claim:HAW2}
Under the original $H'$, every nonzero $YOf$ in the
$\mathrm{PartLinkAlt}$ contract has an original
$Y:\mathrm{PartAltAlt}$ with $Y.Y=YOf$ such that irreducibility
of $Y.r\circ V.\mathrm{pr}$ gives (HAW.4), the entire
unramified Satake assertion, and the local comparison in
(HAW.5).  Moreover the original $\mathrm{AutUOAlt}$ holds:
the actual classical spaces are spanned by parts whose
irreducible pullbacks have such a cusp with
$\pi.\mathrm{IsLocComp}(s.w(v),Y.\sigma_l(v))$ for every
$v\mid\ell$.
\begin{proof}
Use UPA16 on the supplied type, weight, coefficient basis,
algebraic representation, and the given actual $\sigma$.
It constructs the original $Y$ with exactly $Y.Y=YOf$, a
continuous semisimple frame representation, and full split
local comparison.  Pull it back by the original continuous
$V.\mathrm{pr}$.  UPA7 converts that comparison to the
cofinite Frobenius premise of HAW1.  If the pullback is
irreducible, HAW1 applies without reselecting $Y$.
This proves all fields of $\mathrm{PartLinkAlt}$.

For $\mathrm{AutUOAlt}$ use ADP12's decomposition of the
actual vector into these $YOf$ spaces.  The given integral
type has open kernel, as in ITA12--ITA13, so that decomposition
applies with the original finite type.  Zero parts contribute
nothing.  The nonzero parts have the preceding predicate.
For each chosen $v\mid\ell$, the part's $Y.\sigma_l(v)$
is an actual component of $\sigma$.  Formula (HAW.5) says it
is a component of the same $\pi$.  Consequently the spanning
statement has the full predicate of $\mathrm{AutUOAlt}$,
including its dependence on the original part $Y$.
\end{proof}
% END HAW2

% BEGIN HAW3
\paragraph{HAW3. The plain-support dictionary and its original type supplier.}
\label{claim:HAW3}
For every positive-rank actual cusp, all its local components
are $\mathrm{PlainTor}$ at $w$ if and only if
$\mathrm{TU}(\pi,w)$.  A spherical component of a cusp has
this property.  In addition, the full original
$\mathrm{TypePlain}$ implication holds: if the parameter of
the algebraically normalized local component has zero
monodromy and the specified inertia characteristic polynomial
of restrictions of global characters, then that local
component is $\mathrm{PlainTor}$.
\begin{proof}
Write $\nu=|\cdot|_w$.  Two smooth characters have an
unramified quotient determined by its value on $\varpi_w$,
because every element of $L_w^\times$ is a uniformizer power
times a unit.  Since $\nu(\varpi_w)=q_w^{-1}$,
\[
 [\xi]=[\xi']\nu
 \quad\Longleftrightarrow\quad
 \xi'/\xi\text{ is unramified and }
                       (\xi'/\xi)(\varpi_w)=q_w.
 \tag{HAW.6}
\]
Uniqueness of cuspidal support identifies every support witness
and every isomorphic actual component.  Its rank-one entries
are characters.  Thus (HAW.6), with the ordered pair reversed
when necessary, is exactly the equivalence between the source's
no-$\nu$-linked-pair condition and (HAW.1).  This proves the
first assertion with all support witnesses, not just one.

Here is the finite converse from the supplied inertia type.
Apply the retained Langlands multisegment classification and
the $\mathrm{IsRecT}$ sum law.  For any segment beginning in
a supercuspidal $\delta$, the characteristic polynomial of
its irreducible Weil parameter divides that of the total
parameter on every Weil element.  At a commutator
$[\sigma,t]$, $\sigma$ Weil and $t$ inertia, the stipulated
global characters all equal one.  This commutator lies in
inertia, so the type hypothesis makes the total polynomial
$(X-1)^n$.  The constituent parameter is semisimple on
inertia, and a semisimple matrix with characteristic polynomial
dividing $(X-1)^n$ is the identity.  Hence every inertia
matrix of this irreducible Weil constituent commutes with the
entire Weil image.  By the elementary Schur argument it is
scalar: an eigenspace of a commuting matrix is a nonzero
invariant subspace.  Choose an eigenline of Frobenius.  Inertia
acts on it by scalars, and inertia together with Frobenius
generates the Weil group, so irreducibility forces dimension
one.  The correspondence's rank equality then forces every
cuspidal support entry to have rank one.  This is the full
abelian-converse consumer needed here; it uses the original
global characters, not arbitrary inertia characters.

The actual local component of a cusp is generic (BCS5).
The generic classification writes its algebraic normalization
as the irreducible product of mutually unlinked segments.
Zero monodromy in the parameter forces every segment to have
length one: a segment of length greater than one contributes
its nonzero standard nilpotent block, and a direct sum cannot
cancel that block.  Two distinct singleton characters whose
ratio is $\nu$ are linked segments, a contradiction.  A
character cannot equal its own $\nu$ twist, since their
uniformizer values differ by $q_w^{-1}\ne1$.
Thus the support is plain.  The common algebraic normalization
twist changes neither the ranks nor character ratios, proving
the stated $\mathrm{TypePlain}$ implication for the original
component and every permitted correspondence.

Finally a spherical generic irreducible has unramified
rank-one cuspidal support and zero-monodromy parameter, by
GUA12 and its spherical classification argument.  The preceding
generic singleton argument gives plain support.  Equivalently,
GUA12's character support combined with the unlinked generic
classification excludes the same adjacent pair.  This proves
the original $\mathrm{SphUnl}$ assertion.
\end{proof}
% END HAW3

% BEGIN HAW4
\paragraph{HAW4. ConjTor, including ramified character ratios.}
\label{claim:HAW4}
If $\mathrm{DualIsConj}(\pi)$ and $w,w'$ lie over the same
place of the maximal real subfield, then
$\mathrm{TU}(\pi,w)$ implies $\mathrm{TU}(\pi,w')$.
\begin{proof}
If $w=w'$ there is nothing to prove.  Otherwise $w'=cw$.
The original duality and conjugation laws identify an actual
component at $w'$ with the inverse-transpose transport of an
actual component at $w$ through the continuous field
isomorphism $c:L_w\simeq L_{w'}$.  Cuspidal support under
this operation is obtained by contragredients and field
transport.  To justify this exact support rule, dualize a
subquotient of normalized parabolic induction.  Smooth
contragredient is exact on admissible modules: a functional
on a submodule extends linearly and then is averaged over
an open compact subgroup fixing it.  Its average still
extends the functional.  The invariant induction pairing
identifies the dual induction with induction of the dual
factors.  Field transport preserves that induction and its
cuspidal factors.  Uniqueness of support proves the rule.

The field isomorphism carries the valuation ring, its maximal
ideal, and the unit group to their counterparts.  One way to
check the first assertion without a chosen valuation scale is
that elements of the valuation ring have bounded nonnegative
powers; this property and its inverse are preserved by a
continuous field isomorphism of local fields.  The units and
nonunits then distinguish the maximal ideal.  Thus $c$ carries
a uniformizer to a uniformizer times a unit, and the residue
cardinality remains $q_w$.

Rank-one support characters transform as
$\xi\mapsto\xi^{-1}\circ c^{-1}$.  Rank one is preserved.
For a proposed forbidden pair at $w'$, take the corresponding
pair at $w$ in the reverse order.  Their quotient is the
transport of the proposed quotient; it is unramified exactly
when that quotient is unramified.  An unramified quotient
ignores the unit multiplying the transported uniformizer, so
its value is unchanged.  A value $q_{w'}=q_w$ would contradict
(HAW.1) at $w$.  This proves all the support quantifiers of
the literal $\mathrm{ConjTor}$, with no unramifiedness
assumption on the individual characters.
\end{proof}
% END HAW4

% BEGIN HAW5
\paragraph{HAW5. The complete AutAway supplier on the original spaces.}
\label{claim:HAW5}
If $V.S'=s.S$, the original $\mathrm{hset}(Q,\tau_0)$ is
an $\mathrm{AutomorphicHeckeSet}$ for the predicate
\[
 \rho\longmapsto
 [\mathrm{IsIrr}(\rho\circ V.\mathrm{pr})
       \Longrightarrow\mathrm{AwayOk}(s,\iota,
                                      \rho\circ V.\mathrm{pr})].
 \tag{HAW.7}
\]
Here $\mathrm{AwayOk}$ includes an actual regular algebraic
dual-is-conjugate attached cusp and $\mathrm{TU}(\pi,w)$
at \emph{every} $w\nmid\ell$ whose underlying real place
is not $s.v_0$.
\begin{proof}
As in HAW2, decompose every actual admissible classical vector
by ADP12, discard zero $YOf$ spaces, and apply the literal
$\mathrm{PartLinkAlt}$ to each nonzero part.  Fix a part with
irreducible pullback and its cusp $\pi$.  If the underlying
real place of $w$ is outside $V.S'$, LTP12 makes the same
automorphic space unramified there.  HAW2 supplies Satake
parameters for $\pi$ at $w$, and HAW3 gives
$\mathrm{TU}(\pi,w)$.

Otherwise $V.S'=s.S$, together with $w\nmid\ell$ and
$\underline w\ne s.v_0$, puts $\underline w$ in $s.St$.
The chosen $w_0=s.w(\underline w)$ is split.  Choose actual
local components of $\pi$ at all chosen places over $s.St$.
To prove a statement for an arbitrary component $\tau$ at
$w_0$, replace just the $w_0$ member of this family by $\tau$.
Every member is still an actual component of $\pi$, hence of
the original unitary space by HAW2.  LTP17--LTP18 and the
given $\mathrm{TypeForAlt}$ apply to this entire updated
family and every permitted $\mathrm{IsRecT}$ map.  They give
the exact zero-monodromy inertia type whose characters are
the original global Teichmuller characters restricted along
$V.\mathrm{pr}$.  HAW3 makes this arbitrary $\tau$ plain.
Its support dictionary gives $\mathrm{TU}(\pi,w_0)$; HAW4
transports that assertion to $w$.

Thus the same part has (HAW.7).  The original vector lies in
the sum of precisely these parts, which is the full
$\mathrm{AutomorphicHeckeSet}$ assertion.  Neither the
exceptional auxiliary place nor the given type was changed.
\end{proof}
% END HAW5

% BEGIN HAW6
\paragraph{HAW6. One window center for every toral unlinked place.}
\label{claim:HAW6}
For every positive-rank actual cusp $\pi$ over a totally
complex number field, there is one $t\in\mathbb R$ such that
\[
       \mathrm{TU}(\pi,w)\ \Longrightarrow\
                  \mathrm{WinAt}(\pi,t,w)
                        \quad\text{for every finite }w.
 \tag{HAW.8}
\]
This is the full original $\mathrm{WinOf}$, with its strict
half-width and its choice of $t$ before $w$.
\begin{proof}
First let $\sigma$ be a generic unitarisable local irreducible
whose support is toral and unlinked.  By BCS8 it is the actual
normalized product of unitary square-integrable segment
representations $\delta_a\nu^{s_a}$, $|s_a|<1/2$.
Support uniqueness identifies its support with the union of
the corresponding segment supports.  Every such segment has
length one.  Otherwise its first two rank-one entries would
be $\xi$ and $\xi\nu$; their reversed quotient is unramified
with uniformizer value $q$, violating (HAW.1).  Its unitary
length-one factor is therefore a unitary rank-one character
$\zeta_a$.  Every support character has the form
$\zeta_a\nu^{s_a}$, and
\[
             |\zeta_a(\varpi)q^{-s_a}|=q^{-s_a},
                     \qquad q^{-1/2}<q^{-s_a}<q^{1/2}.
 \tag{HAW.9}
\]
Repeated characters and all positive local ranks are included.

Here is the finite global unitarization step.  Let $\zeta$
be the original continuous central Hecke character of the
positive-rank cusp $\pi$.  Its positive absolute value is
trivial on the compact norm-one idele class group: bounded
powers of a point and its inverse force modulus one.
It therefore factors through the idele norm in
$\mathbb R_{>0}$.  Taking logarithms turns a continuous
homomorphism into a continuous additive function on
$\mathbb R$, hence multiplication by one real $b$.
Thus $|\zeta|=|\cdot|_{\mathbb A}^{b}$.
Choose the single character
$\chi=|\cdot|_{\mathbb A}^{-b/n}$ and the actual twisted
cusp $\pi\otimes(\chi\circ\det)$.  Its central character
has modulus $|\cdot|_{\mathbb A}^{b}\chi^n=1$.
This is the original normalized/unitary cuspidal class;
the cuspidal positive form in BCS1 makes its local
components unitarisable.  At every finite place
$|\chi_w(\varpi_w)|=q_w^{b/n}$.  Put $a=b/n$.
This proves the precise single-twist assertion of the
original global unitarization supplier and fixes $a$
before any local place is selected.

Twisting sends each support character $\xi$ to
$\xi\chi_w$ and preserves rank and every quotient of two
characters.  Hence a toral unlinked place of $\pi$ is a
toral unlinked place of $\pi\otimes\chi$.  Its actual
component is generic by BCS5 and unitarisable by BCS1, so
(HAW.9) applies.  For each support character there is
$s\in(-1/2,1/2)$ with
$|\xi(\varpi_w)|q_w^a=q_w^{-s}$.
Consequently $|\xi(\varpi_w)|=q_w^{-a-s}$, and (HAW.2)
holds with the single choice $t=-a$.  This proves (HAW.8)
for every support witness at every eligible place.
\end{proof}
% END HAW6

\subsection{Ordinary extraction at every arithmetic point}

% BEGIN HAW7
\paragraph{HAW7. A nonzero ordinary eigenvector gives the exact principal-series subquotient.}
\label{claim:HAW7}
Let $k$ be any nonarchimedean local field and let $\sigma$ be
a smooth irreducible complex representation of
$\mathrm{GL}_n(k)$, with the original Iwahori data.  Let
$0\ne v\in\mathrm{eig}(\sigma,c,\varphi)$, where
$\varphi(e_i(u))=1$ for $u\in U_c(k,c+1)$, and suppose
\[
 U_{a_j}v=b_jv,\qquad b_j\ne0\quad(1\le j\le n),\qquad
 a_j=\operatorname{diag}(\underbrace{\varpi,\ldots,\varpi}_j,1,\ldots,1).
 \tag{HAW.10}
\]
Under the original local finite-length input there are smooth
characters $\omega_i:k^\times\to\mathbb C^\times$ such that
\[
 \begin{gathered}
 \sigma\text{ is a subquotient of }\mathrm{PS}(\omega),\qquad
 \omega_i(u)=\varphi(e_i(u))^{-1},\\
 b_j=\mathrm{psEig}(\omega,j)
       =q^{-j(j-1)/2}\prod_{i<j}\omega_i(\varpi).
 \end{gathered}
 \tag{HAW.11}
\]
The normalization is exactly
$\mathrm{PS}(\omega)=\operatorname{Ind}^{\rm norm}_B
(\omega_i|\cdot|^{(n-1)/2})_i$.
\begin{proof}
For $n=0$ the group and its irreducible representation are
one-dimensional in the only possible way, and the conclusion
is immediate.  Suppose $n\ge1$.  The positive-operator
product identity of OST5 applied to
$z=a_1\cdots a_{n-1}=\operatorname{diag}
(\varpi^{n-1},\ldots,1)$ makes $U_zv=bv$ for a nonzero
$b$ (the empty product is $1$ when $n=1$).  Thus
$v=U_z^m(b^{-m}v)$ for all $m$.  AFS11 implies that its
Jacquet image $x=p(v)$ is nonzero.

Let $J$ be the normalized Borel Jacquet module.  It has finite
length by the retained Bernstein--Zelevinsky theorem.  As
explained in AFS12, every smooth irreducible representation
of the diagonal torus is a character: an open compact
stabilizer of a cyclic vector acts trivially, and the
remaining finitely generated abelian quotient has
one-dimensional simple complex modules by the weak
Nullstellensatz.  Hence $J$ is finite-dimensional.
The finite-dimensional commutative algebra generated by the
torus action decomposes $J$ into generalized-character blocks.
Choose one block with a nonzero projection of $x$, and let
its character be $\mu$.  If $t x=cx$, projection to this
block gives the same scalar relation for a nonzero vector;
there $t-\mu(t)$ is nilpotent, so necessarily $c=\mu(t)$.
The chosen block has a nonzero character quotient: quotient
by its maximal ideal and use Nakayama, then a nonzero linear
functional.  Extend by zero on the other blocks.  This gives
a nonzero functional $l:J\to\mathbb C$ with
$l(tu)=\mu(t)l(u)$.  It need not be nonzero on $x$; the
scalar-relation property just proved is what is required.
The character is smooth because a vector on which $l$ is
nonzero has an open stabilizer.

Write $\mu(t)=\prod_i\eta_i(t_i)$.  On compact torus units
the Jacquet modulus is one, so the scalar relations of $x$
give $\eta_i(u)=\varphi(e_i(u))^{-1}$.  The explicit
Frobenius map
\[
        v'\longmapsto [g\longmapsto l(p(\sigma(g)v'))]
 \tag{HAW.12}
\]
is nonzero, smooth, equivariant, and has left $B$-law
$\delta_B^{1/2}\mu$.  Compactness of $B\backslash G$
ensures it lies in the original induced representation.
Irreducibility makes it an injection.  Put
$\omega_i=\eta_i|\cdot|^{-(n-1)/2}$; the target is exactly
$\mathrm{PS}(\omega)$ and the unit relation is unchanged.

It remains to keep the scalar normalization.  The actual
upper-unipotent coset index is $I_j=q^{j(n-j)}$.  AFS11's
unnormalized Jacquet identity and normalized torus action give
\[
 p(U_{a_j}v)=I_jp(\sigma(a_j)v),\qquad
 \mu(a_j)=\delta_B(a_j)^{-1/2}I_j^{-1}b_j.
\]
Since $\delta_B(a_j)=q^{-j(n-j)}$, we obtain
\[
 b_j=I_j\delta_B(a_j)^{1/2}\mu(a_j)
     =q^{j(n-j)/2-j(n-1)/2}\prod_{i<j}\omega_i(\varpi),
\]
which is (HAW.11), also for $j=n$.  This proves the full
original $\mathrm{OrdJac}$ without distinct unit characters
or a good-point hypothesis.
\end{proof}
% END HAW7

% BEGIN HAW8
\paragraph{HAW8. The original scalar extraction OrdScalAlt.}
\label{claim:HAW8}
Retain all data and guards in $\mathrm{OrdScalAlt}$:
$H.\mathrm{Ax}$, $\mathrm{IsCO}(O,A)$, $H.\mathrm{IsWt}(\theta)$,
an actual $\mathrm{IsSpec}(\theta,\iota,x,\lambda,\xi,c_0,sp)$,
an original part $Y$, a level $K$, and $e\ge c_0$.
If $0\ne F\in Y.Y\cap\mathrm{ClAt}(K,e)$ is an eigenvector
of every local positive generator at $v$, with nonzero
eigenvalues $b_j$, there are smooth $\omega_i$ with
\[
 \begin{split}
 Y.\sigma_l(v)&\text{ a subquotient of }\mathrm{PS}(\omega),\\
 \omega_i(u)\,\iota(\mathrm{finPart}(\theta,x,\lambda,v,i,u))&=1,\\
 b_j&=\mathrm{lowVal}(\iota,\lambda,v,j+1)^{-1}
                         \mathrm{psEig}(\omega,j+1).
 \end{split}
 \tag{HAW.13}
\]
\begin{proof}
The original $\mathrm{ClAt}$ conditions make $F$ fixed by
the local upper unipotent and the lower depth-$e$ subgroup.
The compact torus acts through the inverse of its finite-part
values.  These values give an actual character: uniqueness of
a scalar acting on the nonzero $F$ makes the identity and
multiplication laws follow from the representation law, and
invertibility makes every scalar nonzero.  The specialization's
triviality condition makes the character trivial on
$U_c(k,e+1)$, since $e\ge c_0$.

The local isotypy $Y.\mathrm{isoL}(v)$ and AFS13 supply an
equivariant linear map to $Y.\sigma_l(v)$ nonzero on $F$.
All invariance and scalar relations survive this map.  There
is no tensor-product completion or scalar extension of the
chosen part in this operation.  The global operator supported
at $v$ is the actual local operator: its coset representatives
are obtained by putting each local representative in the $v$
coordinate and $1$ in every other coordinate.  This gives
a bijection of the original cosets, without extra index.
Thus on $F$,
\[
 \mathrm{opC}_{v,j}F
       =\mathrm{lowVal}_{v,j+1}^{-1}U_{a_{j+1}}F.
\]
Every displayed low value is nonzero.  If it were zero, the
source convention for inversion would make the left side
zero, contradicting $b_jF\ne0$.  Set
$a_{j+1}=\mathrm{lowVal}_{v,j+1}b_j\ne0$ and apply HAW7 to
the nonzero local image of $F$.  Dividing its resulting
$\mathrm{psEig}$ formula by that same low value gives
(HAW.13).  These are precisely all fields of
$\mathrm{OrdScalAlt}$, for the supplied data and every $v$.
\end{proof}
% END HAW8

% BEGIN HAW9
\paragraph{HAW9. PartPtAlt for an arbitrary predicate on the actual part.}
\label{claim:HAW9}
Suppose the original $\mathrm{AutYAlt}(H,\iota,P)$ holds,
where $P(\lambda,d,\xi,Y)$ may depend on the entire part.
At every specified arithmetic specialization and every
specified $O$-algebra point $y$ of the original commuting
Hecke algebra extending $x$, there are an actual part $Y$
and $F\ne0$ such that
\[
 P(\lambda,d,\xi,Y),\qquad
 F\in Y.Y\cap\mathrm{OrdAt}(K,c_0),\qquad
 \mathrm{opC}_iF=\iota(y(\mathrm{gen}_i))F\quad\text{for all }i.
 \tag{HAW.14}
\]
The hypotheses are exactly those of $\mathrm{PartPtAlt}$,
including the specified $\mathrm{IsSpec}$ map and admitted
weight; no property of $P$ beyond $\mathrm{AutYAlt}$ is used.
\begin{proof}
CFD9--CFD12's finite averaging and polynomial generalized-eigenprojection
argument applies verbatim to the collection of parts satisfying
$P$.  Indeed $\mathrm{AutYAlt}$ supplies the needed spanning
statement for every actual classical vector; every part has
the same stability fields, and AFS16 preserves each under the
original Hecke operators.  Average over the original level,
upper and lower compact subgroups and the compact torus
character to preserve each part.  These are finite averages
over cosets of stabilizers; normality of those stabilizers is
not required, and translating cosets proves invariance.
On the resulting finite-dimensional classical space, use the
original normalized distinguished $z$-operator.  By CFD11,
the ordinary space is exactly the sum of its generalized
eigenspaces whose eigenvalues $\alpha$ satisfy
$|\iota^{-1}(\alpha)|_\ell=1$.  The Chinese remainder
theorem supplies a polynomial equal to $1$ modulo the
primary factors for precisely those eigenvalues and $0$
modulo every other primary factor, including any nonzero
nonunit eigenvalue.  This polynomial in the original
operator preserves each part and projects exactly onto
the original ordinary space.  Therefore
\[
       \mathrm{OrdAt}=\sum_{Y:P(Y)}(Y.Y\cap\mathrm{OrdAt}).
\]
ASP11 identifies this space with the complex scalar extension
of the actual finite free ordinary module; ASP13 intertwines
all the original Hecke generators.  The residual idempotent
of OHA6--OHA7 selects its specified ordinary block and
preserves every part by AFS18.  The Hecke algebra acts faithfully on
that block by its definition as a subalgebra of endomorphisms.

AFS9 now applies to the given point $y$: its module fiber is
nonzero, and the corresponding local Artinian block of the
finite commuting algebra is nonzero.  The sum above contains
a nonzero contribution in that block from some $Y$ with $P$.
In a last nonzero power of the block's nilpotent maximal
ideal, a nonzero vector is a simultaneous eigenvector for
the given character.  The specialization map transports it
injectively to the required $F$, and intertwining gives every
equation in (HAW.14).  This proves $\mathrm{PartPtAlt}$ for
the full part-dependent predicate, not only predicates of
the frame representation.
\end{proof}
% END HAW9

% BEGIN HAW10
\paragraph{HAW10. Full CoreOrdAlt with no good-point restriction.}
\label{claim:HAW10}
Retain exactly the $\mathrm{CoreOrdAlt}$ data in
\texttt{Ord\allowbreak Defs}: an original $H,A,\theta$ with its axioms,
an arbitrary $\mathrm{AutYAlt}$ predicate $P$, the specified
level and residual Hecke algebra $T$, its $\mathrm{IsCO}$
and commutativity structures, and the actual
$\mathrm{IsUniv}(D,\psi)$ record.  For every admitted
$\lambda$ and every $y:T\to_O\overline{\mathbb Q}_\ell$
arithmetic of weight $\lambda$, there exist $d,\xi,Y$ with
\[
 \begin{split}
 &P(\lambda,d,\xi,Y),\qquad D\otimes y=\det Y.r,\\
 &\forall v\ \exists\omega:\quad
 Y.\sigma_l(v)\text{ is a subquotient of }\mathrm{PS}(\omega),\\
 &\hspace{30mm}\forall i\quad
 \mathrm{OrdChar}(\iota,H.\mathrm{art}(v),\omega_i,
                       \mathrm{htW}(\lambda,i),y\psi_{v,i}).
 \end{split}
 \tag{HAW.15}
\]
The determinant equality is the full polynomial law and the
ordinary characters have the exact original order.
\begin{proof}
Put $x=y|_A$.  ASP13 gives the original arithmetic
specialization, and HAW9 gives $Y,F$ with all the required
eigenvalue equations.  HEP7's determinant argument gives
$D\otimes y=\det Y.r$; it compares the original spherical
eigenvalues and then the full characteristic-zero determinant
law, including arbitrary scalar algebras.

For each local generator, $\mathrm{IsUniv}.\psi\_\allowbreak unif$
expresses its value as a product of values of the unit-valued
characters $\psi_{v,i}$ at an Artin preimage of a uniformizer.
Thus $y(\mathrm{gen}_{v,j})\ne0$, and so does its image
under $\iota$.  Membership in $\mathrm{OrdAt}$ implies
membership in the corresponding $\mathrm{ClAt}$, so HAW8
applies to these nonzero eigenvalues.  Its inverse-unit
relation and exact $\mathrm{psEig}$ equality are precisely
the hypotheses of HEP7's local ordinary-character statement.
That statement gives the final line of (HAW.15), including
continuity of each $y\psi_{v,i}$.  All fields refer to the
same $Y,D,\psi,y$.  No pairwise distinct finite characters,
goodness condition, or flag on $Y.r$ was used.
\end{proof}
% END HAW10

% BEGIN HAW11
\paragraph{HAW11. EndOrdAlt for the same arbitrary representing rho.}
\label{claim:HAW11}
For every original $R.\mathrm{IsJ}(B,K,J)$, the full
$\mathrm{ENDO}(B,\iota,J)$ holds.  Thus if $J\subseteq\ker x$,
$x.\mathrm{Arith}(h)$, $x.\mathrm{IsRep}(\rho)$ and $\rho$
is irreducible, there is an actual regular algebraic cusp
$\pi$ with $\mathrm{DualIsConj}(\pi)$ and
$\mathrm{AttC}(\iota,\rho,\pi)$ such that, at every chosen
$w=s.w(v)$ over $\ell$, $\mathrm{ExpOf}(\iota,s.hw(v),
\pi,x.\psi_v)$ holds.  Its exact witnesses are an original
local Artin map, smooth $\omega_i$, and integers $h_i(\tau)$
strictly increasing in $i$ for every embedding, with
\[
 \mathrm{OrdChar}(\iota,a,\omega_i,h_i,x.\psi_{v,i}),\qquad
 \forall\tau_w\text{ a local component of }\pi,\quad
                 \tau_w\text{ a subquotient of }\mathrm{PS}(\omega).
 \tag{HAW.16}
\]
Continuity of the supplied $\rho$ is not a premise.
\begin{proof}
The $J=\top$ branch admits no such $x$, because $1$ would
map to zero in a field.  In the nonzero branch retain the
original surjection $f:B\twoheadrightarrow T$, its
$\mathrm{IsMap}$ record, and its compatibility with the
given weight algebra.  The map $x$ factors uniquely as
$y\circ f$, since $\ker f=J\subseteq\ker x$.
AFS6--AFS8 convert the original arithmetic weight to the
admitted dominant $\lambda$ for this same $y$.
Apply HAW10 using HAW2's $\mathrm{AutUOAlt}$ predicate.
Naturality of scalar extension and pushing characters gives
$x.D=\det Y.r$ and exactly the characters $x.\psi$ in
(HAW.15).  Pullback along $V.\mathrm{pr}$ gives
$x.\mathrm{IsRep}(Y.r\circ V.\mathrm{pr})$ as a full law.

The map $V.\mathrm{pr}$ is surjective: lift a frame element
to the original global quotient, then use $V.h\eta$ to lift
it to $G_L$.  An invariant subspace for the pullback is
therefore precisely an invariant subspace for $Y.r$.
Consequently the pullback is semisimple.  The given
irreducible $\rho$ is semisimple too.  Equality of their
determinant laws gives equality of their characteristic
polynomials, and the finite algebra argument in HAW1 gives
an isomorphism, without requiring continuity of $\rho$.
Thus the pullback is irreducible.  The part predicate supplies
$\pi$, and attachment transfers through this isomorphism.

For (HAW.16) keep the Artin map and $\omega$ in (HAW.15)
and take
$h_i(\tau)=\mathrm{htW}(\lambda,i)(\tau)
=\lambda_\tau(n-1-i)+i$.
Dominance makes the first term nondecreasing in $i$, while
$i$ is strictly increasing; hence these integer weights are
strictly increasing.  HAW2 identifies $Y.\sigma_l(v)$ with
an actual component of $\pi$.  Every other component at
that place is isomorphic to it, so the same
principal-series subquotient statement holds for each one.
This proves the full $\mathrm{ExpOf}$ and hence the original
$\mathrm{EndOrdAlt}$, for the given $\rho$ and same ideal.
\end{proof}
% END HAW11

\subsection{Ordered flags and transport between conjugate places}

% BEGIN HAW12
\paragraph{HAW12. Ordered matrix relations force the original flag.}
\label{claim:HAW12}
Let $K$ be a field, $\rho:H\to\mathrm{GL}_n(K)$, and
$\psi_i:H\to K^\times$.  Suppose
\[
 \prod_{i=0}^{n-1}(\rho(g_i)-\psi_i(g_i)I)=0,\qquad
 \det(X-\rho(g))=\prod_i(X-\psi_i(g)),
 \tag{HAW.17}
\]
for all ordered tuples and all $g$, and some $s\in H$
has pairwise distinct $\psi_i(s)$.  Then $\rho$ has the
full upper-triangular flag with ordered diagonal characters
$\psi_0,\ldots,\psi_{n-1}$.  For the original arithmetic
characters $x.\psi_v$, such a separating $s$ always exists.
\begin{proof}
Put $\lambda_i=\psi_i(s)$.  The characteristic polynomial
of $\rho(s)$ has distinct roots, so its Lagrange spectral
idempotents $E_i$ are nonzero rank-one projections whose sum
is $I$.  In (HAW.17) set every $g_j=s$ except $g_b=g$,
and multiply on the left by $E_a$ and on the right by $E_b$.
For $a\ge b$ this gives
\[
 \Bigl(\prod_{j<b}(\lambda_a-\lambda_j)\Bigr)
 E_a(\rho(g)-\psi_b(g)I)E_b
 \Bigl(\prod_{j>b}(\lambda_b-\lambda_j)\Bigr)=0.
 \tag{HAW.18}
\]
Both scalar factors are nonzero.  For $a>b$ this says
$E_a\rho(g)E_b=0$; for $a=b$ it says
$E_b\rho(g)E_b=\psi_b(g)E_b$.  Choosing one nonzero vector
in each $E_iK^n$, in this order, proves the exact flag.
No factors from different positions have been commuted.

At a place over $\ell$, arithmeticity expresses each
$\psi_i$ on inertia as a smooth finite-part character times
the algebraic character $\prod_\tau\tau(u)^{-h_i(\tau)}$.
There are finitely many local embeddings and at least one,
since the completion is a finite extension of $\mathbb Q_\ell$.
Intersect the finitely many open kernels of the finite parts
in the compact inertia group; its normal subgroup $U$ has
finite positive index $e$.  Local reciprocity lifts the
rational unit $1+\ell$ to an inertia element $\gamma$.
Then $\gamma^e\in U$, and
\[
       \psi_i(\gamma^e)
                  =((1+\ell)^e)^{-\sum_\tau h_i(\tau)}.
 \tag{HAW.19}
\]
The weights strictly increase at each embedding, so their
sums strictly increase.  Distinct integer powers of the
positive rational number $(1+\ell)^e>1$ remain distinct
in $\overline{\mathbb Q}_\ell$.  Thus $s=\gamma^e$
separates the original characters.  This uses arithmeticity,
not a good-point condition.
\end{proof}
% END HAW12

% BEGIN HAW13
\paragraph{HAW13. The full original OrdFlagJ for arbitrary irreducible rho.}
\label{claim:HAW13}
Let $B.\mathrm{RingProps}(\Lambda,J)$ hold.  If
$J\subseteq\ker x$, $x.\mathrm{Arith}(h)$,
$x.\mathrm{IsRep}(\rho)$ and $\rho$ is irreducible, then
at each chosen place over $\ell$ its restriction has
$\mathrm{HasFlag}$ with the \emph{ordered} characters
$x.\psi_{v,i}$.  This is the full $\mathrm{FlagJ}$ predicate;
there is no continuity hypothesis on $\rho$ and no
$\mathrm{V0ok}$ guard.
\begin{proof}
Irreducibility implies $n\ge1$; the zero-dimensional vector
space cannot be an irreducible nonzero module.  The original
$\phi=V.\mathrm{pr}$ is surjective: lift a frame element to
the original global quotient and apply $V.h\eta$ to obtain
its preimage in $G_L$.
We first descend $\rho$ to the frame.  If $g\in\ker\phi$,
the determinant identity $x.D\circ\phi=\det\rho$ gives
\[
          \operatorname{tr}(\rho(t)(\rho(g)-I))=0
                                           \quad(t\in G_L).
\]
The matrix span of an irreducible representation over an
algebraically closed field is all $M_n$ (the finite algebra
density theorem).  The trace pairing on $M_n$ is
nondegenerate, as is seen by testing matrix units.  Hence
$\rho(g)=I$.  Thus $\rho=\rho_0\circ\phi$, and $\rho_0$
has the same full matrix span.  Its determinant equals
$x.D$ as a full polynomial law: for any scalar algebra,
surjectivity of $\phi$ makes the induced map of group
algebras surjective, so equality after pullback implies
equality before it.

Factor $x$ through $B/J$.  Specializing $\mathrm{RingProps}.O1$
gives the characteristic polynomial identity in (HAW.17)
for the local restriction of $\rho_0$.  Specializing its
$O2$ gives the literal $\mathrm{Kills}$ assertion for the
ordered group-algebra element
$z=\prod_i([g_i]-x.\psi_i(g_i)[1])$.
To spell out its implication, use the scalar algebra
$\overline{\mathbb Q}_\ell[T]$ and test $\mathrm{Kills}$
on $T[t]$.  The coefficient of $T$ in
$\det(I+T\rho_0(t)\rho_0(z))=1$ is
$\operatorname{tr}(\rho_0(t)\rho_0(z))=0$.
The full matrix span and nondegenerate trace pairing force
$\rho_0(z)=0$.  This is the ordered product identity in
(HAW.17), with the exact source order.

HAW12 provides the separating local element and converts
these two identities into the required flag.  Pullback to
the original local decomposition group is exactly the
restriction of $\rho$ in the assertion.  Only $O1,O2$ of
the ring record were used; the conditional auxiliary
unramified field plays no role.
\end{proof}
% END HAW13

% BEGIN HAW14
\paragraph{HAW14. Polarization and full ConjUnl transport.}
\label{claim:HAW14}
Let $\rho$ be continuous, let $\pi$ satisfy
$\mathrm{DualIsConj}(\pi)$ and $\mathrm{AttC}(\iota,\rho,\pi)$,
and let $w,w'$ lie over the same real place.  Fix one
prolongation $A$ over $w$ and a full flag of $\rho|_{D_A}$
with characters $\psi_i$.  Then
\[
 \begin{split}
 \mathrm{Unl1}(\rho,w')&\Longrightarrow
                         \forall i,j\quad\psi_i\ne\psi_j\varepsilon,\\
 [\forall i,j, m\ge1\quad\psi_i^m\ne(\psi_j\varepsilon)^m]
                 &\Longrightarrow\mathrm{PotUnl}(\rho,w').
 \end{split}
 \tag{HAW.20}
\]
These are both clauses of the literal $\mathrm{ConjUnl}$.
\begin{proof}
We first record the polarization consequence used in the
transport.  For every complex-conjugation lift $c$, the
continuous semisimple representations
$\rho^c$ and $\rho^\vee\varepsilon^{1-n}$ are isomorphic.
Here semisimplicity is included in $\mathrm{AttC}$.
Indeed the original duality dictionary (finite automorphic
type and Gelfand--Kazhdan) converts $\mathrm{DualIsConj}$
to the corresponding conjugate-self-dual cuspidal class.
At cofinally many places, it sends the normalized Satake
multiset to its inverse at the conjugate place.  The attached
geometric Frobenius roots are $q_w^{(n-1)/2}\alpha$.
Dualizing and twisting by $\varepsilon^{1-n}$ sends them to
\[
       (q_w^{(n-1)/2}\alpha)^{-1}q_w^{n-1}
                             =q_w^{(n-1)/2}\alpha^{-1}.
\]
Intersect the attachment set with its conjugate to obtain
the same characteristic polynomials for the two
representations.  Chebotarev and the finite semisimple
comparison in HAW1 give their isomorphism.  Equivalently it
is a nondegenerate pairing with the original multiplier
$\varepsilon^{1-n}$.  No symmetry sign of that pairing is
needed for this consumer.  This reconstructs the finite
polarization comparison from the retained original
automorphic conjugation and Frobenius laws.

For any prolongation $A'$ over $w'$, transitivity of the
absolute Galois action on prolongations over the common real
place gives $g\in G_{L^+}$ carrying $A$ to $A'$.  Conjugation
gives an isomorphism $f:D_{A'}\simeq D_A$ and preserves
the cyclotomic character.  The restriction of $g$ to $L$ is
either the identity or complex conjugation.  In the first
case, the two restrictions of $\rho$ are conjugate after
precomposition by $f$.  In the second, write $g$ as an
element of $G_L$ times a complex-conjugation lift and use
the preceding pairing.  The second restriction is then
conjugate to
\[
          (\rho|_{D_A})^\vee\circ f\ \otimes\mu
 \tag{HAW.21}
\]
for the common character $\mu=\varepsilon^{1-n}$ pulled
through $f$.  These statements use actual decomposition
groups, not just Weil semisimplifications.

A flag transports under a group isomorphism and a conjugacy.
The dual flag has the reverse-ordered inverse characters;
a common twist multiplies all of them by $\mu$.
For each $m\ge1$, existence of a cyclotomic ratio is unchanged.
In the dual case the exact equivalence, with the pair reversed,
is
\[
 \psi_i^m=(\psi_j\varepsilon)^m
 \quad\Longleftrightarrow\quad
 (\psi_j^{-1}\mu)^m=(\psi_i^{-1}\mu\varepsilon)^m.
 \tag{HAW.22}
\]
Surjectivity of $f$ gives the same equivalence before and
after precomposition.

Finally two full flags of the same representation have the
same multiset of characters.  Their traces are the sums of
their diagonal characters.  Distinct group characters are
linearly independent: in a shortest nonzero relation,
translate by an element on which two characters differ and
subtract a scalar multiple of the original relation to
obtain a shorter one.  Thus their multiplicities agree in
characteristic zero.  This argument does not require local
semisimplicity.  If $\mathrm{Unl1}(\rho,w')$ holds, compare
its flag at any $A'$ with the transported flag and use
(HAW.22) for $m=1$; a forbidden pair among the $\psi_i$
would be a forbidden pair in that flag.  If all powers are
excluded among the $\psi_i$, (HAW.22) excludes them in the
transported flag at \emph{every} $A'$.  This is exactly
$\mathrm{PotUnl}$ and proves (HAW.20).
\end{proof}
% END HAW14

\subsection{The exact strict window excludes every torsion ratio}

% BEGIN HAW15
\paragraph{HAW15. The ordinary-character linking calculation.}
\label{claim:HAW15}
Let $w\mid\ell$, fix the original local Artin map
$a:W_{L_w}\to L_w^\times$, and suppose
$\mathrm{OrdChar}(\iota,a,\omega_i,h_i,\psi_i)$ holds,
with $h_i(\tau)$ strictly increasing in $i$ at every
embedding.  Then
\[
          \omega_i=\omega_j|\cdot|_w
                      \quad\Longrightarrow\quad
                         \psi_i=\psi_j\varepsilon.
 \tag{HAW.23}
\]
More precisely, put
$D(\tau)=h_j(\tau)-h_i(\tau)-1$ and
$\delta=\psi_i/(\psi_j\varepsilon)$.  For every Weil
element $s$, the exact identity is
\[
 \delta(s)\,\iota^{-1}(\omega_j(a(s)))q_w^{\deg s}
       =\iota^{-1}(\omega_i(a(s)))
                           \prod_\tau\tau(a(s))^{D(\tau)}.
 \tag{HAW.24}
\]
\begin{proof}
The defining ordinary-character identity is
$\psi_k(s)=\iota^{-1}(\omega_k(a(s)))
\prod_\tau\tau(a(s))^{-h_k(\tau)}$.
Local reciprocity in the original geometric Artin convention
gives
\[
 \varepsilon(s)=q_w^{\deg s}\prod_\tau\tau(a(s)),\qquad
 |a(s)|_w=q_w^{\deg s}.
 \tag{HAW.25}
\]
Divide the identities, keeping the same local embeddings;
all factors are nonzero.  This proves (HAW.24).

A continuous character of a compact group into
$\overline{\mathbb Q}_\ell^\times$ has $\ell$-adic norm
one: otherwise positive powers of it or of its inverse
would be unbounded, contrary to compactness of its image.
The $\psi_i$ are continuous by $\mathrm{OrdChar}$, hence
so is $\delta$ on the compact decomposition group.
Under the premise in (HAW.23), equations (HAW.24)--(HAW.25)
reduce on the Weil group to
$\delta(s)=\prod_\tau\tau(a(s))^{D(\tau)}$.
Choose an Artin preimage of the rational nonzero number
$\ell$.  Its value is $\ell^{\sum_\tau D(\tau)}$,
whose $\ell$-adic norm can be one only if
$\sum_\tau D(\tau)=0$.

The strictly increasing integer weights now force
$D(\tau)=0$ for every local embedding.  If $i<j$, every
$D(\tau)\ge0$ and their sum is zero.  If $i\ge j$,
every $D(\tau)<0$, which is impossible since the set of
local embeddings is nonempty.  Thus $\delta=1$ on the
Weil group.  This subgroup is dense in the decomposition
group, and $\delta$ is continuous with Hausdorff target,
so $\delta=1$ on the entire decomposition group.  This
proves (HAW.23) with no extra continuity assumption on
any ambient global representation.
\end{proof}
% END HAW15

% BEGIN HAW16
\paragraph{HAW16. Torsion ratios have exactly the forbidden factor q; full ExpTor.}
\label{claim:HAW16}
Under HAW15's hypotheses, for every $m\ge1$,
\[
 \psi_i^m=(\psi_j\varepsilon)^m
       \quad\Longrightarrow\quad
                   |\omega_j(\varpi_w)|=q_w|\omega_i(\varpi_w)|.
 \tag{HAW.26}
\]
Consequently the full original $\mathrm{ExpTor}$ holds:
if $\mathrm{ExpOf}(\iota,a,\pi,\psi)$ holds at $w\mid\ell$,
then $\mathrm{TorAt}(\pi,w)$; if also
$\psi_i\ne\psi_j\varepsilon$ for all $i,j$, then
$\mathrm{TU}(\pi,w)$; and any $\mathrm{WinAt}(\pi,t,w)$
implies the exclusions (HAW.3) for these same $\psi_i$.
\begin{proof}
Suppose $\delta^m=1$ with $\delta$ as in HAW15.
Smoothness of $\omega_i,\omega_j$ makes the intersection
of their kernels an open neighborhood of $1$ in
$L_w^\times$.  The rational units $1+\ell^M$ converge to
$1$, so choose $M\ge1$ with $N=1+\ell^M$ in that
intersection.  Here $N>1$ and its local valuation is zero.
An Artin preimage $s$ of $N$ has degree zero.
Equation (HAW.24) gives $\delta(s)=N^{\sum_\tau D(\tau)}$.
Raising to $m$ shows $N^{m\sum D}=1$ in
$\overline{\mathbb Q}_\ell$, hence in $\mathbb Q$.
Since $N>1$ and $m>0$, $\sum D=0$.  The integer-weight
argument in HAW15 again gives $D(\tau)=0$ at every local
embedding.

Take now an Artin preimage of $\varpi_w$.  Its degree is
$-1$, so (HAW.24) becomes
$\delta(s)\iota^{-1}(\omega_j(\varpi_w))q_w^{-1}
=\iota^{-1}(\omega_i(\varpi_w))$.
The complex image of $\delta(s)$ is an $m$-th root of unity
and has modulus one.  Taking complex absolute values proves
(HAW.26), with the exact factor $q_w$.

For the three $\mathrm{ExpTor}$ conclusions, take its
specified smooth characters $\omega$ and any actual local
component of $\pi$.  It is a subquotient of
$\mathrm{PS}(\omega)$, whose normalized inducing characters
are
\[
             \mu_i=\omega_i|\cdot|_w^{(n-1)/2}.
 \tag{HAW.27}
\]
Cuspidal support uniqueness identifies the support of every
component, and every support witness, with the multiset of
these rank-one characters.  This proves $\mathrm{TorAt}$.
If (HAW.1) failed for $\mu_i,\mu_j$, their common shift
cancels, and the unramified quotient of value $q_w$ gives
$\omega_j=\omega_i|\cdot|_w$.  Formula (HAW.23), with
indices interchanged, would give
$\psi_j=\psi_i\varepsilon$.  Thus the stated exclusion
implies $\mathrm{TU}$.

Finally suppose $\mathrm{WinAt}(\pi,t,w)$ and a torsion
ratio existed.  Formula (HAW.26) remains true for $\mu_i,
\mu_j$, since the shift in (HAW.27) is common.  The two
strict window inequalities would give
\[
 |\mu_j(\varpi_w)|=q_w|\mu_i(\varpi_w)|
       >q_wq_w^{t-1/2}=q_w^{t+1/2},
\]
contradicting the strict upper bound for $\mu_j$.
This proves the exclusions simultaneously for every pair
and every positive integer $m$.  There is no bound on $m$
and no loss of either endpoint of the window.
\end{proof}
% END HAW16

% BEGIN HAW17
\paragraph{HAW17. The full BornJ implication above ell.}
\label{claim:HAW17}
Assume the original $\mathrm{RingProps}(\Lambda,J)$ and
$\mathrm{ENDO}(B,\iota,J)$.  Let $x,h,\rho$ satisfy
$J\subseteq\ker x$, $x.\mathrm{Arith}(h)$,
$x.\mathrm{IsRep}(\rho)$, irreducibility, and continuity
of $\rho$.  If $w\mid\ell$ and
$\mathrm{Unl1}(\rho,w)$, then
\[
 \mathrm{PotUnl}(\rho,w),\qquad
 \forall\pi\quad\mathrm{AttC}(\iota,\rho,\pi)
                              \Longrightarrow\mathrm{TU}(\pi,w).
 \tag{HAW.28}
\]
This is the entire original $\mathrm{BornJ}$, with all
prolongations and every attached cusp in its second clause.
\begin{proof}
Use $\mathrm{ENDO}$ to obtain one actual cusp $\pi_0$,
its dual-is-conjugate property, attachment, and all its
$\mathrm{ExpOf}$ witnesses.  Put
$v=\underline w$ and $w_0=s.w(v)$.  The original place
dictionary says $w_0\mid\ell$ and
$\underline{w_0}=\underline w$.  HAW13 gives the ordered
flag $\psi=x.\psi_v$ at the chosen prolongation over $w_0$.
The first clause of HAW14 turns
$\mathrm{Unl1}(\rho,w)$ into
$\psi_i\ne\psi_j\varepsilon$.  HAW16 then makes
$\mathrm{TU}(\pi_0,w_0)$.

Choose the single global window center for $\pi_0$ by
HAW6.  Its window at $w_0$ and the last clause of HAW16
exclude every positive-power cyclotomic ratio among the
same $\psi_i$.  The second clause of HAW14 therefore gives
$\mathrm{PotUnl}(\rho,w)$ at every prolongation.

Let $\pi$ be any other attached cusp.  At cofinally many
places its Satake multiset and that of $\pi_0$ give the
same Frobenius polynomial for $\rho$.  Equality of their
root multisets, followed by cancellation of the common
nonzero scalar $q^{(n-1)/2}$, gives the same normalized
Satake multisets with multiplicities.  The original strong
multiplicity-one theorem identifies their cuspidal classes;
its actual-space version identifies the corresponding
$\mathrm{CuspSub}$, and in either formulation every local
component and support is the same up to isomorphism.
HAW4 transports $\mathrm{TU}(\pi_0,w_0)$ to
$\mathrm{TU}(\pi_0,w)$, and this identification gives
$\mathrm{TU}(\pi,w)$.  No regular-algebraic guard was
needed for this second attached cusp.
\end{proof}
% END HAW17

% BEGIN HAW18
\paragraph{HAW18. The same frame point with the full AwayOk predicate.}
\label{claim:HAW18}
If $V.S'=s.S$ and $R.\mathrm{IsJ}(B,K,J)$, every
$x,h$ with $J\subseteq\ker x$ and $x.\mathrm{Arith}(h)$
has a continuous semisimple frame representation
$\rho_0:V.\mathrm{frame}.G_0\to\mathrm{GL}_n
(\overline{\mathbb Q}_\ell)$ with
\[
 x.\mathrm{IsRep}(\rho_0\circ V.\mathrm{pr}),\qquad
 \mathrm{IsIrr}(\rho_0\circ V.\mathrm{pr})
       \Longrightarrow\mathrm{AwayOk}(s,\iota,
                                     \rho_0\circ V.\mathrm{pr}).
 \tag{HAW.29}
\]
This is the literal $\mathrm{coreA}$ output, including its
predicate and original frame.
\begin{proof}
Again $J=\top$ is impossible at the given point.  Retain
the actual surjection $B\twoheadrightarrow T$, factor
$x$ through it, and use AFS6--AFS8 to obtain the same
admitted dominant weight.  HAW5 is an
$\mathrm{AutomorphicHeckeSet}$ spanning assertion.  Viewed
as a predicate of the full part via its frame representation,
it is exactly the $\mathrm{AutYAlt}$ hypothesis in HAW10.
Apply HAW10 to that predicate and the given point of $T$.
Put $\rho_0=Y.r$ for the resulting part.  Continuity and
semisimplicity are fields of that original part.  Its
predicate gives the implication in (HAW.29), while its full
determinant law, composed with the original surjection and
$V.\mathrm{pr}$, gives the first assertion.  Thus this
construction also proves the stronger point predicate
without changing $T$, the point, the frame, or the ideal.
\end{proof}
% END HAW18

% BEGIN HAW19
\paragraph{HAW19. The full literal IrrArithPointsAutomorphicAway.}
\label{claim:HAW19}
Under exactly the original $H'$, for every original datum
with $n\ge1$, $V.S'=s.S$ and $R.\mathrm{IsJ}(B,K,J)$,
\[
              \mathrm{IrrArithPointsAutomorphicAway}(B,\iota,J)
 \tag{HAW.30}
\]
holds.  Explicitly, for every $x,h,\rho$ with
$J\subseteq\ker x$, $x.\mathrm{Arith}(h)$,
$x.\mathrm{IsRep}(\rho)$ and $\rho$ irreducible, the full
$\mathrm{AutOffN}(s.v_0,\iota,\rho)$ holds.  Its clauses are
\[
 \begin{split}
 &\mathrm{AutCsd0}(\iota,\rho),\\
 &\mathrm{Continuous}(\rho)\ \Longrightarrow\\
 &\quad\left\{\begin{array}{l}
     \forall w\mid\ell,\quad
       \mathrm{Unl1}(\rho,w)\Longrightarrow\mathrm{PotUnl}(\rho,w),\\
     \forall\pi,\ \mathrm{RegAlg}(\pi)\land
                  \mathrm{AttC}(\iota,\rho,\pi)\Longrightarrow\\
     \qquad\exists t\in\mathbb R\ \forall w,\quad
     \underline w\ne s.v_0\ \land\
       (w\nmid\ell\ \lor\ \mathrm{PotUnl}(\rho,w))
          \Longrightarrow\mathrm{WinAt}(\pi,t,w).
   \end{array}\right.
 \end{split}
 \tag{HAW.31}
\]
\begin{proof}
HAW18 provides $\rho_0$ with the same determinant law.
Surjectivity of $V.\mathrm{pr}$ makes its pullback
semisimple, as in HAW11.  The given $\rho$ is irreducible
and hence semisimple.  HAW1's finite algebra comparison
therefore gives
$\rho_0\circ V.\mathrm{pr}\simeq\rho$, without a
continuity premise on the latter.  Transport irreducibility
to the former and apply its $\mathrm{AwayOk}$ predicate.
Its regular algebraic dual-is-conjugate cusp $\pi_0$ is
attached also to $\rho$ by isomorphism.  This is exactly
$\mathrm{AutCsd0}$.

To prove the second clause, now assume
$\mathrm{Continuous}(\rho)$, exactly where the source
places this implication.  JRP15 supplies the full
$\mathrm{RingProps}(\Lambda,J)$ at this same ideal, and
HAW11 supplies $\mathrm{ENDO}$.  Apply HAW17.  It gives
the required $\mathrm{Unl1}\Rightarrow\mathrm{PotUnl}$
at every place above $\ell$.

Fix any regular algebraic attached cusp $\pi$.  Choose
$t$ once, before choosing a place, from HAW6 applied to
this $\pi$.  Let $w$ satisfy the two guards in (HAW.31).
If $w\nmid\ell$, the $\mathrm{AwayOk}$ property gives
$\mathrm{TU}(\pi_0,w)$.  The same strong-multiplicity-one
comparison as in HAW17 transports it to
$\mathrm{TU}(\pi,w)$.  If $w\mid\ell$, the other guard
is $\mathrm{PotUnl}(\rho,w)$.  Its exclusions for $m=1$
give $\mathrm{Unl1}(\rho,w)$ at every prolongation, so
HAW17 gives $\mathrm{TU}(\pi,w)$ directly.  In both cases
the already chosen $t$ and HAW6 give
$\mathrm{WinAt}(\pi,t,w)$.  This proves the exact quantifier
order and all the clauses of (HAW.31).

The argument retains the original rank, coefficient ring,
auxiliary place, level, point, representation, and ideal.
All odd primes allowed by the original Hida setting,
including $\ell=3$, remain allowed.  The branch $J=\top$
is vacuous at points and is included.  No additional
continuity, good-point, or $\mathrm{V0ok}$ hypothesis has
been added to (HAW.30).
\end{proof}
% END HAW19

\subsection{What this completion supplies}

HAW19 is the entire original Hida $\mathrm{end2}$ field.
HAW10 independently supplies the literal $\mathrm{CoreOrdAlt}$
needed by the reverse point and twist arguments, and HAW13
supplies its full ordered-flag consequence.  The proof retains
the original published foundations listed at the start and
the earlier chapters' named foundation boundaries.  In
particular it does not count a new proof of definite unitary
base change, polarized cuspidal Galois attachment, local
classification, or strong multiplicity one.  The new work
is the explicit human reconstruction of all the finite
comparisons, ordinary extraction, flag and conjugate-place
transport, and strict uniform-window consumers leading to
the literal endpoint.



\section{Returning an automorphic ordinary point to the original ideal}
\label{sec:hida-point-return-v45}

This chapter proves the original \texttt{Point\allowbreak Theorem\allowbreak Base\allowbreak Object}.
The final implication is restricted to ranks $2,3,4$, exactly as in
the source.  Its ordinary return construction is valid in every
positive rank.  All coefficient maps, the fixed algebraic isomorphism
$\iota:\overline{\mathbb Q}_\ell\simeq\mathbb C$, the integral type,
the auxiliary depth-one subgroup, and the original ideal are retained.
No containment in the kernel of the point is an input to the return.

Fix the original $\mathrm{HypothesesH}'$, Hida setting $s$, level $V$,
auxiliary data $Q$, coefficient DVR $O=c.O$, residual datum $r$,
universal ordinary determinant ring $B$, weight ring $\Lambda$,
unitary Hecke datum $R_H$, and $R_H.\mathrm{IsJ}(B,K,J)$.
In the return and final theorem assume $V.S'=s.S$.
Write $p=V.\mathrm{pr}:G_L\to G_0$ and $\varepsilon=\mathrm{cycBar}$.
The original setting assumes $\ell$ odd; $\ell=3$ is included.
Every local correspondence mentioned below means every map satisfying
the full original $\mathrm{IsRecT}$, with its actual local Artin map.
OFL1 identifies that map with the stored map $a$.

Here are the normalizations used throughout.  At a selected place $w$,
put $q=\#k_w$, choose the stored uniformizer $\varpi$, and set
\[
 \begin{split}
 \mathrm{PS}(\omega)&=\operatorname{Ind}_{B_w}^{\mathrm{GL}_n(L_w),\,\mathrm{norm}}
       \bigl(\omega_i|\,\cdot\,|^{(n-1)/2}\bigr)_{i=0}^{n-1},\\
 \mathrm{psEig}(\omega,j)&=\prod_{0\leq i<j}\omega_i(\varpi)q^{-i},
 \qquad \mathrm{lowW}(\lambda,i)_\tau=\lambda_\tau(n-1-i),\\
 h_i(\tau)&=\mathrm{htW}(\lambda,i)_\tau
                       =\mathrm{lowW}(\lambda,i)_\tau+i,\\
 \mathrm{lowVal}(\lambda,w,j)&=
       \iota\left(\prod_{i<j}\mathrm{algPow}_w
                            (\mathrm{lowW}(\lambda,i))(\varpi)\right).
 \end{split}                                                   \tag{HPT.1}
\]
Thus $h_i(\tau)$ is strictly increasing in $i$.
Our Artin convention is $\operatorname{ord}_w a(s)=-\deg s$;
on the Weil group
\[
 \varepsilon(s)=q^{\deg s}\mathrm{algPow}_w(1)(a(s)).              \tag{HPT.2}
\]
The predicate $\mathrm{OrdChar}(\omega,h,\psi)$ means smoothness of
$\omega$, continuity of $\psi$, and, for every Weil element,
\[
 \psi(s)=\iota^{-1}\omega(a(s))\,
                                \mathrm{algPow}_w(-h)(a(s)).    \tag{HPT.3}
\]
All values in (HPT.3) use the specified embedding and the actual
inclusion of the Weil group into the decomposition group.

\paragraph{Retained foundations.}
We retain exactly the original local classification, support,
local Artin reciprocity, period and local-monodromy foundations
already recorded in GUA, OFL and ITA; the existence of a full
$\mathrm{IsRecT}$ map is retained, as there.  We also retain
$\mathrm{HypothesesH}'.\mathrm{descU}:\mathrm{DescWeak}$ and the
original Jacquet--Shalika isobaric rigidity input
$\mathrm{litN.aut.js}.1$.  Neither is declared newly proved here.
The former applies only to $2\leq n\leq4$, a CM field unramified
over its maximal real subfield at every finite place, and a finite
exception set consisting of split places and containing all places
above $\ell$.  It returns a unitary-group automorphic representation
unramified at nonsplit places and a \emph{cofinite split-place}
Satake comparison.  HPT9 proves the stronger comparison needed here
from that precise output.  The latter input identifies two finite
positive-rank cuspidal families when their unions of normalized
Satake parameters agree at cofinitely all finite places.

\subsection{Point uniqueness and local reverse comparison}

% BEGIN HPT1
\paragraph{HPT1. The evaluated pair determines the entire point.}
\label{claim:HPT1}
Two $O$-algebra points $x,y:B\to\overline{\mathbb Q}_\ell$ are equal
if their evaluated determinant polynomial laws and all their ordered
characters $\psi_{v,i}$ agree.
\begin{proof}
ASP2 puts both entire images in the integer ring $O_E$ of one finite
extension $E/\mathbb Q_\ell$.  It also proves that maximal-ideal
elements have norm less than one.  Both residue maps therefore have
image in the prescribed copy of $k_O$.  Set
\[
                A_0=\{a\in O_E:\bar a\in k_O\}=O+\mathfrak m_E.
\]
This is a finite $O$-module, being an $O$-submodule of the finite
module $O_E$.  Its units are exactly the elements outside
$\mathfrak m_E$: the inverse of a nonzero residue in $k_O$ again
has residue in $k_O$.  It is consequently local with maximal ideal
$\mathfrak m_E$ and residue field $k_O$.  It is Noetherian and
complete: finite modules over the complete DVR $O$ are complete,
and the $\mathfrak m_{A_0}$-adic and $\mathfrak m_O$-adic topologies
agree.  For the last assertion, $A_0/\mathfrak m_OA_0$ is an
Artinian local ring, so a power of its maximal ideal is zero.
Thus $A_0$ satisfies the literal original $\mathrm{IsCO}(O,A_0)$.

Regard $x,y$ as local maps to $A_0$.  Mapping $B.u$ along either
map gives a pair over $A_0$: continuity follows because a local
map sends each maximal-ideal power into the corresponding power;
the residual, polarization, ordinary and inertia relations follow
by applying the map.  The two resulting pairs agree.  Indeed the
characters agree in the subring $A_0\subset\overline{\mathbb Q}_\ell$.
For the determinant law, evaluate on a generic finite group-algebra
sum $\sum T_i g_i$.  Its value lies in $A_0[T_1,\ldots,T_m]$,
which injects into $\overline{\mathbb Q}_\ell[T_1,\ldots,T_m]$.
The assumed full-law equality therefore identifies all coefficients
already in $A_0$.  Evaluation on any scalar algebra proves equality
of the laws over $A_0$.  The uniqueness clause of the original
universal property of $B$ now gives $x=y$.  This argument uses no
topological-generation assertion about selected traces.
\end{proof}
% END HPT1

% BEGIN HPT2
\paragraph{HPT2. Trace relations give the original inertia type.}
\label{claim:HPT2}
Suppose $x.\mathrm{IsRep}(\rho)$ and $\rho$ is semisimple.  If a
subgroup is killed by $p$, it is killed by $\rho$.  At each original
type place $u\in s.S_t$, and for \emph{every} original Teichm\"uller
family $t$, the full inertia action of $\rho$ is conjugate to the
ordered diagonal family $t_i\circ p$ after the specified coefficient
map.  Its characters have open kernels.  Every corresponding
away-$\ell$ Weil--Deligne parameter has precisely this inertia type.
\begin{proof}
If $p(g)=1$, the full-law equality gives
$\operatorname{tr}(\rho(h)(\rho(g)-1))=0$ for every $h\in G_L$.
The algebra generated by a semisimple representation over an
algebraically closed field is a product of full matrix algebras on
its distinct simple constituents.  The trace pairing of its given
representation is the sum of the ordinary matrix trace pairings,
with positive integral multiplicities.  It is nondegenerate in
characteristic zero.  Thus $\rho(g)=1$.

The same proof applied to the literal $B.u.I$ relation shows that
$\rho$ kills the common kernel of the residual characters on
inertia at $u$.  The quotient embeds into a product of copies of
$k_O^\times$, hence is finite abelian of exponent dividing
$N=|k_O|-1$.  Its characteristic-zero representation is a sum of
characters of order dividing $N$.  Residual characteristic
polynomials, given by $B.u.\mathrm{red}$, agree with those of the
ordered residual family.  Reduction is injective on $N$-th roots
of unity, since $\ell\nmid N$.  The eigenvalue multisets therefore
agree with the Teichm\"uller lifts at every inertia element.
Linear independence of distinct characters of the finite abelian
quotient identifies their multisets.  Ordering the summands gives
one conjugating matrix valid on \emph{all} inertia elements.

For completeness, each lifted character has open kernel.  The
residual kernel is open.  On it a Teichm\"uller value $z$ satisfies
$z^N=1$ and $z\equiv1$; the factor $1+z+\cdots+z^{N-1}$ is a unit,
so $(z-1)(1+\cdots+z^{N-1})=0$ forces $z=1$.
For an away-$\ell$ $\mathrm{IsWDOf}$ witness, finite inertia makes
monodromy zero.  Its raw Weil action agrees with $\rho$ on inertia.
The Frobenius semisimplification correction commutes with inertia
and is unipotent; finite-order inertia matrices are semisimple,
so this correction does not change them.  This is the literal
inertia conclusion used in GUA14.  Applying $\iota$ gives the
original $\mathrm{IsTypeWD}$ for every allowed correspondence.
\end{proof}
% END HPT2

% BEGIN HPT3
\paragraph{HPT3. Locally algebraic characters give the exact ordinary characters.}
\label{claim:HPT3}
A continuous decomposition-group character $\psi$ is locally
algebraic of weight $h$ precisely when there is a smooth $\omega$
satisfying (HPT.3).  For such $\omega$, the finite inertia part
$\phi$ in $\mathrm{IsWtPart}$ satisfies
\[
                 \iota\phi(\gamma)=\omega(a(\gamma))^{-1}.
                                                               \tag{HPT.4}
\]
Two continuous characters satisfying (HPT.3) with the same
$\omega,h$ are identical on the entire decomposition group.
\begin{proof}
Given the locally algebraic decomposition, on the Weil group put
\[
       \xi(s)=\iota(\psi(s)/\mathrm{algPow}_w(-h)(a(s))).
\]
On inertia its kernel is exactly the open kernel of the finite
part.  Local Artin reciprocity therefore factors this smooth
character through $a$, giving $\omega$ and (HPT.3).
Conversely define $\phi$ by (HPT.4).  Its kernel is open by
smoothness of $\omega$ and continuity of the Artin map on inertia.
Multiplying (HPT.3) by $\phi$ gives exactly the algebraic-weight
identity in $\mathrm{IsWtPart}$.  This identity holds for every
permitted local Artin map, since OFL1 identifies all of them with
the stored map.  These arguments prove both directions with the
same finite part.  Finally the Weil group is dense in the
decomposition group.  The equality in (HPT.3) and continuity of
the two characters extend their equality from that dense subgroup.
No topology is transported through $\iota$.
\end{proof}
% END HPT3

% BEGIN HPT4
\paragraph{HPT4. Strict weights and the forbidden cyclotomic ratio exclude linkage.}
\label{claim:HPT4}
Suppose $h_i(\tau)$ is strictly increasing in $i$, each
$\mathrm{OrdChar}(\omega_i,h_i,\psi_i)$ holds, and
$\psi_i\ne\psi_j\varepsilon$ for every $i,j$.
Then for $i\ne j$ the characters $\omega_i,\omega_j$ are distinct
and neither is the other times $|\,\cdot\,|^{\pm1}$.
\begin{proof}
A continuous character of a compact group into
$\overline{\mathbb Q}_\ell^\times$ has norm one.  Powers of an
element and its inverse lie in a compact set, so both norms must
be at most one.  If on the Weil group such a character is
$\mathrm{algPow}_w(d)\circ a$, choose a Weil element mapping to
the rational element $\ell\in L_w^\times$.  Its norm is
$|\ell|_\ell^{\sum_\tau d(\tau)}$, so
$\sum_\tau d(\tau)=0$.  The local embedding set is finite and
nonempty because $w\mid\ell$.

If $\omega_i=\omega_j$, (HPT.3) makes $\psi_i/\psi_j$ equal to
the algebraic character with exponents $h_j-h_i$.  Strict increase
makes all of these exponents have the same nonzero sign, contradicting
their zero sum.  If $\omega_i=\omega_j|\,\cdot\,|$, equations
(HPT.2)--(HPT.3) make $\psi_i/(\psi_j\varepsilon)$ have exponents
$h_j-h_i-1$.  For $i>j$ they are all negative.  For $i<j$ they are
all nonnegative, so their zero sum forces every one to vanish.
The ratio is then one on the Weil group and, by density, on the
decomposition group.  This is the prohibited ratio.  Interchanging
$i,j$ treats the inverse norm twist.
\end{proof}
% END HPT4

% BEGIN HPT5
\paragraph{HPT5. A given flag determines jumps and Weil polynomials at every bridge.}
\label{claim:HPT5}
Let $\rho_w$ have a full ordered flag with characters $\psi_i$
satisfying $\mathrm{OrdChar}(\omega_i,h_i,\psi_i)$.
Suppose the original $\mathrm{DRat}$ datum for $\rho_w$ has jumps
$H_\tau$ and Weil characteristic polynomial $P_s(X)$.
Then
\[
 H_\tau=\{h_i(\tau):0\leq i<n\},\qquad
 P_s(X)=\prod_{i=0}^{n-1}(X-\omega_i(a(s))).                    \tag{HPT.5}
\]
These identities hold for every original period ring, bridge,
induced embedding and Weil element in the definition of
$\mathrm{DRat}$.
\begin{proof}
Fix any of those original choices.  The supplied $\mathrm{DRat}$
already makes $\rho_w$ de Rham; closure under subobjects and
quotients makes all terms of its given flag de Rham.  We do not
infer this from the characters alone.  The rank-one calculation OFL4
applies to each character, using HPT3 to identify the actual finite
part: its de Rham jump is $h_i(\tau)$ and its raw Weil character is
$\omega_i\circ a$.  The original period foundations give exact
de Rham and potentially semistable functors on these representations,
with strict induced filtrations.  One may choose a common finite
coefficient field and common finite extension making the finitely
many successive terms semistable; their existence and coefficient
compatibility are the same retained laws used in OFL12.
Apply exactness successively to the supplied flag.  Strictness says
that the multiplicity of each jump is the sum of its multiplicities
in the one-dimensional quotients.  In a basis adapted to the
resulting Weil flag the raw action is upper triangular, with these
rank-one diagonal characters, so its characteristic polynomial is
the product in (HPT.5).  Frobenius semisimplification preserves that
polynomial.  Descent of the chosen finite extension preserves the
flag maps and the polynomial identities.  Finally $\mathrm{DRat}$
uses these same functors and the same representation; uniqueness
of their filtered dimensions and characteristic polynomials gives
the displayed identities.  Since the choices at the start were
arbitrary, the entire original quantifier is preserved.
\end{proof}
% END HPT5

% BEGIN HPT6
\paragraph{HPT6. Recovering the principal series from a toral parameter.}
\label{claim:HPT6}
Let $\tau$ be an actual irreducible smooth local representation with
the original $\mathrm{PlainTor}$ property.  Suppose for every
$\mathrm{IsRecT}$ map the parameter of $\mathrm{algNorm}(\tau)$
has characteristic polynomial $\prod_i(X-\omega_i(a(s)))$ at
every Weil element.  Then $\tau$ is a subquotient of
$\mathrm{PS}(\omega)$.  If the characters satisfy the conclusion
of HPT4, $\tau\simeq\mathrm{PS}(\omega)$.
\begin{proof}
Choose a cuspidal support of $\mathrm{algNorm}(\tau)$.
The original $\mathrm{PlainTor}$ statement makes all its members
characters.  The multisegment exhaustion and support laws of the
original local classification, together with the product law of
$\mathrm{IsRecT}$, identify the characteristic polynomial of its
parameter with the product of the rank-one support parameters.
The rank-one correspondence law makes these parameters
$\xi\circ a$.  OFL1 ensures that the Artin map here is precisely
the one in the hypothesis.

Taking traces gives equality of sums of the support characters
with $\sum_i\omega_i$ on $L_w^\times$, since $a$ is onto.
Linear independence of distinct characters identifies their
multisets, including repetitions.  Undoing
$\mathrm{algNorm}=|\det|^{(1-n)/2}\otimes(-)$ shows that the
support of $\tau$ is the multiset
$\{\omega_i|\,\cdot\,|^{(n-1)/2}\}_i$.  The support theorem
therefore places $\tau$ among the irreducible subquotients of
the normalized induction of these characters, which is exactly
our $\mathrm{PS}(\omega)$.

Under HPT4 these singleton segments are pairwise unlinked.
The original totally disjoint induction theorem, starting with
rank one and applying induction in stages, makes their product
irreducible.  An irreducible subquotient of this irreducible
representation is the representation itself.  This proves the
claimed actual intertwining isomorphism, not merely a support
or semisimplified-parameter equality.
\end{proof}
% END HPT6

% BEGIN HPT7
\paragraph{HPT7. The original automorphic window gives plain torality and reverse unramifiedness.}
\label{claim:HPT7}
If $\mathrm{WinAt}(\pi,t,w)$ holds, every actual local component
$\tau$ of $\pi$ is $\mathrm{PlainTor}$ and every
$\mathrm{rc}(\mathrm{algNorm}(\tau))$ is toral with zero
monodromy.  At $w\nmid\ell$, if this is the Weil--Deligne
parameter of a representation trivial on full inertia, then
$\tau$ is spherical.
\begin{proof}
The literal $\mathrm{WinAt}$ first makes every support member a
rank-one character and then puts the norm of its uniformizer value
strictly between $q^{t-1/2}$ and $q^{t+1/2}$.  Two such characters
cannot differ by $|\,\cdot\,|^{\pm1}$: their norm ratio would be
$q$ or $q^{-1}$, contradicting the two strict inequalities.
Equal characters cause no linkage of singleton segments.
The local component of a cusp is generic by BCS5.  Apply the
original $\mathrm{GenericUnlinkedGL}$ classification: it is an
irreducible product of mutually unlinked segments.  Each segment
has rank-one cuspidal support, and a segment of length at least
two would contribute an adjacent pair differing by
$|\,\cdot\,|$, which the strict window forbids.  All segments
therefore have length one, including any repeated identical
singletons.  Thus this actual irreducible product is the full
principal series of the support characters.  This uses the
generic unlinked-segment theorem, rather than the stronger
pairwise-distinct hypothesis in the totally disjoint theorem.
It proves $\mathrm{PlainTor}$
with its actual inducing data.  The rank-one and product laws of
every $\mathrm{IsRecT}$ give a direct sum of rank-one Weil
characters, with zero monodromy.  Commutators act identically,
which is the original torality condition.  GUA2 supplies $\mathrm{Gd}$.

Now suppose the full Galois inertia is trivial.  Its raw
Weil--Deligne inertia is trivial and stays so under Frobenius
semisimplification.  In the just computed direct sum each support
character is consequently trivial on Artin images of inertia.
These images are precisely $O_{L_w}^\times$.  The common
normalizing norm twist is also trivial there, so all inducing
characters are unramified.  In their normalized principal series,
the Iwasawa formula gives a nonzero $\mathrm{GL}_n(O_{L_w})$-fixed
function with value one at the identity.  The irreducible induction
is $\tau$, proving its sphericity.  This proves the reverse
implication rather than appealing to injectivity of a spherical
parameter assignment.
\end{proof}
% END HPT7

% BEGIN HPT8
\paragraph{HPT8. Compatibility for the supplied representation and its supplied flag.}
\label{claim:HPT8}
Suppose $\rho$ is continuous and $\mathrm{IsAutCsdA}(\iota,\rho)$.
Choose the regular algebraic conjugate-self-dual cusp $\pi$ in its
actual $\mathrm{AutCsd0}$ witness.  For every $w\nmid\ell$, and
for every $w\mid\ell$ where the supplied flag has no ratio
$\psi_i=\psi_j\varepsilon$, the actual local components are
$\mathrm{PlainTor}$ and the full $\mathrm{LocGlobT}$ comparison
holds for this same $\rho$.  At a selected $\ell$-place with an
arithmetic ordered flag of weights $h_i$, the component is
$\mathrm{PS}(\omega)$ with (HPT.3), and its Hodge multiset is
$\{h_i\}_i$.
\begin{proof}
The $\mathrm{AutCsd0}$ witness includes semisimplicity and the
cofinite attachment of this $\rho$ to $\pi$.  Polarize $\pi$ by
the original character $\delta_{L/L^+}^{n}$: its pullback along
the norm is one and its sign is $(-1)^n$.  Thus every hypothesis
of the original RowT application GUA7 holds.  It produces a
continuous semisimple representation $\rho_0$, its Hodge
multisets and its full guarded local compatibility.
At cofinitely many places the cusp is spherical and away from
$\ell$.  GUA12 and HED8 with $t=(1-n)/2$ convert the RowT
comparison into its cofinite attachment to $\pi$, with the
original algebraic Satake normalization.  Its Frobenius
polynomials thus equal those of $\rho$.  Chebotarev
and continuity identify the characteristic polynomials everywhere;
Brauer--Nesbitt gives one constant intertwiner
$\rho\simeq\rho_0$.  Enlarge a finite coefficient field to contain
this matrix and its inverse.  Tensoring that matrix with every
period ring transports the filtered modules and Weil actions,
so the full Hodge and local comparison statements hold for $\rho$.
This is the same finite-coefficient descent used in ADK1 and GUA10;
it preserves all bridges and every $\mathrm{IsRecT}$ quantifier.

Continuity of each character in the supplied flag follows by
conjugating the matrix representation into upper triangular form
and taking a diagonal entry.  A flag at one prolongation supplies
the same kind of flag at every prolongation: conjugate the
decomposition subgroup by a global Galois element and conjugate
the flag basis by its matrix.  This group isomorphism is onto and
the cyclotomic character is conjugation invariant.  The forbidden
ratio therefore stays forbidden.  Hence the precise $\mathrm{Unl1}$
premise of $\mathrm{WitA}$ holds.  Its first field upgrades this
to $\mathrm{PotUnl}$; its second field provides a single window
for this attached cusp at all away-$\ell$ places and these
above-$\ell$ places.  HPT7 gives plain torality and $\mathrm{Gd}$,
so RowT's toral, zero-monodromy branch supplies the stated
$\mathrm{LocGlobT}$.  No rank-dependent alternative branch is used.

Arithmeticity gives local algebraicity of the ordered characters.
HPT3 supplies $\omega_i$.  Apply HPT5 to the actual Hodge and
local comparison just obtained, then HPT6 and HPT4.
This gives the full principal series and the ordered-weight
Hodge multiset.  All steps use the original $\rho$, with no
assumption that the point already kills $J$.
\end{proof}
% END HPT8

% BEGIN HPT9
\paragraph{HPT9. Upgrading weak descent to every split place and infinity.}
\label{claim:HPT9}
Let $\pi$ be the cusp of HPT8 and let $\sigma$ be a unitary-group
automorphic representation with the cofinite split Satake
comparison supplied by $\mathrm{DescWeak}$.  Then every actual
split local component of $\sigma$ is the corresponding component
of $\pi$, up to an actual intertwiner, and the infinity multiset
of $\sigma$ equals that of $\pi$.
\begin{proof}
Use BCS18 and GUA18 on this actual $\sigma$.  Write its cuspidal
pieces as $\Pi_b$ of ranks $a_b>0$, with lengths $d_b>0$, and
write $R_\sigma$ for the continuous semisimple representation
produced there.  At the cofinite split places where all objects
are spherical and away from $\ell$, GUA14 and the weak comparison
identify the characteristic polynomials of $R_\sigma$ and
$\rho$.  Every degree-one place is split by HED6.  HED7 therefore
extends this equality to every element of $G_L$.

At cofinitely \emph{all} finite places, BCS19 and GUA14 compute
the same polynomial as the union of the strings
\[
       \{\alpha q_w^{k-(d_b-1)/2}:
                  \alpha\in\mathrm{Sat}(\Pi_b,w),\ 0\leq k<d_b\}.
                                                               \tag{HPT.6}
\]
The attachment of $\rho$ to $\pi$ identifies (HPT.6) with
$\mathrm{Sat}(\pi,w)$.  Interpret the left side as the family
of actual cuspidal twists
$\Pi_b\otimes|\det|^{(d_b-1)/2-k}$ indexed by $(b,k)$.
The Satake sign is correct because twisting by $|\det|^u$
multiplies a Satake parameter by $q_w^{-u}$.
The retained Jacquet--Shalika rigidity input applies to these
two finite positive-rank cusp families over the totally complex
field $L$.  It gives a bijection of this index set with the
singleton indexing $\pi$.  Since all $d_b$ are positive, there
is exactly one $b$, its $d_b$ is one, its rank is $n$, and
its cusp is $\pi$ itself.

For these witnesses the full split product in BCS18 is just the
length-one local component of $\pi$.  The length-one Speh and
one-factor induction maps are identities.  Local component
uniqueness LTP3 identifies every other actual split component
of $\sigma$ with it.  The BCS18 infinity equality becomes
exactly the infinity multiset of $\pi$.  This argument proves
the strong comparison for this weak descent; it does not assert
that arbitrary unitary-group automorphic representations have
only length-one Speh factors.
\end{proof}
% END HPT9

\subsection{Constructing the nonzero vector in the original coefficient space}

% BEGIN HPT10
\paragraph{HPT10. Simultaneous global compact fixed vectors.}
\label{claim:HPT10}
Let $\sigma$ be the actual irreducible automorphic unitary-group
representation.  Assume it is unramified outside $V.S'$ and its
actual component at $v_0$ has a nonzero $K_c(1)$-fixed vector.
There is a nonzero vector fixed by the original subgroup $K_0$
which is maximal compact outside $V.S'$, is $K_c(1)$ at $v_0$,
and is the identity at the selected $\ell$-places, type places
and auxiliary places.
\begin{proof}
For a finite place of $L^+$, restriction to its coordinate factors
and replacement of the other coordinates by one gives a retraction
of the finite adelic unitary group onto its local factor $H_w$.
The unitary equation is preserved because the entire pair of
conjugate coordinates is retained.  Thus the group is a product
of $H_w$ with its commuting complementary factor.  The retraction
preserves $K_0$.  Its local factors in $K_0$ are compact: the
integral unitary factor is a closed subgroup of the corresponding
matrix compact group, and $K_c(1)$ is compact as well.

We first record a fixed-vector transfer valid also at nonsplit
places.  If $W$ is a nonzero $H_w$-stable subspace of $\sigma$,
global irreducibility says that the span of its global translates
is $\sigma$.  By the product decomposition it is the sum of
translates $cW$ by the complementary factor.  Normalized averaging
over a compact local subgroup $C_w$ is an exact projector on a
smooth complex representation: on each vector its orbit map is
locally constant, so its integral is a finite sum.  This projector
commutes with $c$.  If $W^{C_w}=0$, the projector vanishes on every
$cW$ and hence on $\sigma$.  Therefore $\sigma^{C_w}\ne0$
implies $W^{C_w}\ne0$.  At the places under consideration the
required nonzero fixed vector in $\sigma$ comes from the stated
unramified condition or from the actual $v_0$ component; LTP5
transfers the latter vector to $\sigma$.

Choose any nonzero smooth vector in $\sigma$.  Its open stabilizer
and the restricted-product topology give a finite set $T$ of
places such that every element of $K_0$ which is one at $T$ fixes
it.  Induct on $|T|$.  For $w\in T$, the nonzero subspace fixed
by these elements of $K_0$ is $H_w$-stable, because the supports
are disjoint.  The preceding averaging argument produces in it
a nonzero vector fixed also by $K_0\cap H_w$.  Every element of
$K_0$ which is one at $T\setminus\{w\}$ is a product of its
local retraction at $w$ and an element which is one at all of $T$.
The new vector is fixed by both.  Removing one place at a time
gives the asserted nonzero $K_0$-fixed vector.  This proof requires
no choice of a nonsplit local Langlands parameter.
\end{proof}
% END HPT10

% BEGIN HPT11
\paragraph{HPT11. The integral type and the infinity dual give the actual nonzero part.}
\label{claim:HPT11}
Assume the hypotheses of HPT10, the original infinity weights
$\lambda_\tau^\vee$, and, at every type place, an actual local
component whose parameter has the prescribed type for every
$\mathrm{IsRecT}$ and every original Teichm\"uller family.
For the original $\mathrm{TypeForAlt}$ representation $\tau_0$,
every specified algebraic basis $e$ with $\mathrm{IsBasis}(\iota,e,\xi)$
gives
\[
                         Y\mathrm{Of}(Q,\iota,\tau_0,e,\xi,\sigma)\ne0.
                                                               \tag{HPT.7}
\]
\begin{proof}
Choose the original Teichm\"uller family using ITA2.  Put
$W=\sigma^{K_0}$, which is nonzero by HPT10.  The actual product
section $J_t$ of the type groups supplied by LTP16 commutes with
$K_0$.  Local component isotypy LTP3 and the given local parameter
types permit the full product type-detection statement ITA12 on
$W$.  It gives a nonzero column $F_0=(F_{0,i})\in W^m$ with
\[
       \sum_a\tau_C(k)_{ia}\,R_U(J_tk)F_{0,a}=F_{0,i}
                   \quad(k\in K_T, 0\leq i<m),                \tag{HPT.8}
\]
where $\tau_C$ is the specified coefficient map applied to
$\tau_0$.  The inverse transpose in ITA9 is exactly what produces
this column equation; no dual of $\tau_0$ is substituted.

Let $N\subset W^m$ be the space of all columns satisfying
(HPT.8).  It is nonzero and stable under the infinity group,
since finite and infinite adelic factors commute.  Every selected
infinity factor on each of its components has the given dual
highest weight.  The source embeddings index the infinity places
bijectively.  LTP13 proves injectivity.  For surjectivity, at each
selected split place $w\mid v\mid\ell$ their number is
$[L_w:\mathbb Q_\ell]=[L_v^+:\mathbb Q_\ell]$: characteristic
zero gives exactly that many embeddings of the local field,
all continuous for its unique extended absolute value.
Summing over $v\mid\ell$ gives $[L^+:\mathbb Q]$, by the
finite algebra decomposition of $L^+\otimes\mathbb Q_\ell$.
This is also the number of infinite places of the CM field $L$.
The injective map between these finite sets is therefore onto.
LTP13's explicit dual identification,
and the successive-Hom argument of ITA11, therefore give a
nonzero intertwiner
\[
       \Phi:D^*\longrightarrow N,\qquad
       D=\bigotimes_\tau\nabla_{\lambda_\tau},\qquad
       R_U(z)\Phi(\ell)=\Phi(\ell\circ\xi_\infty(z^{-1})).       \tag{HPT.9}
\]
More explicitly, take a nonzero copy of the required irreducible
dual at the first factor; its Hom multiplicity space remains
isotypic for every commuting remaining factor.  Repeating over
the finite set and evaluating the tensor of these Hom maps gives
(HPT.9).  At each step a nonzero map out of an irreducible module
is injective, so the final map is nonzero.  This also proves that
no finite-factor condition was lost.

Apply the actual inverse scalarization ADP8--ADP9 to each
component of $\Phi$, using the specified basis $e$.  It gives
coefficient-valued functions $F_i$ in the original $\mathrm{AF}(\xi)$
such that every scalarization of $F_i$ is $\Phi(\ell)_i$.
It remains to check the \emph{actual} $U_Q$ relation.  For
$u\in U_Q$, lift $((u: G^a),1)$ to a finite adelic element $z$
using the surjectivity in ADP9.  Coordinate by coordinate,
\[
                      z=J_t(\mathrm{toKT}(u))\,k_0,
                                      \qquad k_0\in K_0.        \tag{HPT.10}
\]
Indeed the selected $\ell$-coordinates of $z$ are one; at type
places the section removes precisely its type coordinate;
outside $V.S'$ the remaining coordinate is integral; at $v_0$
it lies in the original $K_c(1)$; and at auxiliary coordinates
the $U_Q$ condition makes it one.  These are exactly the
defining coordinates of $K_0$.

The translation identity ADP7, the trivial infinity coordinate
of $z$, equation (HPT.8), and the $K_0$ invariance give, after
every scalarization,
\[
       \sum_a\tau_C(\mathrm{toKT}(u))_{ia}
                     F_a(g((u: G^a),1))=F_i(g).               \tag{HPT.11}
\]
Scalarization is injective by ADP7, so (HPT.11) is an equality
of the original coefficient functions.  All their scalarizations
belong to $\sigma$.  Thus $F\in Y\mathrm{Of}$ by its literal
definition.  If $F=0$, all scalarizations and hence $\Phi$ would
vanish, contradicting its construction.  This proves (HPT.7).
\end{proof}
% END HPT11

% BEGIN HPT12
\paragraph{HPT12. A joint eigenfunctional produces a joint eigenvector.}
\label{claim:HPT12}
Let a commuting family $(T_i)$ act on a finite-dimensional vector
space over an algebraically closed field.  If a nonzero functional
$L$ satisfies $LT_i=\mu_iL$ for every $i$, there is a nonzero vector
$v$ with $T_iv=\mu_iv$ for every $i$.
\begin{proof}
The algebra generated by the family is finite-dimensional and
commutative, even if the indexing set is infinite.  Decompose the
space into its finitely many simultaneous generalized-character
spaces.  On a space of character $\chi\ne\mu$, choose $i$ with
$\chi(T_i)\ne\mu_i$.  A power of $T_i-\chi(T_i)$ vanishes on
that space, while applying $L$ multiplies by the nonzero scalar
$(\mu_i-\chi(T_i))$ to the same power.  Hence $L$ vanishes there.
The $\mu$ space is consequently nonzero.  On it the commuting
operators $T_i-\mu_i$ generate the nilpotent maximal ideal of a
finite local algebra.  The last nonzero term of its powers applied
to the space is killed by that ideal.  Any nonzero vector there
is the required simultaneous eigenvector.
\end{proof}
% END HPT12

% BEGIN HPT13
\paragraph{HPT13. A principal series has the exact local ordinary eigenline.}
\label{claim:HPT13}
Suppose each smooth $\omega_i$ is trivial on $U_c(c_0+1)$.
In $\mathrm{PS}(\omega)$ there is a nonzero $I_1(c_0)$-fixed
vector on which the compact diagonal torus acts by
$\prod_i\omega_i(t_i)$ and, for every $0\leq j\leq n$,
the original positive double-coset operator at
$a_j=\operatorname{diag}(\varpi I_j,I_{n-j})$ acts by
$\mathrm{psEig}(\omega,j)$.
\begin{proof}
Realize the principal series by functions $f(bg)=\chi(b)f(g)$,
where
\[
       \chi(t)=\delta_B(t)^{1/2}\prod_i\omega_i(t_i)|t_i|^{(n-1)/2}.
\]
The function supported on $BI_1(c_0)$ with value one at the
identity is well defined: $\chi$ is trivial on $B\cap I_1(c_0)$,
whose diagonal entries lie in $U_c(c_0+1)$.  Thus evaluation at
the identity is a nonzero functional on the $I_1(c_0)$-fixed
space.  That space is finite-dimensional.  Iwasawa decomposition
reduces its functions to finitely many coordinates on
$K/I_1(c_0)$, and $I_1(c_0)$ is open in the compact group $K$.

The compact torus and all positive operators preserve this fixed
space.  OST5 and OST26 give their actual coset representatives
and their pairwise commutation.  Evaluation at the identity has
torus eigenvalue $\chi(t)=\prod_i\omega_i(t_i)$ and positive
operator eigenvalue $[N:aNa^{-1}]\chi(a)$: all representatives
lie in $N$, on which the inducing character is one.
For $a=a_j$, the index and character are
\[
 [N:a_jNa_j^{-1}]=q^{j(n-j)},\qquad
 \chi(a_j)=\prod_{i<j}\omega_i(\varpi)q^{-(n-1-i)}.
\]
Since $\sum_{i<j}(n-1-i)=j(n-j)+\sum_{i<j}i$, their product is
exactly $\mathrm{psEig}(\omega,j)$.  Apply HPT12 to obtain a
simultaneous eigenvector, including all compact-torus operators.
In the source $\mathrm{IsInvOn}(\phi,\omega)$ notation this
torus character is $\phi^{-1}$.  The depth is $c_0+1$ as stated;
no increase in depth or restriction on the residue prime is made.
\end{proof}
% END HPT13

% BEGIN HPT14
\paragraph{HPT14. Simultaneous local eigenlines give an actual global ordinary vector.}
\label{claim:HPT14}
Let $H$ be the original Hecke setting, let $\mathrm{IsSpec}$ be
its actual specialization at an arithmetic point $x_A:A\to
\overline{\mathbb Q}_\ell$, with data $\lambda,\xi,c_0,\mathrm{sp}$,
and let $Y$ be an original nonzero $\mathrm{PartAltAlt}$.
Suppose its selected local components are $\mathrm{PS}(\omega_v)$,
the auxiliary components have nonzero $K_q$-fixed vectors, and
\[
 \omega_{v,i}(u)\,\iota\mathrm{finPart}_{v,i}(u)=1,\qquad
 \left|\iota^{-1}\bigl(\mathrm{lowVal}(v,j+1)^{-1}
                 \mathrm{psEig}(\omega_v,j+1)\bigr)\right|_\ell=1.
                                                               \tag{HPT.12}
\]
Then $Y$ contains a nonzero vector in the actual specialized
ordinary space at $K,c_0$, with all spherical eigenvalues $Y.a$
and all normalized ordinary eigenvalues
\[
             \alpha_{v,j}=\mathrm{lowVal}(v,j+1)^{-1}
                                     \mathrm{psEig}(\omega_v,j+1).
                                                               \tag{HPT.13}
\]
\begin{proof}
The triviality field of $\mathrm{IsSpec}$ and (HPT.12) make each
$\omega_{v,i}$ trivial on $U_c(c_0+1)$.  Apply HPT13 locally.
To impose these conditions simultaneously, start with $Y$ and
process the finite set of selected and auxiliary factors.  The
current subspace is stable and isotypic for every remaining factor:
actions at different selected coordinates commute; selected and
auxiliary actions commute by the literal right-translation
formula; two auxiliary actions commute by the original disjoint
coordinate condition in $H.\mathrm{Ax}$.  Isotypy means that any
nonzero stable subspace contains an embedded copy of its given
irreducible local representation (LTP3).  Map the nonzero vector
of HPT13, or the supplied auxiliary fixed vector, into this copy.
Intersecting with that factor's conditions therefore leaves a
nonzero subspace and preserves all previous conditions.  Finite
induction gives one nonzero $F$ satisfying every condition.

The local $N$ and lower-unipotent fixed conditions combine into
the corresponding product conditions.  The $U_Q$ relation already
holds on $Y$, and the auxiliary fixed conditions extend it to
$H.\mathrm{lev}(K)$.  The torus character in HPT13 and (HPT.12)
is exactly the finite-part character in $H.\mathrm{ClAt}$.
Thus $F\in H.\mathrm{ClAt}(K,c_0)$.  OST26--OST27 compute its
positive Hecke eigenvalues and give (HPT.13).

For ordinary projection one must use the original indexing:
the element $z$ is the product of $a_j$ for $j=0,\ldots,n-1$,
whereas the displayed generators are indexed by $j+1=1,\ldots,n$.
The normalized $z$ eigenvalue on $F$ is
\[
              \prod_v\prod_{j=0}^{n-1}
                  \mathrm{lowVal}(v,j)^{-1}\mathrm{psEig}(\omega_v,j).
\]
The $j=0$ factor is one and all remaining factors have norm one
by (HPT.12).  Hence $F$ belongs to the original unit-eigenvalue
ordinary part.  Its spherical eigenvalues are those of $Y$
because $F\in Y$.  This proves the complete simultaneous assertion
on the original function space.
\end{proof}
% END HPT14

% BEGIN HPT15
\paragraph{HPT15. Exact unit and residue normalization at the original point.}
\label{claim:HPT15}
Let $x:B\to\overline{\mathbb Q}_\ell$ be arithmetic of weight
$h=\mathrm{htW}(\lambda)$, and suppose (HPT.3) holds for its
ordered characters.  Then the first identity in (HPT.12) holds,
and every $\iota^{-1}\alpha_{v,j}$ is a unit value of $x$,
congruent to the specified ordinary residual eigenvalue
$R_H.a(v,j)$ modulo elements of norm less than one.
If $Y$ has the same Frobenius polynomials as $\rho$ with
$x.\mathrm{IsRep}(\rho)$ at the original spherical places, its
spherical eigenvalues have the analogous residual congruence.
\begin{proof}
For a local unit $u$, choose a Weil element mapping under $a$
to $u$.  Its degree is zero.  Insert (HPT.3) in the original
$\theta$ relation and use the exact weight map
$\Lambda.\mathrm{toPair}(B.u)$; (HPT.2) cancels the $i$ copies
of the cyclotomic algebraic factor against
$h_i=\mathrm{lowW}_i+i$.  The remaining equality is
$\omega_{v,i}(u)\iota\mathrm{finPart}_{v,i}(u)=1$.
This calculation is HEP4--HEP6 in the reverse direction and uses
the given finite part, not a new choice of it.

Choose a Weil lift $s_\varpi$ with $a(s_\varpi)=\varpi$ and
degree $-1$, and let $g$ be its actual decomposition-group image.
The unit of $B$ relevant to the $(j+1)$-st operator is
\[
             b_{v,j}=\prod_{i=0}^{j}
                          B.u.\psi_{v,i}(g)\,\varepsilon(g)^i.
\]
Using (HPT.2)--(HPT.3) and (HPT.1) term by term gives
\[
                          x(b_{v,j})=\iota^{-1}\alpha_{v,j}.
                                                               \tag{HPT.14}
\]
Indeed the algebraic exponents are $-h_i+i=-\mathrm{lowW}_i$
and the numerical factors are $q^{-i}$.  ASP2 says that the value
of a unit and its inverse both have norm at most one, so this
value has norm one.  The residual character field of $B.u$ and
the literal $R_H.\mathrm{res.unif}$ identify its residue with
$R_H.a(v,j)$.  ASP2 then gives the strict norm inequality.

At a spherical place $w\nmid\ell$, compare the coefficient of
$X^{n-j}$ in the common Frobenius polynomial.  With the source
Hecke convention it is
$(-1)^j q_w^{j(j-1)/2}Y.a(w,j)$.  The residual determinant law
and $R_H.\mathrm{res.frob}$ give the same coefficient with
$R_H.a(w,j-1)$ for $1\leq j\leq n$.  The factor
$q_w^{j(j-1)/2}$ has $\ell$-adic norm one; division preserves
the strict congruence.  This proves every residual congruence
needed for the actual residual block.
\end{proof}
% END HPT15

% BEGIN HPT16
\paragraph{HPT16. The actual ordinary vector gives a point of the original residual Hecke algebra.}
\label{claim:HPT16}
In HPT14 suppose all its eigenvalues are congruent to the prescribed
residual function $a$.  Then the original residual module
$H.MK(\theta,a,K)$ is nonzero, and there is an $O$-algebra point
$y:H.TK(\theta,a,K)\to\overline{\mathbb Q}_\ell$ extending $x_A$
and taking every generator to the eigenvalue in HPT14, transported
by $\iota^{-1}$.  This is the entire literal
\texttt{core\allowbreak Up\_\allowbreak alt} conclusion.
\begin{proof}
Let $N$ be the original ordinary module before taking the residual
block.  OHA4 and OHA8 make $N$ finite free and its given operators
commuting.  The span and independence fields of ASP13 identify
$N\otimes_{A,\iota x_A}\mathbb C$ with the specialized ordinary
space.  The nonzero vector $F$ of HPT14 therefore determines an
eigenline in this scalar extension of the actual $N$.
Its character on the actual algebra $T\subset\operatorname{End}_A(N)$
exists because every generator preserves this line.  It is
well defined on $T$: any relation acting as zero on $N$ acts
as zero after scalar extension and hence on that line.
Transport its values through $\iota^{-1}$ to obtain $\chi$.
It extends $x_A$ and takes the generators to precisely the stated
eigenvalues.  All its values are integral: $T$ is generated over
$A$ by these integral eigenvalues, and ASP2 bounds $x_A(A)$ by one.
Reduction of $\chi$ kills $\mathfrak m_A$ and every $T_i-a_i$.

OHA5--OHA6 decompose $T$ into its actual finite collection of
local idempotent factors.  The ideal
$(\mathfrak m_A,T_i-a_i)_i$ singles out at most one factor; the
reduction of $\chi$ shows that this factor exists.  Its idempotent
$e$ maps to one under $\chi$, whereas all other idempotents map
to zero.  Consequently $eF=F\ne0$ and $eN=H.MK(\theta,a,K)$
is nonzero.  The action of $eT$ on $eN$ is faithful: it is zero
on $(1-e)N$ already, so an element killing $eN$ kills $N$ and
is zero in the original operator algebra.  Thus $eT$ is exactly
the original block algebra $H.TK$.  Restrict $\chi$ to it to get
$y$.  This proves the coefficient-map identity and all spherical
and normalized ordinary generator identities simultaneously.
\end{proof}
% END HPT16

% BEGIN HPT17
\paragraph{HPT17. The full original ordinary return theorem in every positive rank.}
\label{claim:HPT17}
Retain the original data with $n\geq1$ and $V.S'=s.S$.
Suppose $x$ is arithmetic with $h=\mathrm{htW}(\lambda)$,
$\rho$ is continuous semisimple and $x.\mathrm{IsRep}(\rho)$.
Suppose an actual automorphic $\sigma$ satisfies all of the
following original return inputs: $\mathrm{SphEq}(\iota,\sigma,\rho)$;
infinity weights $\lambda^\vee$; unramifiedness outside $s.S$;
a nonzero $K_c(1)$-fixed vector in its actual $v_0$ component;
at each selected $\ell$-place an actual component
$\mathrm{PS}(\omega_v)$ satisfying (HPT.3) for $x.\psi_{v,i}$;
and $\mathrm{StrongEquivalenceHida}$ at all type places, with
its quantifier over every actual local component and every
$\mathrm{IsRecT}$.  Then $J\subseteq\ker x$.
\begin{proof}
AFS8 translates the arithmetic weight to the original admitted
$\lambda$.  Use ASP13 to choose the actual specialization map,
algebraic basis, depth and coefficient structure.  HPT2 and the
strong comparison supply all parameter types at the type places,
for every Teichm\"uller family.  HPT11 constructs a nonzero
$Y\mathrm{Of}$ with this basis.  UPA16 now produces the original
part $Y$ and its continuous semisimple frame representation $R_Y$,
with all local components still those of this same $\sigma$.

At each original spherical place $w\in H.W_p$, the hypotheses
make $\sigma$ spherical.  Choose an actual local component and
its Satake parameters.  $\mathrm{SphEq}$ computes the Frobenius
polynomial of $\rho$ from those parameters; the local comparison
of $Y$ and GUA14 compute the polynomial of $R_Y\circ p$ from
the same parameters.  They agree for every original prolongation
and Frobenius witness.  HEP1--HEP2, using the given full determinant
law of $x$ and the original density of $H.W_p$, give the full
frame equality
\[
                              x.D=\det R_Y.                    \tag{HPT.15}
\]
This is an equality of polynomial laws, not only of Frobenius
characteristic polynomials.  HPT15 gives every spherical residual
congruence and every normalized ordinary unit and congruence.
Actual local-component uniqueness identifies $Y$'s selected
components with the given principal series.

There are no auxiliary factors here.  Indeed $Q.\mathrm{range}$
would put the place of any $q\in Q$ in $V.S'=s.S$, contrary
to $Q.\mathrm{notS}$.  Thus the auxiliary fixed-vector requirement
of HPT14 is vacuous.  Apply HPT14--HPT16 to get the nonzero
original residual module and its point $y$ with the exact
weight and eigenvalues.

In the zero-module branch of $\mathrm{IsJ}$ this nonzero module
is a contradiction.  In the other branch take its actual
surjection $f:B\twoheadrightarrow H.TK$ with $J=\ker f$ and
the stated determinant, ordinary-character and weight comparisons.
Apply HEP7 to a nonzero vector of $Y$, the just constructed
Hecke point, and its spherical and ordinary eigenvalue identities.
It gives
\[
                (yf).D=\det R_Y,\qquad
       \mathrm{OrdChar}(\omega_{v,i},h_i,(yf).\psi_{v,i}).       \tag{HPT.16}
\]
The finite-part premise in that application is exactly HPT15,
since $y$ has the weight $x|_\Lambda$.  Equations
(HPT.15)--(HPT.16) identify the two full laws.  HPT3 identifies
all ordered characters.  HPT1 therefore gives $x=yf$ as maps on
the entire original ring $B$.  For $b\in J=\ker f$ we obtain
$x(b)=y(f(b))=0$, proving the conclusion.  This is the literal
\texttt{up\allowbreak U\_\allowbreak w\_\allowbreak alt} endpoint with all of its source inputs,
proved here on the actual module and actual ideal.
\end{proof}
% END HPT17

% BEGIN HPT18
\paragraph{HPT18. The complete low-rank point theorem.}
\label{claim:HPT18}
For every original point $x$, weight $h$ and representation
$\rho:G_L\to\mathrm{GL}_n(\overline{\mathbb Q}_\ell)$ with
$2\leq n\leq4$, assume
\[
 \begin{gathered}
 x.\mathrm{Arith}(h),\qquad x.\mathrm{IsRep}(\rho),\qquad
 \mathrm{Continuous}(\rho),\qquad \mathrm{IsAutCsdA}(\iota,\rho),\\
 \rho|_{I_{s.A(v_0)}}=1,\qquad
 \mathrm{HasFlag}(\rho|_{D_{s.A(v)}},x.\psi_v)
                  \quad\text{for every selected }v\mid\ell,\\
 x.\psi_{v,i}\ne x.\psi_{v,j}\varepsilon|_{D_{s.A(v)}}
                          \quad\text{for every }v,i,j .
 \end{gathered}                                                 \tag{HPT.17}
\]
Then $J\subseteq\ker x$.  Hence the same original $J$ satisfies
the entire \texttt{Point\allowbreak Theorem\allowbreak Base\allowbreak Object} field.
\begin{proof}
Use AFS8 to obtain the admitted dominant $\lambda$ with
$h_i=\lambda(n-1-i)+i$.  The $\mathrm{IsAutCsdA}$ witness
supplies a regular algebraic conjugate-self-dual cusp $\pi$,
the attachment and semisimplicity of $\rho$, and its $\mathrm{WitA}$
guard.  HPT8 gives the full compatible local parameters and
the selected ordinary principal series for this supplied
$\rho$.  Outside $s.S=V.S'$, the map $p$ kills inertia.
HPT2 makes $\rho$ unramified there; these places are away from
$\ell$.  HPT7 makes every local component of $\pi$ there
spherical, and the original Satake theorem supplies its parameters.

Apply the retained $\mathrm{DescWeak}$ with $E=L$, prime $\ell$
and exception set $s.S$.  The original Hida setting supplies
finite unramifiedness of $L/L^+$, finiteness of $s.S$, splitting
of every place above $s.S$, and inclusion of every place above
$\ell$.  The rank inequalities are exactly (HPT.17).  The
cusp and its Satake condition were just verified.  Thus it gives
an actual $\sigma$, unramified at all nonsplit places, with the
cofinite split comparison.  HPT9 upgrades this to every split
local component and the infinity multiset.

We verify each input of HPT17.  At selected $\ell$-places the
local component is the principal series already obtained in HPT8,
with the same ordered $\omega_i$, and the corresponding Hodge
multiset is $\{\lambda(n-1-i)+i\}_i$.  If $\nu$ denotes the
dominant infinity weight of $\sigma$, the original Hodge/infinity
normalization and HPT9 say that the same multiset is
$\{-\nu_i+i\}_i$.  Both sequences are strictly increasing.
Their minima agree; remove that common term and repeat to get
termwise equality.  Hence $\nu_i=-\lambda(n-1-i)$, exactly
$\lambda^\vee$.  This holds at every source embedding, including
its original local embedding identification.

Outside $s.S$, nonsplit unramifiedness is the descent output.
At split places it follows from HPT9, the spherical components
of $\pi$, and LTP20--LTP21.  At $v_0$ the residue characteristic
is different from $\ell$ by the original Hida-setting condition.
The full-inertia hypothesis in (HPT.17) and HPT7 make its
$\pi$ component spherical.  Its nonzero maximal-compact fixed
vector is fixed by the original subgroup $K_c(1)$.  HPT9
transports this vector to the actual $\sigma$ component.

At every split away-$\ell$ place, the full parameter comparison
of HPT8 and the same-component statement of HPT9 give the
spherical Frobenius law GUA14, which is precisely
$\mathrm{SphEq}$.  At every type place, all actual components
match by HPT9 and all correspondences match by HPT8, giving
the literal $\mathrm{StrongEquivalenceHida}$, with no restriction
to one correspondence or one choice of component.  All hypotheses
of HPT17 have now been proved from (HPT.17) and the original
data.  Its conclusion is $J\subseteq\ker x$, as required.
In particular this containment has not been used in producing
$\sigma$, the local vectors, or the point $y$.
\end{proof}
% END HPT18


\section{Finite-order twisting preserves the same ordinary ideal}
\label{sec:hida-twisting-v45}

This chapter proves the literal $B.\mathrm{TwistOk}(J)$, in every
original positive rank.  It reuses HPT17's ordinary return theorem;
the rank-$2$--$4$ descent input in HPT18 is not used.  Fix all
original Hida, coefficient, level, type and Hecke data as in HPT,
with $V.S'=s.S$, $n\geq1$, $\ell$ odd and
$R_H.\mathrm{IsJ}(B,K,J)$.  The isomorphism $\iota$ remains fixed.
We retain the same original global Artin, local classification,
local correspondence and period/monodromy foundations as before.
Every occurrence of a local correspondence means the complete
original $\mathrm{IsRecT}$ quantifier; no new global descent
foundation is introduced.

The exact input consists of two original points $x,x'$, a weight
$h$, and a character $\beta:G_0\to\overline{\mathbb Q}_\ell^\times$
such that
\[
 \begin{gathered}
 J\subseteq\ker x,\quad x.\mathrm{Good}(h),\quad
 \ker\beta\text{ open},\quad \beta^{\ell^k}=1\text{ for some }k,\quad
 \beta\circ\mathrm{conj}=\beta^{-1},\\
 I_u\subseteq\ker\beta\quad(u\in s.S_t),\qquad
 I_0\subseteq\ker\beta,\qquad
 x.D.\mathrm{IsTwist}(\beta,x'.D),\\
 x'.\psi_{v,i}=x.\psi_{v,i}
                 \bigl(\beta\circ p|_{D_{s.A(v)}}\bigr)
                                        \quad\text{for all }v,i.
 \end{gathered}                                                  \tag{HTS.1}
\]
Here $p=V.\mathrm{pr}$, and all inertia groups are the original
frame groups.  Goodness includes arithmeticity and injectivity
of the ordered finite-part characters at each selected place.
The desired conclusion is $J\subseteq\ker x'$ on the same ring.

% BEGIN HTS1
\paragraph{HTS1. An odd-order conjugate-dual character supplies the required unitary character.}
\label{claim:HTS1}
Put $\mu=\beta\circ p$ and $m=\ell^k=2N+1$.  There is a
finite-order Hecke character $\nu$ with Galois character
$\delta=\mu^{N+1}$, and $\kappa=\nu^2$ has Galois character
$\mu$ with weight zero.  The character
$\chi=\nu\circ\det$ on the actual adelic unitary group has
open kernel and is trivial on rational and infinite elements.
At a split local factor its character is $\kappa_w\circ\det$;
at a nonsplit local factor it is $\nu_w\circ\det$.
\begin{proof}
An open kernel makes $\beta$ locally constant, hence continuous;
the same holds after composing with the continuous $p$.
Since $m$ is odd, $\delta^2=\mu^{m+1}=\mu$ and $\delta^m=1$.
It is conjugate-dual because $p$ carries the specified complex
conjugation to the original frame involution.

Choose the original global Artin map.  The continuous character
$\delta$ kills the closure of the commutator subgroup, so factors
through the abelianized Galois group.  Compose this factor with
Artin and $\iota$ to define $\nu$ on ideles.  Its image is finite,
so its kernel is open despite the absence of any continuity
assumption on $\iota$.  The Artin kernel contains rational
ideles, so $\nu$ is an actual Hecke character.  At weight zero
the defining $\mathrm{IsLAdicOf}$ equality is simply
\[
         \delta(g)=\iota^{-1}\nu(a)
                       \quad\text{if }\mathrm{Art}(a)=\bar g.
                                                               \tag{HTS.2}
\]
This proves the stated attachment, its order bound, and the
attachment of $\kappa$ to $\mu$.  Uniqueness follows by applying
(HTS.2) to every idele, since every element of the abelianized
Galois group has a lift to the Galois group.

The conjugate-dual identity gives $\nu(\bar a)=\nu(a)^{-1}$.
One may check this without presupposing functoriality of the
chosen Artin map: outside the union of the two finite conductor
sets, conjugating a Frobenius lift and its prime preserves its
residue power and gives the claimed inverse values.  Finite
characters of the idele class group correspond to continuous
finite characters of the Artin quotient.  Indeed their open
kernels contain rational ideles and the connected infinity
factor, hence the closure specifying the original Artin kernel;
the quotient map factors them.  Chebotarev equality at the above
cofinite Frobenius set then identifies these characters, proving
the identity on all ideles.

The determinant character $\chi$ has open kernel because the
determinant is continuous and $\nu$ has open kernel.  On a
rational matrix its determinant is a rational idele, so it is
trivial.  It is trivial on infinity: $L$ is totally complex,
and every infinity idele has an $m$-th root with finite part one;
applying the order-$m$ character gives value one.
At a split pair $w,\bar w$, the unitary equation says that the
conjugate of the determinant at $\bar w$ is the inverse of the
determinant at $w$.  The two local $\nu$ values are consequently
equal, and their product is $\nu_w^2=\kappa_w$.  At a nonsplit
place there is one coordinate, giving $\nu_w$ itself.
This proves the precise split/nonsplit local formula on the
actual adelic group, including the square in the split case.
\end{proof}
% END HTS1

% BEGIN HTS2
\paragraph{HTS2. Multiplication by the actual determinant character twists automorphic spaces.}
\label{claim:HTS2}
For an actual irreducible automorphic subspace $\sigma$, set
\[
          M_\chi f(g)=\chi(g)f(g),\qquad \sigma'=M_\chi\sigma.
                                                               \tag{HTS.3}
\]
Then $\sigma'$ is actual irreducible automorphic, with the same
infinity weights.  At every split place its actual components
are precisely the twists of those of $\sigma$ by
$\kappa_w\circ\det$.  Wherever $\mu$ is trivial on inertia,
unramifiedness is preserved; so is every fixed vector for a
subgroup of the local maximal compact, in particular $K_c(1)$.
\begin{proof}
$M_\chi$ is a linear isomorphism with inverse $M_{\chi^{-1}}$,
and direct substitution gives
\[
          R_U(g)M_\chi f=\chi(g)M_\chi R_U(g)f.                 \tag{HTS.4}
\]
Continuity and left rational invariance follow from HTS1.
Intersect an open stabilizer of $f$ with $\ker\chi$ to get
an open stabilizer of $M_\chi f$.  The infinity orbit span
is the image of the original finite-dimensional span, so it
is finite-dimensional.  Thus the original automorphic-form
conditions all hold.  Equation (HTS.4) preserves invariant
subspaces in both directions, proving irreducibility and
nonvanishing.  Triviality of $\chi$ at infinity preserves
every infinity intertwiner and hence its weight.

Compose any local-component embedding with $M_\chi$.
Equation (HTS.4) and the split formula in HTS1 give exactly
the twisted local representation.  Conversely use the inverse
map and inverse twist.  This proves the statement for every
actual component, not just a chosen family.  A finite-order
Galois character is trivial on inertia exactly when its
weight-zero Hecke character is trivial on local units, by
local Artin surjectivity.  If $\mu$ is trivial on inertia,
so is its power $\delta$; thus both $\kappa_w$ and $\nu_w$
are unramified.  The determinant of a maximal-compact matrix
is a unit, so the local character of HTS1 is one on the
compact subgroup in question.  The same nonzero vector
remains fixed under the twisted action, at split and nonsplit
places alike.
\end{proof}
% END HTS2

% BEGIN HTS3
\paragraph{HTS3. The two exact guards for local twisting.}
\label{claim:HTS3}
Two local situations provide precisely the original
$\mathrm{Good}(\tau)$ and $\mathrm{Gd}(\mathrm{rc}(\mathrm{algNorm}\tau))$
guards: first, a prescribed zero-monodromy inertia type given by
restrictions of global characters; second, an actual twist
$\tau\otimes(\kappa_w\circ\det)$ which is spherical.
These statements hold for every permitted correspondence and
do not assert $\mathrm{PlainTor}$ for a nongeneric component.
\begin{proof}
For the first situation, write the parameter using the original
Langlands multisegment law.  The type hypothesis gives zero
monodromy and the inertia polynomial of restrictions of global
characters.  For a Weil element $s$ and an inertia element $t$,
their commutator belongs to inertia and each such global character
is one on it.  Hence its total characteristic polynomial is
$(X-1)^n$.  The polynomial of any irreducible Weil constituent
of any supercuspidal support factor divides this polynomial.
Its inertia action is semisimple, so the commutator acts as the
identity.  Thus inertia commutes with the entire irreducible
Weil image and is scalar by Schur's lemma.  An eigenline for
Frobenius is then stable under inertia and Frobenius, which
generate the Weil group.  Irreducibility forces dimension one.
The rank equality of the correspondence makes every cuspidal
support factor rank one.  This is enough for the original
$\mathrm{Good}$ disjunction, including repeated or linked
characters.  The support product law makes the polynomial on
every commutator of two Weil elements equal to $(X-1)^n$;
the type already gives zero monodromy.  These are exactly
the two fields of $\mathrm{Gd}$.

For the second situation, GUA12 identifies the spherical
component's support as unramified characters and its parameter
as a sum of rank-one unramified characters with zero monodromy.
Untwisting its support by the single character $\kappa_w$
again gives rank-one support, hence $\mathrm{Good}(\tau)$.
The twist law of $\mathrm{IsRecT}$ applies because this guard
holds.  It identifies the untwisted parameter with the same
sum tensored by a rank-one Weil character.  Such a character
is one on every commutator and does not change whether the
monodromy is zero.  The untwisted parameter is therefore
$\mathrm{Gd}$ at every correspondence.  Neither argument
uses genericity or a stronger all-local compatibility theorem.
\end{proof}
% END HTS3

% BEGIN HTS4
\paragraph{HTS4. Finite-character twisting respects the full away local comparison.}
\label{claim:HTS4}
Suppose $w\nmid\ell$, $R$ is continuous, $\mu$ is the finite
character of HTS1, $\tau$ satisfies the original $\mathrm{Good}$
guard, and $\tau'\simeq\tau\otimes(\kappa_w\circ\det)$.
For every original $\mathrm{IsRecT}$,
\[
 \mathrm{IsWDOf}(\iota,R,\mathrm{rc}(\mathrm{algNorm}\tau))
 \ \Longrightarrow\
 \mathrm{IsWDOf}(\iota,R\mu,\mathrm{rc}(\mathrm{algNorm}\tau')).
                                                               \tag{HTS.5}
\]
Here $(R\mu)(g)=\mu(g)R(g)$ on the same vector space.
\begin{proof}
Choose a finite coefficient field containing the representation's
compact image and the finite image of $\mu$, as in ADK1.
Intersect the open subgroup used for the original away local
monodromy witness with $\ker\mu$.  On it $R\mu=R$, so the
same nilpotent monodromy operator works.  On the Weil group
the resulting raw action is the old action multiplied by
$\iota\mu(s)$.  This follows immediately from the defining
formula $R(s)\exp(-t_\ell(s)N)$.  A scalar factor commutes with
the action and its semisimplification correction; hence the
Frobenius-semisimplified parameter is the original one tensored
by this rank-one character.  The same coefficient lattice,
after the harmless finite enlargement, supplies an actual
original $\mathrm{IsWDOf}$ witness.

Local-global Artin compatibility and (HTS.2) identify
$\iota\mu(s)=\kappa_w(a(s))$.  OFL1 identifies this $a$ with
the map in the chosen correspondence.  Its twist law applies
under $\mathrm{Good}(\tau)$; the common algebraic normalizing
norm twist commutes with the character twist.  It gives exactly
the parameter on the right of (HTS.5).  The choices were made
for an arbitrary correspondence and prolongation, so all those
quantifiers remain intact.
\end{proof}
% END HTS4

% BEGIN HTS5
\paragraph{HTS5. Good finite parts give a full principal series and its exact twist.}
\label{claim:HTS5}
Suppose $\tau$ is a smooth irreducible subquotient of $\mathrm{PS}(\omega)$,
$\mathrm{OrdChar}(\omega_i,h_i,\psi_i)$ holds, and the finite
parts $\phi_i$ of these characters are pairwise distinct.
Then $\tau\simeq\mathrm{PS}(\omega)$.  If $\tau'$ is its twist
by $\kappa_w\circ\det$, then
\[
 \tau'\simeq\mathrm{PS}(\omega'),\qquad
 \omega'_i=\omega_i\kappa_w,\qquad
 \mathrm{OrdChar}(\omega'_i,h_i,\psi_i\mu|_{D_w}).               \tag{HTS.6}
\]
The latter characters are locally algebraic with the same weights.
\begin{proof}
By HPT3 the restrictions of $\omega_i$ to units are inverses,
through $\iota$, of the finite parts on Artin images of inertia.
Surjectivity of local Artin on units implies that the restrictions
are pairwise distinct.  No two $\omega_i$ therefore differ by
any unramified norm twist, including twists $-1,0,1$.
The original totally disjoint product theorem and induction
in stages make $\mathrm{PS}(\omega)$ irreducible.  Its nonzero
irreducible subquotient $\tau$ is isomorphic to it.

The explicit inducing-function map
$f(g)\mapsto\kappa_w(\det g)f(g)$ identifies the twist of the
normalized induction with the normalized induction of
$(\omega_i\kappa_w)|\,\cdot\,|^{(n-1)/2}$.
This proves the actual intertwiner in (HTS.6).  The weight-zero
Artin identity says $\mu(s)=\iota^{-1}\kappa_w(a(s))$.
Multiply it by (HPT.3) to get the last assertion of (HTS.6),
including smoothness and continuity.  Equivalently the finite
part becomes $\phi_i\mu|_{I_w}^{-1}$: its kernel contains the
intersection of the two open kernels, and the product of this
finite part with $\psi_i\mu$ has the same algebraic expression.
This directly checks every field of local algebraicity and
does not require a new choice of the algebraic weight.
\end{proof}
% END HTS5

% BEGIN HTS6
\paragraph{HTS6. Extracting an actual automorphic part with all ordinary subquotients.}
\label{claim:HTS6}
For $J\subseteq\ker x$ and $x.\mathrm{Good}(h)$, there are an
admitted $\lambda$ with $h=\mathrm{htW}(\lambda)$ and one actual
part $Y$ arising from a nonzero
$Y\mathrm{Of}(Q,\iota,\tau_0,e,\xi,\sigma)$, such that
$x.D=\det Y.r$ and every selected local component of $\sigma$
is a subquotient of $\mathrm{PS}(\omega_v)$ with
$\mathrm{OrdChar}(\omega_{v,i},h_i,x.\psi_{v,i})$.
The same $\sigma$ has the full split $\mathrm{LocGlobT}$
comparison with $Y.r\circ p$, its original infinity weights,
unramifiedness outside $V.S'$, and its original $K_c(1)$ vector.
\begin{proof}
Use the part-dependent predicate which says that $Y$ is obtained
from such a $\sigma,e,\xi$ and has its actual local components
and full split compatibility.  This predicate satisfies the
literal $\mathrm{AutYAlt}$ spanning condition: ADP12 decomposes
each original coefficient function into its $Y\mathrm{Of}$
pieces; discard the zero pieces and apply UPA16 to each nonzero
piece, retaining exactly its same $\sigma,e,\xi$ and local
components.  The integral type has open kernel by ITA12 and
the specified injective coefficient map, as required by ADP12.

The top-ideal branch has no such point.  Otherwise HED1 factors
$x$ through the original Hecke quotient.  AFS8 translates its
weight, and HAW10, with the just proved part-dependent predicate,
gives this actual $Y$, its full determinant equality and all
ordinary subquotients.  Naturality HED2 pulls the equality back
to the identical $x.D$, and the exact IsJ character maps make
the characters precisely $x.\psi$, in their given order.
Finally LTP12, LTP14 and LTP15 apply to this same nonzero
$Y\mathrm{Of}$ to give unramifiedness, infinity weights and
the depth-one fixed vector.  LTP17--LTP18 also give its full
type statement for every Teichm\"uller family.  This explains
why the arbitrary-part version of HAW10 is used: a representation
with a cofinite attachment alone would not supply these data.
\end{proof}
% END HTS6

% BEGIN HTS7
\paragraph{HTS7. The twisted point and every input of the ordinary return.}
\label{claim:HTS7}
For the data in (HTS.1), put $R=Y.r\circ p$ using HTS6 and
$R'=R\mu$ on the same vector space.  Then $x'$ is arithmetic
of weight $h$, $R'$ is continuous semisimple,
$x'.\mathrm{IsRep}(R')$, and the actual twisted automorphic
representation $\sigma'$ of HTS2 satisfies every input of HPT17
for $x',R',\lambda$.
\begin{proof}
HED5 gives continuity and semisimplicity of $R$.  Scalar twisting
preserves invariant subspaces and invariant complements, so
$R'$ is semisimple, and its continuity follows from that of
$R,\mu$.  The last line of (HTS.1) and the finite-part formula
in HTS5 show local algebraicity for $x'.\psi$ with the same
strictly increasing weights.  Thus $x'.\mathrm{Arith}(h)$.

The determinant twist relation is a full polynomial-law identity.
Evaluate it on an arbitrary scalar algebra and a finite sum
$\sum a_g g$.  It replaces each term by
$a_g\beta(g)g$; the matrix under $Y.r$ is therefore
$\sum a_g\beta(g)Y.r(g)$.  Pulling along $p$ gives exactly
the matrix of $R'$ on that sum.  Since $x.D=\det Y.r$,
this proves the full $x'.\mathrm{IsRep}(R')$, including scalar
algebras with nilpotents.

HTS6 gives actual components of $\sigma$ with the original
ordinary subquotients.  Goodness in (HTS.1) and HTS5 promote
them to full principal series, then twist them to the full
principal series and ordered characters in (HTS.6).  The
character identity for $x'$ in (HTS.1) identifies those characters
exactly.  HTS2 preserves the infinity weights.  Outside $V.S'$
the quotient $p$ kills inertia; hence so do $\mu$ and $\delta$,
and HTS2 preserves unramifiedness.  At $v_0$ the stipulated
$I_0\subseteq\ker\beta$ gives the same conclusion and preserves
the original nonzero $K_c(1)$-fixed vector, since this subgroup
lies in the same maximal compact.

We check the guarded comparisons rather than asserting that
twisting preserves an unconditional local statement.  At a
split away-$\ell$ place choose any spherical component of
$\sigma'$, with its actual Satake multiset.  Untwist it to the
corresponding component $\tau$ of $\sigma$ by HTS2.
HTS3's second case gives $\mathrm{Good}(\tau)$ and the
$\mathrm{Gd}$ parameter for every correspondence.  These
activate the away branch of the full $\mathrm{LocGlobT}$ from
HTS6.  HTS4 twists that comparison, and GUA14 converts it to
the original spherical Frobenius law.  This proves the entire
$\mathrm{SphEq}$ for $\sigma',R'$.

At a type place choose any actual component $\tau'$ of
$\sigma'$ and its untwist $\tau$.  LTP17--LTP18, applied to
the original nonzero part, give the inertia type of
$\mathrm{rc}(\mathrm{algNorm}\tau)$ for every correspondence
and every Teichm\"uller family.  The type characters are
restrictions of the global characters $t_i\circ p$, so HTS3's
first case gives exactly $\mathrm{Good}(\tau)$ and
$\mathrm{Gd}$ there.  The full original local comparison
from HTS6 and HTS4 therefore yield
$\mathrm{IsWDOf}(\iota,R',\mathrm{rc}(\mathrm{algNorm}\tau'))$
for every such component and correspondence.  This is the
whole $\mathrm{StrongEquivalenceHida}$, with no omitted guard.
All inputs of HPT17 have now been verified on the same objects.
\end{proof}
% END HTS7

% BEGIN HTS8
\paragraph{HTS8. The literal original TwistOk endpoint.}
\label{claim:HTS8}
For the original data in every positive rank, the same ideal
$J$ satisfies $B.\mathrm{TwistOk}(J)$.
\begin{proof}
Take arbitrary $x,x',h,\beta$ satisfying all the original
hypotheses (HTS.1), including the type-inertia and depth-one
conditions.  HTS1--HTS7 construct the actual twisted automorphic
representation and verify each hypothesis of HPT17 for $x'$.
Apply that all-rank return theorem to conclude
$J\subseteq\ker x'$.  The argument uses $\ell$ odd only to
write $\ell^k=2N+1$ and take its unique square-root power;
it includes $k=0$ and $\ell=3$.  The given coefficient DVR,
fixed $\iota$, integral type, level and ideal were never changed.
This proves every quantifier of the original \texttt{Twist\allowbreak Ok};
it makes no low-rank restriction and uses no new descent premise.
\end{proof}
% END HTS8


\section{One actual ordinary module family at every auxiliary level}
\label{sec:hida-level-modules-v45}

\status{HLM1--HLM26 reconstruct the full original \texttt{Level\allowbreak Alt}
producer.  The same Taylor--Wiles base is chosen before all auxiliary
levels, including the empty auxiliary set.  Every module vector is
separated by the original classical points, with their full local
correspondence quantifiers.  The original \(\mathrm{HypothesesH}'\)
and its named local, automorphic, reciprocity and cohomological
foundations remain hypotheses.  No formal theorem is substituted
for the arguments below.}

\subsection{Notation and retained inputs}

Use the original field \(L\), closure \(\Omega\), setting \(s\),
coefficient prime \(\ell\), coefficient DVR \(O=c.O\), rank
\(n\geq1\), and fixed algebraic isomorphism
\(\iota:\overline{\mathbb Q}_\ell\simeq\mathbb C\).
There is no continuity assumption on \(\iota\), no restriction
\(\ell>n\), and no exclusion of \(\ell=3\).
For an original Hecke setting \(H\), a complete Noetherian local
\(O\)-algebra \(A\) in \(\mathcal C_O\), and the given weight
and residual eigenvalues \(\theta,a\), put
\[
 N_K=H.X.\operatorname{Ord}(H.\tau_A,\theta,H.\operatorname{lev}K),
 \qquad M_K=H.\operatorname{MK}(\theta,a,K),
 \qquad T_K=H.\operatorname{TK}(\theta,a,K).
 \tag{HLM.1}
\]
These are the original submodules of one coefficient-valued
function space, namely \(F=H.X.\operatorname{Fn}(A,H.m)\).
An auxiliary level is a tuple of compact open subgroups
\(K_q\subset G_q=\operatorname{GL}_n(L_{H.wq(q)})\);
write \(K=\prod_qK_q\) when no confusion results.
The residual part of a finite stable module \(N\), for commuting
operators \(T_i\) and scalars \(a_i\), means exactly
\[
 \operatorname{loc}_a(N)=
 \{z\in N:\ \forall i,k\ \exists r\geq0,
                    (T_i-a_i)^rz\in\mathfrak m_A^kN\}.
 \tag{HLM.2}
\]
Thus \(M_K=\operatorname{loc}_a(N_K)\), with the scalar
\(a_i\) mapped from \(O\) to \(A\).

OST19, OST24 and OHA5--OHA11 supply the actual ordinary projector,
finite freeness, arbitrary coefficient base change and residual
idempotents.  ASP13 supplies one specialization map on all levels.
AFS5, AFS19--AFS20 and UPA17 supply the density, part action and
actual automorphic parts when the setting comes from
\(\mathrm{UnitaryHeckeData}\).  LTP supplies the original sections
and their exact local actions.  Their hypotheses are retained,
including the original finite-level and neatness assertions FNL16
and NTL17.  The abstract statements below print any such hypotheses
when they are not yet installed.

The use of local class field theory is its original full
\(\mathrm{LocCFT}\)/Artin interface.  Existence of the auxiliary
universal determinant ring uses the original cohomological
\(\Phi\)-finiteness package, already reconstructed in KX6--KX9
from its stated inputs.  Neither a new local correspondence nor a
new cohomological finiteness theorem is asserted here.

\subsection{Projectivity and the actual directed union}

% BEGIN HLM1
\paragraph{HLM1. An integral trace criterion for projectivity.}
\label{claim:HLM1}
Let a finite group \(\Delta\) act on an \(A\)-projective module
\(N\) by \(\rho\).  If an \(A\)-linear map \(u:N\to N\)
satisfies
\[
             \sum_{g\in\Delta}\rho(g)u\rho(g^{-1})=1,
 \tag{HLM.3}
\]
then \(N\) is projective over \(A[\Delta]\).
Its norm map has image \(N^\Delta\) and kernel generated over
\(A\) by the differences \(\rho(g)z-z\).
\begin{proof}
The induced module \(A[\Delta]\otimes_A N\), with action on
the first factor, is projective: a splitting of an \(A\)-free
module onto \(N\) induces a splitting of free group-algebra
modules.  The equivariant surjection
\(g\otimes z\mapsto\rho(g)z\) has equivariant section
\[
               z\longmapsto
                 \sum_g g\otimes u(\rho(g^{-1})z).
\]
Equivariance follows by replacing \(g\) by \(hg\), and the
composite is (HLM.3).  This proves projectivity.
If \(z\) is invariant, (HLM.3) reads
\(z=\sum_g\rho(g)u(z)\), proving surjectivity of the norm
onto the invariants.  If \(\sum_g\rho(g)z=0\), then
\(\sum_g u(\rho(g^{-1})z)=0\); subtract this from (HLM.3)
applied to \(z\).  The result expresses \(z\) as a sum of
differences.  Conversely every difference is killed by the norm,
by reindexing its sum.  No division by \(|\Delta|\) occurs.
\end{proof}
% END HLM1

% BEGIN HLM2
\paragraph{HLM2. A neat fundamental domain gives the required trace.}
\label{claim:HLM2}
For the original frame \(X\), let \(V\triangleleft U\) be
finite-index subgroups of \(X.U_b\), and assume \(X.\mathrm{Neat}(U)\).
There is a function \(\chi:G_a\times G_p\to\{0,1\}\subset A\)
which is invariant under left \(\Gamma\), right \(V\), and
right \(I(0)\), and satisfies
\[
     \sum_{q\in U/V}\chi(x(q,1))=1\quad\hbox{for every }x.
 \tag{HLM.4}
\]
Multiplication by \(\chi\) preserves the original space
\(\operatorname{Sp}(\tau,\theta,V)\).
\begin{proof}
Choose one representative in each equivalence class for
\(y=X.\mathrm{glob}(\gamma)x(u,b)\),
\(u\in U,b\in I(0)\).  If two presentations from the same
chosen representative have respective \(u,u'\), cancellation
gives a stabilizer relation with away component \(u'u^{-1}\).
Neatness says that component is one.  Thus the away component
\(u\), and in particular its coset modulo \(V\), is unique.
Let \(\chi\) be one precisely when this coset is trivial.
Left global translation and right \(I(0)\) do not change it;
right \(V\) does not change whether it is trivial.  For
\(xq\), its coset is the coset of \(uq\), so precisely one
\(q\in U/V\) contributes to (HLM.4).

The defining equations of \(\operatorname{Sp}\) are global
invariance, invariance under the positive unipotent, the specified
torus character, congruence invariance under a sufficiently deep
opposite unipotent at each precision, and the full matrix-valued
\(V\)-equation.  Every subgroup occurring in the first four
right equations is contained in \(I(0)\).  Hence multiplication
by \(\chi\) commutes with those translations and preserves
their congruences.  In the last equation the scalar \(\chi\)
commutes with the coefficient matrix \(\tau(u)\).  This checks
all equations, including the original coefficient action.
\end{proof}
% END HLM2

% BEGIN HLM3
\paragraph{HLM3. The full ordinary finite-quotient projectivity argument.}
\label{claim:HLM3}
Assume the original finite-level, continuous-weight and neatness
hypotheses at \(U\), and \(V\triangleleft U\) of finite index.
Then \(\operatorname{Ord}(\tau,\theta,V)\) is projective
over \(A[U/V]\) for the actual right-translation action.
The original double-coset trace has image
\(\operatorname{Ord}(\tau,\theta,U)\) and kernel the span
of the translation differences at level \(V\).
\begin{proof}
Finite index transfers the finite class bound to \(V\);
neatness transfers the stabilizer freeness condition.  OST19
therefore gives finite \(A\)-freeness and the ordinary projector
\(e:\operatorname{Sp}(V)\to\operatorname{Ord}(V)\).
Right translations by \(U\) preserve \(\operatorname{Sp}(V)\)
and its ordinary subspace, because \(V\) is normal.  They
commute with \(U_z\), and preserve every value-congruence
filtration.  The factorial-limit characterization of \(e\)
in OST12 consequently gives \(e\rho(u)=\rho(u)e\): compare
the finite powers at each precision and then use separatedness.
No topology on the coefficient isomorphism is involved.

Let \(m_\chi\) be HLM2's multiplication map and put
\(u=e m_\chi\) on \(\operatorname{Ord}(V)\).
Before applying \(e\), (HLM.4) gives, pointwise and with the
matrix factors cancelling,
\(\sum_q\rho(q)m_\chi\rho(q^{-1})=1\).
Commuting \(e\) through this finite sum gives (HLM.3).
HLM1 proves projectivity and the two norm assertions.
Finally the original unnormalized double-coset operator at one
is exactly \(\sum_{q\in U/V}\rho(q)\) on \(V\)-fixed
vectors: its chosen representatives enumerate those right
cosets, and changing a representative multiplies by an element
fixing the vector.  The \(U\)-invariants of \(\operatorname{Ord}(V)\)
are \(\operatorname{Ord}(U)\), since their ambient and
ordinary congruences are identical.  This proves all parts of
the original \texttt{ord\allowbreak Proj} statement at its stated scope.
\end{proof}
% END HLM3

% BEGIN HLM4
\paragraph{HLM4. Residual parts commute with injections and intersections.}
\label{claim:HLM4}
Let \(N\subset N'\) be stable finite modules over a complete
Noetherian local ring, with the same commuting operators.
Then
\[
       \operatorname{loc}_a(N)=N\cap\operatorname{loc}_a(N').
 \tag{HLM.5}
\]
An operator-intertwining linear map carries residual parts into
residual parts.  The canonical residual projector commutes with
every endomorphism commuting with the operators.
\begin{proof}
Functoriality follows directly from (HLM.2).  For the reverse
inclusion in (HLM.5), the Artin--Rees argument gives a fixed
\(c\) such that
\(N\cap\mathfrak m^{k+c}N'\subseteq\mathfrak m^kN\)
for all \(k\).  Here is the needed finite-generation argument:
the Rees ring \(\bigoplus_k\mathfrak m^kt^k\) is Noetherian,
being generated by the finitely many degree-one generators of
\(\mathfrak m\); its finite module
\(\bigoplus_k\mathfrak m^kN't^k\) has finite graded submodule
\(\bigoplus_k(N\cap\mathfrak m^kN')t^k\).
Take \(c\) at least its generator degrees.  Multiplying those
generators proves the displayed inclusion.  If \(z\in N\)
satisfies (HLM.2) in \(N'\), apply its condition at precision
\(k+c\); the relevant operator power is also in \(N\),
so it is in \(\mathfrak m^kN\).
Finally OHA5--OHA6 construct the residual projector as an
idempotent of the finite algebra generated by the operators.
Every commuting endomorphism commutes with this idempotent.
\end{proof}
% END HLM4

% BEGIN HLM5
\paragraph{HLM5. The actual smooth union and all its fixed modules.}
\label{claim:HLM5}
Under the original \(H.\mathrm{Ax}\) and \(\mathrm{IsWt}\),
put \(\mathcal M=\bigcup_KM_K\subset F\).
It is a module with a representation \(\rho\) of
\(G_Q=\prod_qG_q\), characterized by
\[
 \rho(g)f=H.X.\mathrm{ra}(H.\tau_A)(j(g))f,
 \qquad j(g)=\prod_qH.\mathrm{sec}(H.wq(q))(g_q).
 \tag{HLM.6}
\]
Every finite set of its vectors lies in one \(M_K\), and
\(\mathcal M^K=M_K\) as original functions.  In particular
these fixed modules are finite over \(A\).
If all \(H.\mathrm{lev}K\) are neat and the original section
faithfulness condition holds, then for every \(K'\triangleleft K\)
the module \(\mathcal M^{K'}\) is projective over
\(A[K/K']\), with the quotient action induced by \(\rho\).
\begin{proof}
The section images at different places commute by \(H.\mathrm{Ax}\).
Thus the finite product \(j\) is a homomorphism and
\[
       H.\mathrm{lev}K=H.U_Q\,j(K),
              \qquad [H.U_Q,j(G_Q)]=1.
 \tag{HLM.7}
\]
Intersections and conjugates of compact open level tuples are
again such tuples; their indices in a compact level are finite
because the corresponding compact quotient is discrete.
If \(K'\subset K\), then \(N_K\subset N_{K'}\), and
HLM4 gives \(M_K\subset M_{K'}\).  Intersections of levels
therefore prove directedness and the finite-set assertion.

Right translation by \(j(g)\) takes \(N_K\) to
\(N_{gKg^{-1}}\).  This follows from (HLM.7), the original
ambient equations, and commutation with \(U_z\).  It also
commutes with every original Hecke operator: away sections
commute at different places, while the \(\ell\)-adic
operators act on the other factor of the function argument.
HLM4 gives stability of the residual parts, proving (HLM.6).
Its order is the source order
\(\rho(g)\rho(h)=\rho(gh)\), inherited from
\(\mathrm{ra}(u)f(x)=\tau(u)f(x(u,1))\).

Certainly \(M_K\subset\mathcal M^K\).  Conversely take
\(f\in M_{K_1}\cap\mathcal M^K\) and refine both levels
to \(K_2\).  The vector lies in \(N_K\), since it is fixed
by both \(U_Q\) and \(j(K)\); all other defining conditions
are already those of \(N_{K_2}\).
HLM4, applied to \(N_K\subset N_{K_2}\), now gives
\(f\in M_K\).  OHA8 gives the asserted finiteness.

For projectivity, (HLM.7) and faithfulness identify
\(K/K'\) with
\(H.\mathrm{lev}K/H.\mathrm{lev}K'\): surjectivity removes
the common \(U_Q\), and the kernel is exactly \(K'\).
Normality follows from commutation in (HLM.7).
HLM3 makes \(N_{K'}\) projective over this group algebra.
Its residual projector commutes with \(K/K'\) by HLM4.
Consequently \(M_{K'}=\mathcal M^{K'}\) is an equivariant
direct summand and is projective.  This proves the complete
\texttt{pi\allowbreak Basic} output, including the actual fixed-module
identification, rather than just existence of some smooth module.
\end{proof}
% END HLM5

\subsection{Universal Hecke maps at every level}

% BEGIN HLM6
\paragraph{HLM6. Integral determinant interpolation in a compact coefficient ring.}
\label{claim:HLM6}
Let \(T\in\mathcal C_O\), and let a family
\(x:T\to_O\overline{\mathbb Q}_\ell\) separate its elements.
Suppose \(s_j\) is dense in a topological group \(\Gamma\),
and polynomials \(P_j\in T[X]\) specialize, for each \(x\),
to the characteristic polynomials of one continuous
\(n\)-dimensional representation \(r_x(s_j)\).
There is a continuous degree-\(n\) determinant law \(D\) over
\(T\) with \(D\otimes x=\det r_x\) and
\(\chi_D(s_j)=P_j\).
Two determinant laws over any ring are equal if all their
characteristic coefficients on group elements are equal.
\begin{proof}
We recall the integral word identity used in JSP9.  For finitely
many noncommuting letters, the unique ordered Lyndon factorization
gives a formal factorization of \(1-\sum Z_i[g_i]\) into
factors \(1-Z_w[g_w]\) indexed by Lyndon words.  This identity
can be checked degree by degree: at each degree, the leading
remaining word determines the next factor, and unique Lyndon
factorization gives coefficient one for every word of the inverse
series.  Only finitely many factors contribute in any fixed
degree.  Work in each finite-degree truncated scalar algebra:
there the product and every group-algebra element to which the
determinant is applied have finite support.  Apply the
multiplicative homogeneous determinant law there and discard
terms above its degree.  No law on an infinite-support
group-algebra element is presumed.  The coefficients of its value on
\(1-\sum Z_i[g_i]\) are integral polynomials in the
characteristic coefficients of the group words \(g_w\).
The top homogeneous part gives the value on
\(\sum Z_i[g_i]\).  Naturality then permits arbitrary
coefficients in every scalar algebra.  This proves the last
assertion and supplies the required integral formulas.  In
particular, no Newton division in \(T\) is used.

The maximal-ideal quotients of \(T\) are finite: its residue
field is finite, and each successive ideal quotient is finite
dimensional over it.  Completeness and separatedness identify
\(T\) with their inverse limit, so its adic topology is compact
Hausdorff.  ASP2 shows that each given \(x\) is continuous;
indeed it bounds the image of \(\mathfrak m_T^k\) by
\(b_x^k\) with \(b_x<1\).  Thus
\(T\hookrightarrow\prod_x\overline{\mathbb Q}_\ell\)
is a homeomorphism onto a closed subring.
For each \(g\in\Gamma\), the tuple of the \(i\)th
characteristic coefficients of \(r_x(g)\) lies in that subring:
this is true at the dense \(s_j\), and the tuple varies
continuously.  Its unique inverse image \(c_i(g)\in T\)
is therefore continuous.

Use the integral word formulas just proved to define the universal
polynomial associated with every finite sum of group elements,
using these \(c_i(g_w)\).  After every \(x\), it is exactly
the polynomial for \(\det r_x\).  Separation applied to
polynomial coefficients shows that these polynomials are
consistent on overlapping lists, homogeneous of degree \(n\),
multiplicative and unital.  Substitution defines a natural law
on every commutative \(T\)-algebra.  This proves existence of
\(D\), its full scalar-extension identities and continuity.
The specified Frobenius polynomials are recovered coefficientwise
by the same separation.  This proof does not assert injectivity
after tensoring \(T\) into the product of fields.
\end{proof}
% END HLM6

% BEGIN HLM7
\paragraph{HLM7. The full original \texttt{core\allowbreak Det} producer.}
\label{claim:HLM7}
Assume the original \(H.\mathrm{Ax}\) and \(\mathrm{IsWt}\),
the supplied automorphic-set hypothesis
\[
            H.\mathrm{AutomorphicHeckeSet}(\iota,P),
\]
an admitted \(\lambda_0\) with \(\mathrm{DenseWt}(\theta,\lambda_0)\),
and the original local Hecke algebra \(T_K\in\mathcal C_O\).
There are \(D_K,\psi_{K,v,i}\) satisfying every field of
\(H.\mathrm{IsUniv}\) on this very algebra.  This abstract
supplier also permits \(n=0\), although the final unitary
\(\mathrm{LevelAlt}\) producer assumes \(n\geq1\).
\begin{proof}
If \(n=0\), take the constant degree-zero determinant
law \(1\); its characteristic polynomial and the original
degree-zero Hecke polynomial are both \(1\), and the local
character family is empty.  All five fields follow directly.
We may therefore assume \(n\geq1\).
OHA10 makes every positive generator
\(U_{v,j}=\mathrm{gen}(\mathrm{inr}(v,j))\) a unit of
\(T_K\).  Put \(U_{v,-1}=1\) and
\(u_{v,i}=U_{v,i}U_{v,i-1}^{-1}\).
On the unit group of \(L_{H.S_l(v)}\), prescribe
\(\xi_{v,i}(y)=\theta(e_{v,i}(y))^{-1}\), now mapped
to \(T_K\).  Extend to the multiplicative group by
\[
 \xi_{v,i}(y\varpi^a)=\xi_{v,i}(y)u_{v,i}^a.
 \tag{HLM.8}
\]
The valuation decomposition makes this a character.  The original
\(\mathrm{ContChar}\) condition says its reduction modulo each
\(\mathfrak m_{T_K}^k\) is continuous.  That quotient is finite,
so the full local reciprocity interface extends it uniquely to
a continuous character of the local abelian Galois group.
Taking the inverse limit and multiplying by
\(\epsilon^{-i}\) defines \(\psi_{K,v,i}\) on the actual
decomposition group.  The unit formula is immediate, and the
uniformizer formula is the telescoping identity
\[
 \prod_{i\leq j}\psi_{K,v,i}(g)\epsilon(g)^i=U_{v,j}
      \quad\text{if }\operatorname{Art}(g)=\varpi.
 \tag{HLM.9}
\]
This proves both original character formulas and their adic
continuity, on every permitted Weil lift.

AFS20 says that the good \(T_K\)-points at \(\lambda_0\)
separate \(T_K\).  For each specified such point, ASP13 and
AFS19(3) give a nonzero eigenvector in an actual part.  Its
representation is continuous, and its spherical eigenvalues are
the point's values of the exact original generators.  The part's
Frobenius equation therefore says that its characteristic
polynomial at \(H.\mathrm{prj}(\sigma^{-1})\) is
\[
 \operatorname{heckePol}\bigl(n,\operatorname{Norm}w,
                               x(\operatorname{genW}(w,-))\bigr).
\]
The degree-zero term and indices \(j+1\) are those of the
original \(\operatorname{genW}\); no spherical normalization
is changed.  \(H.\mathrm{Ax.dense}\) and HLM6 produce
\(D_K\) with this polynomial and adic continuity.  Along
with (HLM.8)--(HLM.9), these are exactly the five
\(\mathrm{IsUniv}\) fields.  The construction precedes any
use of an \(\mathrm{Out}\) record that assumes such a pair.
\end{proof}
% END HLM7

% BEGIN HLM8
\paragraph{HLM8. Uniqueness and functoriality of the original universal action.}
\label{claim:HLM8}
Two \(\mathrm{IsUniv}\) pairs on the same original Hecke
algebra are equal.  If an injective \(A\)-linear map
\(f:M_K\to M_{K'}\) intertwines every original Hecke
generator, and a ring \(R\) with pair \((D_0,\psi_0)\)
has the source's unique-map property into \(\mathcal C_O\),
then any given \(\mathrm{IsMap}\) homomorphisms from \(R\)
to the two Hecke algebras induce intertwining actions on \(f\).
\begin{proof}
The two determinant characteristic coefficients agree on the
specified dense Frobenius set.  Modulo each maximal-ideal power
they are continuous functions to a finite discrete ring, so they
agree everywhere.  Separatedness and HLM6 give equality of full
determinant laws.  The unit formulas determine each local
character on units.  Equality of the prefix products in the
uniformizer formulas determines the individual values by
cancellation in the unit group.  Units and one uniformizer
generate the local multiplicative group, so reciprocity gives
equality on the Weil subgroup.  Its density, continuity modulo
every ideal power, and separatedness give equality on the
decomposition group.

For the action assertion, first construct
\(\Phi:T_{K'}\to_A T_K\).  Every polynomial in the target
operators preserves \(f(M_K)\); its restriction there is the
same polynomial in the source operators.  If two polynomials
give the same target operator, injectivity of \(f\) makes
their source operators equal.  This defines \(\Phi\) on the
algebra generated by the operators, preserving all scalars and
generators, and gives
\(f\Phi(t)=tf\).  It is an \(O\)-algebra map between
\(\mathcal C_O\) objects, hence local and adically continuous:
the common residue maps give
\(\Phi(\mathfrak m^k)\subset\mathfrak m^k\).
Pushing the target \(\mathrm{IsUniv}\) pair through \(\Phi\)
preserves all five fields, since it preserves the relevant
scalars, units and generators.  Uniqueness makes it the source
pair.  The unique-map property of \(R\) then gives
\(f_K=\Phi\circ f_{K'}\), which is precisely the asserted
intertwining.  If the source module is zero the assertion is
automatic; an injection from a nonzero source excludes a zero
target.  This proves the literal \texttt{TCore\allowbreak Uniq} and
\texttt{Level\allowbreak UR2Compat} consumers without an extra density
hypothesis on \(A\).
\end{proof}
% END HLM8

% BEGIN HLM9
\paragraph{HLM9. Original universal Hecke maps and ideals exist at every level.}
\label{claim:HLM9}
For every original supplied \(\mathrm{UnitaryHeckeData}\),
universal ordinary determinant ring \(B\), weight object
\(\Lambda\), and level \(K\), there is an ideal satisfying
the full original \(\mathrm{IsJ}\).  If \(M_K\ne0\), its
map is a surjection \(f_K:B\to T_K\), is an
\(\mathrm{IsMap}\), and satisfies
\(f_K\Lambda.\mathrm{toPair}(B.u)=A\to T_K\).
\begin{proof}
If \(M_K=0\), choose \(J=B\), the literal zero branch.
Otherwise OHA11 gives the nonzero \(\mathcal C_O\) algebra
\(T_K\) with its exact residual eigenvalues.  FNL16,
UPA17 and AFS5 install HLM7 at the admitted zero weight.
It gives \(D_K,\psi_K\); HEP11 now applies to this supplied
pair and gives its actual separating good points and their
continuous representations and flags.

We check that this is an original ordinary determinant pair.
Modulo \(\mathfrak m_{T_K}\), the Frobenius formula and
\(\mathrm{IsRes.frob}\) agree.  Their coefficient functions
are locally constant on the actual group.  Density and HLM6
give \(\overline D_K=r.\overline D\) as full laws.
On inertia, compare the \(\mathrm{IsTheta}\) equation with
the unit equation (HLM.8); cancellation gives
\(\psi_{K,v,i}|_{I_v}=\kappa_{v,i}\) after the given
scalar map.  Reduction of (HLM.9) and the literal residual
uniformizer formula identify the residual characters on a
uniformizer.  Their agreement on inertia, then on the dense
Weil subgroup, gives the full prescribed reduction.  For the
torsion condition, the original \(\Lambda.\mathrm{inert}\)
or \(\mathrm{IsKap}\) field supplies a continuous extension
\(\psi^0_{v,i}\) of \(\kappa_{v,i}\) having the prescribed
residual and torsion values; map it to \(T_K\).  The quotient
\(\psi_{K,v,i}/\psi^0_{v,i}\) is trivial on inertia.
The actual local decomposition group modulo inertia is the
unramified group \(\widehat{\mathbb Z}\), which is
torsion-free.  An element of the source \(\mu_G\) has
\(\ell\)-power-torsion image in the local abelianization,
so its image in this unramified quotient is zero.  The
quotient character therefore kills the full preimage
\(\mu_G\), exactly as required by the original
\(\mathrm{IsPsi.mu}\) field.  Thus every \(\mathrm{IsPsi}\) field
holds, and the weight map is the original one.

The remaining pair relations can be checked at the separating
good points.  The part's form makes the determinant invariant
under the prescribed star.  The ordered flag kills the ordered
ordinary product, in its source order.  For each residual type
place choose the original Teichmuller lifts of the residual
characters.  The part's type diagonalization makes every
element of the source \(I_p(u)\) act trivially, because its
residual characters are all one and the Teichmuller lift of one
is one.  Thus the required trace relations hold at every point.
The universal coefficient argument of JRP6--JRP12, or directly
the polynomial-coefficient argument of JSP2, descends these
relations to \(T_K\).  This constructs all fields of the
original pair, not just its values at points.

Universality of \(B\) gives \(f_K\) carrying its pair to
\((D_K,\psi_K)\).  Universality of \(\Lambda\), using the
inertia identity already proved, gives weight compatibility.
Every \(U_{v,j}\) is the prefix product (HLM.9), hence lies
in the image.  The \(j\)th away generator is the
\((n-j-1)\)st characteristic coefficient multiplied by
\((-1)^{j+1}q^{-(j+1)j/2}\).  This scalar is an original
\(O\)-unit because that place is away from \(\ell\).
Consequently every Hecke generator and every \(A\)-scalar
is in the image.  They generate \(T_K\), so \(f_K\) is
surjective.  Taking \(J=\ker f_K\) proves the complete
nonzero \(\mathrm{IsJ}\) branch.  This argument is valid
for all original ranks, not just the rank-four application
of JSP3--JSP5.
\end{proof}
% END HLM9

% BEGIN HLM10
\paragraph{HLM10. One universal-ring action on the entire union.}
\label{claim:HLM10}
Assume the original abstract \texttt{umod\allowbreak Abs} hypotheses:
\((R,D_0,\psi_0)\) has the unique-map property of HLM8,
\(\lambda:A\to_O R\) is given, and for every \(K\)
there exists an original \(\mathrm{IsJ}\) for these data.
Then \(\mathcal M\) of HLM5 has one ring action
\(\alpha:R\to\operatorname{End}_A(\mathcal M)\), with
\[
 \alpha(\lambda(b))z=bz,\qquad
 \alpha(t)\rho(g)z=\rho(g)\alpha(t)z.
 \tag{HLM.10}
\]
Its restriction to each nonzero \(M_K\) is induced by an
original \(\mathrm{IsMap}\) homomorphism \(R\to T_K\).
\begin{proof}
Choose the action \(\mu_K\) furnished by the stated
\(\mathrm{IsJ}\), using the unique zero action if \(M_K=0\).
It has the first identity of (HLM.10).  For an inclusion
\(M_K\subset M_{K'}\), the operators are the same actual
operators on functions.  HLM8 says that the actions agree.
For \(z\in\mathcal M\), define \(\alpha(t)z=\mu_K(t)z\)
at any level containing it.  A common refinement proves
independence, additivity, \(A\)-linearity, and both ring laws;
all involve finitely many vectors, so one such refinement
suffices for each equality.

For a group element \(g\), translation is an injective
Hecke-intertwining map \(M_K\to M_{gKg^{-1}}\), with
inverse translation by \(g^{-1}\).  HLM8 again proves
the second identity of (HLM.10).  The restriction formula is
the definition, including its equality in the original ambient
function space.  Thus no level-dependent action is being
silently identified with another one.
\end{proof}
% END HLM10

\subsection{Classical points separate all module vectors}

% BEGIN HLM11
\paragraph{HLM11. One scalar-character projection on a directed union.}
\label{claim:HLM11}
Let an additive map \(s:M\to V\) take values in a complex vector
space, let \(M=\bigcup_KM_K\) be directed, and suppose the
actions \(\alpha(t)\) and \(\rho(g)\) commute.  Suppose
\(s\rho(g)=\pi(g)s\).  Put
\(E_K=\operatorname{span}_{\mathbb C}s(M_K)\).
Assume every \(E_K\) under consideration has an action of
\(R\) agreeing with \(s\alpha\), and is spanned by
subspaces contained in specified parts \(Y_i\), on each of
which this action is scalar through a character
\(\chi_x:R\to\mathbb C\) related to that part.
It suffices to have these assertions at all levels containing
one specified \(z\).  If \(s(z)\ne0\), there are a character
\(\chi_{x_0}\), a \(\pi\)-stable subspace \(E_0\), and
a complex linear projection on \(E=\operatorname{span}s(M)\)
onto \(E_0\), such that
\[
 p(s(z))\ne0,\qquad p(s(\alpha(t)m))=\chi_{x_0}(t)p(s(m)),
 \qquad p\pi(g)=\pi(g)p\quad\hbox{on }E.
 \tag{HLM.11}
\]
Moreover \(E_0\) is contained in the sum of those original
parts related to a character with the same values as
\(\chi_{x_0}\).
\begin{proof}
Any finite linear relation among the \(s(m)\) lies at one
level containing \(z\).  Applying the given action on that
\(E_K\) shows that the same relation holds among
\(s(\alpha(t)m)\).  Thus these values define a unique
linear operator \(T_t\) on \(E\).  The ring laws and the
commutation with \(\pi\) follow on the spanning set \(s(M)\).
Each vector belongs to an \(E_K\), hence is a finite sum
of simultaneous scalar vectors of the stated kinds.

Eigenspaces for distinct characters are linearly independent.
Indeed, in a finite purported relation and for a chosen
character \(\chi_0\), choose for each other character
\(\chi_j\) some \(t_j\) with
\(\chi_0(t_j)\ne\chi_j(t_j)\).  The commuting operator
\(\prod_j(T_{t_j}-\chi_j(t_j))\) kills all other summands
and is a nonzero scalar on the chosen one.  Thus the relation's
chosen summand is zero.  It follows that \(E\) is the
algebraic direct sum of its character spaces.  At least one
component of \(s(z)\) is nonzero.  Take \(p\) to be that
component projection.  Since \(\pi(g)\) commutes with all
\(T_t\), it preserves every character space and commutes
with \(p\).  Grouping the original spanning subspaces by
character proves the last assertion.  If a map on the entire
ambient \(V\) is desired, extend \(p\) by zero on any
linear complement of \(E\); every asserted identity concerns
\(E\), so this extension introduces no equivariance claim
outside it.
\end{proof}
% END HLM11

% BEGIN HLM12
\paragraph{HLM12. Finite local isotypic refinement preserves all bridges.}
\label{claim:HLM12}
Let \(Q\) be finite.  If a representation of \(\prod_qG_q\)
is, for each factor, a sum of copies of specified irreducible
representations, then every nonzero vector survives a group
equivariant linear map to a representation isotypic for one
specified irreducible at each factor.  Suppose the specified
irreducibles arise from original parts having a common full
determinant law.  One may use the Galois representation of any
one such part in all the resulting local \(\mathrm{WOT}\)
assertions, for every prolongation and every original
\(\mathrm{IsRecT}\) map.
\begin{proof}
A sum of simple modules is semisimple: select a maximal direct
sum of simple submodules; a simple module not contained in it
would have zero intersection with it and enlarge the sum.
Every submodule of a semisimple module is semisimple, as follows
by the same maximal-sum argument, or by selecting complements
to its intersections with finite sums of simple modules.
Group the simple summands by isomorphism class.  A homomorphism
between nonisomorphic simple modules is zero, so their
isotypic sums form a canonical direct decomposition.  Every
endomorphism commuting with this factor preserves that
decomposition.  Consequently its component projections commute
with all the other factors.  Project onto a component on which
the chosen vector is nonzero, and repeat for the finitely many
factors.  Previous isotypic properties persist under taking
submodules and images.  This proves the first assertion,
including \(Q=\varnothing\), when no projection is needed.

For the second assertion, the finite-dimensional Galois
representations of the parts are semisimple.  Equality of their
determinants gives equality of their characters.  Here the
coefficient field has characteristic zero.  To recall the
character argument, take the finite-dimensional algebra generated
by the simultaneous representation on their direct sum.
Its radical acts trivially on both semisimple modules, and its
semisimple quotient is a product of matrix algebras over the
algebraically closed field.  The traces of the central primitive
idempotents give the multiplicities times the simple dimensions;
they determine the multiplicities.  Equal characters on the
group give equal traces on its linear span, hence on these
idempotents.  Thus the two representations are isomorphic.
Conjugation induces the same isomorphism on the period modules
defining \(\mathrm{WOT}\).  It transports the part's local
comparison for each given prolongation and each given
\(\mathrm{IsRecT}\) map.  Nothing in this transport chooses
one preferred local correspondence or weakens the universal
quantifier over them.
\end{proof}
% END HLM12

% BEGIN HLM13
\paragraph{HLM13. The full original \texttt{pi\allowbreak Pts\_\allowbreak n\_\allowbreak alt}.}
\label{claim:HLM13}
Under the original hypotheses of HLM5, an admitted
\(\lambda_0\) with \(\mathrm{DenseWt}(\theta,\lambda_0)\),
the supplied \(\mathrm{AutomorphicHeckeSet}(\iota,P)\),
and the levelwise \(\mathrm{IsUniv}\) restriction property
of HLM10, every nonzero \(z\in\mathcal M\) admits
\[
 (x,r,Y,\pi,\mathrm{sp},(\sigma_q)_q),
 \quad x:R\longrightarrow\overline{\mathbb Q}_\ell,
 \quad r:H.G_q\longrightarrow\mathrm{GL}_n(\overline{\mathbb Q}_\ell),
 \tag{HLM.12}
\]
where \(Y\) is a complex vector space, \(\pi\) is a
representation of \(\prod_qG_q\), \(\mathrm{sp}\) is
additive, and each \(\sigma_q\) is the original smooth
irreducible local representation, with all the following properties:
\[
 \begin{gathered}
 \mathrm{sp}(z)\ne0,\qquad r\text{ is continuous},\qquad
 D_0\mathbin{\mathrm{map}}x=\det r,\\
 \mathrm{sp}(\alpha(t)m)=\iota(x(t))\mathrm{sp}(m),\qquad
 \mathrm{sp}(\rho(g)m)=\pi(g)\mathrm{sp}(m),\\
 \operatorname{IsoSum}(\pi|_{G_q},Y,\sigma_q)\quad(q\in Q).
 \end{gathered}
 \tag{HLM.13}
\]
For every \(q\), every valuation subring \(B_q\) over
\(H.wq(q)\), and every \(rc\) satisfying its original
\(\mathrm{IsRecT}\), one has
\(\mathrm{WOT}(\iota,r\circ H.\mathrm{prj},
rc(\operatorname{algNorm}\sigma_q))\).
\begin{proof}
Choose a level containing \(z\), and an \(A\)-basis of
its finite free ordinary module \(N_K\).  Some coefficient
of \(z\) is nonzero.  Density supplies a good point
\(x_0:A\to_O\overline{\mathbb Q}_\ell\) at which that
coefficient remains nonzero.  ASP13 gives one specialization
\(s\) on all ambient functions, whose images of this basis
are linearly independent.  Expansion in the basis proves
\(s(z)\ne0\).  On the whole union \(s\) is additive:
any two vectors lie at one level, where ASP13 asserts
additivity.  The same argument gives its scalar rule.

For each level put \(E_K=\operatorname{span}_{\mathbb C}s(M_K)\).
AFS19 gives the actual Hecke action on \(E_K\), its spanning
by intersections with original parts, and a scalar character
on every nonzero such intersection.  At any level containing
\(z\), its Hecke algebra is nonzero and is in \(\mathcal C_O\)
by OHA11.  Compose its specified map \(R\to T_K\) with
each of these scalar characters.  HEP7, applied to a nonzero
vector in the intersection, identifies the resulting full
determinant with that of the part.  Thus these are exactly
the related characters required in HLM11.

The finite product of the actual sections defines the complex
right-translation representation \(\pi_0\) on the ambient
classical functions.  ASP13's action field, applied after
putting a vector and its translate at a common level, gives
\(s\rho(g)=\pi_0(g)s\).  For a single factor the relevant
operators commute: away operators are finite sums of translations
at disjoint sections, and the ordinary operators are finite
sums on the other factor of \(G_a\times G_p\).
Their order is the original right-translation order.  This
also proves the scalar action commutation if only the
levelwise restriction property, rather than HLM10's stronger
global conclusion, was assumed.

Apply HLM11.  Its surviving character is
\(\iota\circ x\) for an actual map \(x:R\to\overline{\mathbb Q}_\ell\),
and its image lies in the sum of parts whose full determinant
is \(D_0\mathbin{\mathrm{map}}x\).  Their local
\(\operatorname{IsoSum}\) fields imply the same spanning
assertion on this subrepresentation, since submodules of a
semisimple module are semisimple.  HLM12 now refines the
finitely many local factors while retaining a nonzero image
of \(s(z)\).  Compose the two projections with \(s\).
Use the continuous Galois representation of the part originally
selected in HLM11.  HLM12 transports all local comparisons
from the parts selected during refinement.  Scalar and group
equivariance survive both projections, proving every field
in (HLM.13) and every stated local comparison.  In particular,
the conclusion is for each module vector, not just for
eigenvectors already chosen at a finite level.
\end{proof}
% END HLM13

% BEGIN HLM14
\paragraph{HLM14. The complete abstract \texttt{umod\allowbreak Abs\_\allowbreak alt} interface.}
\label{claim:HLM14}
Assume precisely the supplied original \(H.\mathrm{Ax}\),
\(\mathrm{IsWt}\), \(\mathrm{AutomorphicHeckeSet}(\iota,P)\),
admitted dense weight, neatness of all levels, and section
faithfulness of HLM5, together with the original literature
hypothesis and the unique-map and all-level \(\mathrm{IsJ}\)
hypotheses of HLM10.  Then the original union \(\mathcal M\)
has the representation \(\rho\) and action \(\alpha\)
of \texttt{umod\allowbreak Abs\_\allowbreak alt}: they commute, the weight algebra
acts by its original scalars, the \(K\)-fixed module is
exactly \(M_K\) and is finite, every normal compact-level
quotient has the specified projective group-algebra action,
the restrictions are actual \(\mathrm{IsMap}\) actions,
and all nonzero vectors have exactly the points of HLM13.
\begin{proof}
HLM5 constructs the representation, fixed-space identity and
all the quotient actions, with the displayed compatibility
on each quotient coset.  HLM10 constructs the single ring
action and proves its restriction identities.  HLM13 uses
those same maps and the same union, so its scalar and group
equivariance are the required ones.  These conclusions do
not require changing either action after choosing a point.
This lists all output fields of the abstract theorem.  In
particular the abstract statement retains its supplied
automorphic and all-level-ideal hypotheses; HLM24 below
constructs those hypotheses for the actual auxiliary levels.
\end{proof}
% END HLM14

\subsection{Changing the allowed ramification without changing the data}

% BEGIN HLM15
\paragraph{HLM15. The original quotient, residual object and pair pullback.}
\label{claim:HLM15}
Let \(V_0,V\) be original levels and
\(h:\operatorname{globOf}(r,V)\to\operatorname{globOf}(r,V_0)\)
an original global-frame homomorphism.  Assume its stated
compatibility with the absolute Galois quotient maps.  Its
restriction
\(p_0:V.\mathrm{frame}.G_0\to V_0.\mathrm{frame}.G_0\)
is continuous and surjective, and
\[
 p_0(V.\mathrm{pr}(g))=V_0.\mathrm{pr}(g).
 \tag{HLM.14}
\]
The original \(r^{\mathrm{up}}\) has global residual
characters \(\bar\chi_i\circ p_0\), unchanged local
\(\bar\psi_{v,i}\) and torsion lifts, and unchanged
weight object \((A,\kappa)\).  A pair on \(r\) pulls
back to a pair on \(r^{\mathrm{up}}\) by
\((D,\psi)\mapsto(D\circ p_0,\psi)\).
It also gives the original \(r_N\), compatibility record,
and pair on \(r_N\) used by \(\mathrm{LevMod}\).
\begin{proof}
Continuity is restriction of \(h.\mathrm{cont}\).  Lift a
point in the target kernel of the quadratic character through
the surjection \(h.p\).  Compatibility of quadratic
characters says that its lift is in the source kernel; this
proves surjectivity of \(p_0\).  The assumed compatibility
of quotient maps gives (HLM.14).

The residual characters remain pairwise distinct because
\(p_0\) is surjective, and their kernels are open by
continuity.  The star involution on group algebras commutes
with the map induced by \(p_0\): on a group basis element
this is exactly preservation of conjugation and of the
cyclotomic character by \(h\), and linearity proves it
on all elements.  Thus residual polarization pulls back.
The local groups, their inertia subgroups, residual characters
and torsion lifts have not changed.  Consequently their
\(\mathrm{IsPsi}\) and inertial-extension conditions are
identical, and the same universal \((A,\kappa)\) represents
the weight functor.

Continuity, reduction, polarization and the ordered ordinary
relations of a pair pull back along these same group-algebra
maps.  For a type place, an element of the new \(I_p(u)\)
comes from the same absolute inertia element and lies in
every \(\ker(\bar\chi_i\circ p_0)\).  Its image under
\(p_0\) therefore belongs to the old \(I_p(u)\).
The old trace relation gives the new one.  This proves the
pair pullback, including all scalar algebras in its polynomial
law fields.  The residual object \(r_N\) in the patching
frame has the same global characters and local data.  Its
type inertia subgroup is contained in that of
\(r^{\mathrm{up}}\), directly from its definition by
the restricted inertia domain.  Restricting that one relation
therefore makes the same \(D,\psi\) a pair on \(r_N\).
These are precisely \texttt{r\allowbreak Up}, \texttt{lam\allowbreak Up},
\texttt{pair\allowbreak Up}, \texttt{r\allowbreak N\_\allowbreak compat} and \texttt{pair\allowbreak N};
no new residual representation or new weight ring is chosen.
\end{proof}
% END HLM15

% BEGIN HLM16
\paragraph{HLM16. The identical type works for every lifted Teichmuller family.}
\label{claim:HLM16}
Under HLM15, the given \(\mathrm{TypeForAlt}(r,\iota,\tau_0)\)
implies \(\mathrm{TypeForAlt}(r^{\mathrm{up}},\iota,\tau_0)\).
If \(V_0.S'\subset V.S'\), let \(j=\mathrm{opUp}\)
keep all ordinary indices and include every surviving away
index.  Then the same \(\theta\), the residual eigenvalues
\(a'=a\circ j\), and the same \(\tau_0\) satisfy the
original new-level \(\mathrm{IsTheta}\),
\(\mathrm{IsWt}\), \(\mathrm{IsRes}\), and
\(\mathrm{UnitaryHeckeData}\) conditions.
\begin{proof}
Write \(k\) for the finite residue field of \(O\) and
\(N=|k|-1\).  Two units \(u,v\) with equal reduction
and \(u^N=v^N=1\) are equal.  Indeed
\[
 (u-v)\sum_{a=0}^{N-1}u^av^{N-1-a}=u^N-v^N=0,
\]
and the sum reduces to \(N\bar v^{N-1}\ne0\), so is
a unit in the local ring.  Apply this with \(v=1\) to
an arbitrary original Teichmuller family \(t'_i\) lifting
\(\bar\chi_i\circ p_0\).  Each \(t'_i\) kills
\(\ker p_0\), hence factors uniquely as \(t_i\circ p_0\).
Surjectivity verifies that \(t_i\) is a Teichmuller lift
of the old character.  Equation (HLM.14) identifies the
restrictions of these two families to the original absolute
inertia at every type place.  Apply the old
\(\mathrm{TypeForAlt}\) to this descended family.  Its
assertion about every original \(\mathrm{IsRecT}\) map
has exactly the same local characters and representation
\(\tau_0\), so is the required new assertion.  The argument
began with an arbitrary new family and retains every such
family and every local bridge.

The torus, inertia and reciprocity maps at \(\ell\) are
unchanged.  The \(\mathrm{IsTheta}\) equation and
\(\mathrm{ContChar}\) condition therefore remain literal
identities.  NTL17 supplies neatness at each new level.
In a stabilizer equation neatness forces both acting elements
to be one, and then both coefficient characters are one;
this is exactly the freeness field of \(\mathrm{IsWt}\).
For \(\mathrm{IsRes.frob}\), the new residual determinant
is the old one composed with \(p_0\), and (HLM.14)
identifies its Frobenius argument.  The generator is indexed
by \(j\), so its prescribed coefficient is \(a\circ j\).
The residual uniformizer formula is unchanged because the
local characters, uniformizer and ordinary indices are
unchanged.  All fields of the new unitary data follow from
these checks.  This proves the literal \texttt{Type\allowbreak For\allowbreak Up}
and the data part of \texttt{step1\_\allowbreak alt} without requiring
a new integral type.
\end{proof}
% END HLM16

% BEGIN HLM17
\paragraph{HLM17. Exact level equality and faithful auxiliary sections.}
\label{claim:HLM17}
For the original unitary groups, membership in
\(\operatorname{lev}K\) is equivalent to the original
maximal-compact conditions away from \(V.S'\) and the
conditions \(\operatorname{cmp}_{w_q}(u)\in K_q\).
If \(V_0.S'=s.S\) and every new \(K_q\) is maximal
compact, then this subgroup equals the original
\(\operatorname{lev}K_0\), for every original \(K_0\).
The section product also satisfies the full faithfulness
hypothesis in HLM14, for every compact auxiliary level.
\begin{proof}
The forward membership assertion holds for every generator
of \(U_Q\vee\bigvee_q\mathrm{sec}_{w_q}(K_q)\): the
sections have identity component away from their underlying
place and the indicated component at their own place.
For the converse, multiply \(u\) on the left successively
by
\(\mathrm{sec}_{w_q}(\operatorname{cmp}_{w_q}u)^{-1}\).
Each factor is in the proposed level.  It makes the selected
component one and leaves every other selected component,
and all components outside \(V.S'\), unchanged.  After the
finite list is exhausted, the remainder belongs to \(U_Q\)
by its defining conditions.  Reversing the multiplications
proves membership.  This argument uses the actual sections,
not an abstract isomorphic product.

The old auxiliary index set is empty: its places would have
underlying place both in \(V_0.S'=s.S\) and outside
\(s.S\).  At each new selected place the maximal-compact
condition also imposes it at its conjugate.  Indeed the unitary
identity identifies that component with conjugated inverse
transpose, which preserves the integral invertible matrices;
this is the opposite-place check of LTP2.  Places above
\(\ell\) have trivial away component.  Thus the displayed
membership criterion is exactly the old one outside \(s.S\),
in both directions.  This proves literal subgroup equality.

Finally let \(j:\prod_qG_q\to U_b\) be any homomorphism
having the stated section on each individual factor.  Since
the product is finite, it is the product of those sections.
Their support and right-inverse identities give
\(\operatorname{cmp}_{w_q}(j(g))=g_q\).  Applying the
forward membership criterion to \(j(g)\in\operatorname{lev}K\)
therefore yields \(g_q\in K_q\) for every \(q\).
This proves the exact \texttt{level\allowbreak UR2Faith} quantifier over
all such \(j\).  All assertions also hold for an empty
index set, using the empty product and no deletion steps.
\end{proof}
% END HLM17

\subsection{The exact old-level identification}

% BEGIN HLM18
\paragraph{HLM18. Every ordinary Hecke character occurs in an original part.}
\label{claim:HLM18}
Let \(N=N_K\) be the full ordinary module and
\(T=\operatorname{heckeAlg}(N,H.\mathrm{op})\), before
selecting a residual part.  Let \(x:A\to_O\overline{\mathbb Q}_\ell\)
be an arithmetic specialization of admitted weight, and
\(z:T\to\overline{\mathbb Q}_\ell\) a ring map
restricting to \(x\).  Under the original supplied
automorphic decomposition there is an original part \(Y\)
with
\[
  Y.a(w,j+1)=\iota\bigl(z(T_{w,j})\bigr)
       \quad\hbox{for every away index }(w,j).
 \tag{HLM.15}
\]
\begin{proof}
The algebra \(T\) is finite over \(A\), and \(N\) is
a finite faithful \(T\)-module.  Its support is all of
\(\operatorname{Spec}T\): a finite module localized to
zero is killed by some element outside that prime, which
would put the prime outside the zero set of its annihilator.
Thus the point defined by \(z\) has nonzero fiber by
Nakayama.  Equivalently its maximal-ideal block in
\(N\otimes_{A,x}\overline{\mathbb Q}_\ell\) is nonzero.
In a nonzero module over a local finite-dimensional algebra,
take the last nonzero power of the maximal ideal acting on
the module.  It is killed by that maximal ideal.  A nonzero
vector there is an eigenvector for every \(t\in T\)
with eigenvalue \(z(t)\).  This is the faithful-fiber
argument of AFS9, and does not require flatness of the
coefficient extension.

ASP13 identifies this base-changed module with the complex
span of its specialized basis, so the vector remains nonzero
after applying \(\iota\).  Its operator equations are
exactly the specialized original equations.  CFD12's
classical ordinary decomposition expresses it as a finite
sum of vectors in original parts.  All away operators act
on each part by the stated \(Y.a(w,j+1)\).  Group those
finitely many summands by their away scalar characters.
The character-independence argument in HLM11 shows that
at least one occurring character equals the eigenvector's
character: otherwise a product of separating operators
would annihilate every summand but act as a nonzero scalar
on the eigenvector.  This proves (HLM.15) simultaneously
at every away index.  It is the full
\texttt{level\allowbreak UR2Ord\allowbreak Pt\_\allowbreak alt} assertion; it does not first
assume the character is a point of the prescribed residual
Hecke block.
\end{proof}
% END HLM18

% BEGIN HLM19
\paragraph{HLM19. A dense retained Frobenius set determines the residual part.}
\label{claim:HLM19}
In HLM18 suppose a collection of retained Hecke indices
contains every ordinary index, and its away places have
dense inverse-Frobenius image in the original global group.
Suppose \(a\) has the original residual Frobenius
polynomials.  Then selecting the residual part using only
the retained indices gives exactly the full residual part.
This holds over the original \(A\) with a dense admitted
weight and does not assume a flat specialization.
\begin{proof}
Let \(T_0\subset T\) be the finite algebra generated by
the retained operators.  OHA5--OHA6 give idempotents
\(e_0\in T_0\) and \(e\in T\) selecting the indicated
residual characters.  Full residual conditions imply the
retained conditions, so \(eN\subset e_0N\) and
\(e_0e=e\), by faithfulness.  If equality fails, the
idempotent \(d=e_0(1-e)\) has \(dN\ne0\).  This is
a finite projective, hence finite free, \(A\)-module of
positive rank.  Choose a good point \(x\) at the admitted
weight, which exists by density applied to \(1\).
Then \(dN\otimes_{A,x}\overline{\mathbb Q}_\ell\ne0\).
An eigencharacter of its finite commuting operator algebra
gives \(z:T\to\overline{\mathbb Q}_\ell\) restricting
to \(x\), with \(z(e_0)=1\) and \(z(e)=0\).

Every element of \(T\) is integral over \(A\).  ASP2
gives \(|x(A)|\leq1\), so a monic integral equation
shows \(|z(T)|\leq1\): a root of norm greater than one
would make its top power strictly dominate all lower terms.
The finite set of generators of \(\mathfrak m_A\) has
images bounded by some \(b<1\), hence
\(|z(\mathfrak m_AT)|\leq b\).  In the \(e_0\)-factor,
each retained \(T_i-a_i\) has nilpotent reduction modulo
\(\mathfrak m_A\).  Some power times \(e_0\) thus
lies in \(\mathfrak m_AT_0\), and
\[
                 |z(T_i)-a_i|<1
                 \quad(i\text{ retained}).
 \tag{HLM.16}
\]

HLM18 supplies a part with all the away eigenvalues of
\(z\).  Its continuous Galois representation has the
original Frobenius polynomials.  Choose lifts in \(O\)
of the finitely many values of the residual characteristic
coefficient functions; the resulting functions on the
global group are locally constant.  At every retained
inverse Frobenius, (HLM.16) says that each characteristic
coefficient of the part differs from its prescribed lift
by an element of norm less than one.  The set on which
each such inequality holds is closed as well as open:
the coefficient functions are continuous and the open
unit ball of \(\overline{\mathbb Q}_\ell\) is also
closed.  The dense retained set therefore gives the same
inequalities at every group element.  Apply them at every
old away Frobenius.  In the original Hecke polynomial the
coefficient indexed by \(j+1\) is
\((-1)^{j+1}q^{(j+1)j/2}z(T_{w,j})\).
Here \(|q|_\ell=1\); cancellation gives (HLM.16) at
every old away index as well.  The ordinary indices were
retained, so the inequality now holds at every generator.

Reduce the bounded map \(z\) in the residue field of
\(\overline{\mathbb Q}_\ell\).  Its values on all
generators are exactly the prescribed residual scalars,
and its values on \(A\) are the prescribed residue map.
Since these generate \(T\), its kernel is precisely the
maximal ideal selected by \(e\).  The corresponding
idempotent must map to one, contrary to \(z(e)=0\).
Thus \(d=0\), and the two modules coincide.  The use of
a positive-rank direct summand before specialization is
what ensures that no unsupported faithfulness assertion
after base change enters this argument.
\end{proof}
% END HLM19

% BEGIN HLM20
\paragraph{HLM20. Literal equality of the old module and maximal-compact invariants.}
\label{claim:HLM20}
For the data of HLM15--HLM17, let \(K^{\max}\) denote
the new product of maximal compact groups.  Then every
surviving operator is the same original endomorphism,
\[
 H'.\mathrm{op}_A(i)=H.\mathrm{op}_A(j(i)),\qquad
 M'_{K^{\max}}=M_{K_0}
           \quad\hbox{inside the original function space}.
 \tag{HLM.17}
\]
Here \(H'\) denotes the new Hecke setting, not the
literature-hypothesis record.
\begin{proof}
An away index retains the same split place, section,
uniformizer, matrix and finite-coset sum.  An ordinary
index retains the same local group and positive element.
Thus the operator identity is literal, with \(j\) only
changing the membership proof of an index.  HLM17 and
the unchanged \(\theta,\tau_0\) identify the full
ordinary modules.  Selecting the new residual part on
that module amounts to selecting only the old operators
indexed by \(j\).

The new inverse-Frobenius image is dense by FNL16.
Its image under the continuous surjection \(p_0\) is
dense in the old group.  By (HLM.14) these are exactly
the retained old Frobenius arguments; their residual
polynomials are the old ones by HLM16.  The admitted
dense weights and actual parts are supplied by AFS5
and UPA17 for the old data.  HLM19 consequently
identifies the selected part with the full old part.
This proves (HLM.17), including the zero-module case.
It supplies both statements of the original
\texttt{step2\_\allowbreak alt}, rather than an unspecified module
isomorphism at maximal compact level.
\end{proof}
% END HLM20

% BEGIN HLM21
\paragraph{HLM21. Exact comparison of the determinant pairs.}
\label{claim:HLM21}
Suppose \(T',T\in\mathcal C_O\) are the new and old
nonzero Hecke algebras, with their \(\mathrm{IsUniv}\)
pairs \((D',\psi')\), \((D,\psi)\), and
\(\Phi:T'\to_A T\) preserves every surviving generator.
Then
\[
          D'\mathbin{\mathrm{map}}\Phi=D\circ p_0,
       \qquad \Phi_*\psi'_{v,i}=\psi_{v,i}.
 \tag{HLM.18}
\]
\begin{proof}
An \(O\)-map of these local coefficient rings preserves
the residue map, and therefore sends every power of the
maximal ideal into the corresponding power.  Thus the
pushed determinant and characters remain adically
continuous.  At every new away Frobenius, the two
determinant polynomials in (HLM.18) agree by the generator
identity and (HLM.14), with the same \(q\), coefficient
index and inverse Frobenius.  Density, comparison modulo
each ideal power, and separatedness give equality of
all characteristic coefficients on the whole new group.
HLM6 gives equality as full determinant laws.

The two local characters have the same unit equations
because \(\Phi\) preserves \(A\) and \(\theta\).
Their prefix products at the original uniformizer agree
because all ordinary indices are retained.  Cancel
successive prefix products and the common cyclotomic
factor to get equality of the individual values.
Units and the uniformizer give equality on the local
multiplicative group.  Reciprocity, Weil density and
adic separatedness give equality on the whole
decomposition group.  This is the full original
\texttt{level\allowbreak UR2Base\allowbreak Eq\_\allowbreak alt}, with its original maps.
\end{proof}
% END HLM21

% BEGIN HLM22
\paragraph{HLM22. The old-level identification respects the universal-ring map.}
\label{claim:HLM22}
Let \(B\) be the given old universal ring, \(B'\) the
new one, and \(\beta:B'\to_OB\) carry its pair to
the pullback of the old pair along \(p_0\), with the
same local characters.  The identity of function modules
in (HLM.17) intertwines the new universal action with
the old action through \(\beta\).  Therefore it gives
the required \(\beta\)-equivariant additive equivalence
between the new maximal-compact fixed space and the old
base module.
\begin{proof}
If the common module is zero the conclusion follows
directly.  Otherwise its two Hecke algebras are in
\(\mathcal C_O\).  Restricting the new generated
algebra of endomorphisms gives an \(A\)-algebra map
\(\Phi:T'\to T\); (HLM.17) identifies each generator
with the corresponding old one.  It need not be
assumed surjective.  Let \(f':B'\to T'\) and
\(f_0:B\to T\) be the two specified maps.  HLM21
shows that \(\Phi f'\) and \(f_0\beta\) take the
new universal pair to exactly the same determinant
and local characters.  The unique-map property of
\(B'\) therefore gives
\(\Phi f'=f_0\beta\).  By definition of \(\Phi\),
this equality is equality of the two actions on the
same functions.  HLM5 identifies the maximal-compact
fixed space with that function submodule, and HLM10
identifies its action with \(f'\).  Their composite
is the asserted additive equivalence and satisfies
the exact scalar equation of \(\mathrm{LevMod.eM\_\allowbreak smul}\).
This proves \texttt{level\allowbreak UR2Smul\allowbreak Key\_\allowbreak alt} without
changing either ideal or coefficient action.
\end{proof}
% END HLM22

\subsection{One base and the complete all-level package}

% BEGIN HLM23
\paragraph{HLM23. Choose the original Taylor--Wiles base before every auxiliary input.}
\label{claim:HLM23}
Fix the given old \(\mathrm{IsJ}\), its ideal \(J\subset B\),
and its level \(K_0\), with \(V_0.S'=s.S\).  Put
\(D=n[L^+:\mathbb Q]\) and
\(S=O[[\mathrm{Empty}\sqcup\mathrm{Fin}(D)]]\).
There is an original
\(b:\mathrm{TWBase}(\mathrm{pairOf}(B.u),J,\mathrm{Empty},D)\)
whose module is the same old \(M_{K_0}\), with its exact
universal action.  This choice depends only on the fixed
old data, and is made before \(V,h,q,v,w,B_q\).
\begin{proof}
The original full local reciprocity input and KX3 give
coordinates on each local pro-\(\ell\) abelianization,
with inertia the \(d_v\) ramified coordinates and
\(\sum_vd_v=[L^+:\mathbb Q]\).  KX4 fixes extensions
\(\psi^0_{v,i}\) with the prescribed residual and torsion
values.  KX5 identifies the given weight ring with the
power series ring on the triples \((v,i,j)\),
\(0\leq i<n\), \(1\leq j\leq d_v\).  Choose an
inertia lift \(\gamma_{v,j}\) of each coordinate basis
vector.  Enumerate the \(D\) triples once, writing
\(v(k),i(k),\gamma_k\), and put
\(c_k=\psi^0_{v(k),i(k)}(\gamma_k)^{-1}\in O^\times\).
The universal character coordinate formulas give an
\(O\)-algebra isomorphism \(e:S\simeq A\) with
\[
 e(X_k)=c_k\kappa_{v(k),i(k)}(\gamma_k)-1,
 \qquad
 \bar c_k\bar\psi_{v(k),i(k)}(\gamma_k)=1.
 \tag{HLM.19}
\]
This is the source's \texttt{weight\_\allowbreak variables\_\allowbreak exist}
for arbitrary rank \(n\); the number of variables is
not specialized to four.

The given \(\mathrm{IsJ}\) supplies
\(\mu:B\to\operatorname{End}_A(M_{K_0})\).
In its nonzero branch it is the given surjective Hecke
map followed by the faithful inclusion into endomorphisms,
so \(\ker\mu=J\).  In its zero branch it is the unique
zero-module action and \(J=B\), giving the same equality.
Weight compatibility gives
\(\mu(\Lambda.\mathrm{toPair}(B.u)(a))m=am\).
Transport the \(A\)-action along \(e\) and put
\(\lambda_b=\Lambda.\mathrm{toPair}(B.u)\circ e\).
Then (HLM.19) and the inertia restriction identity give
\[
 \lambda_b(X_k)=c_kB.u.\psi_{v(k),i(k)}(\gamma_k)-1.
 \tag{HLM.20}
\]
The resulting \(S\)- and \(B\)-actions have the
required scalar equality.  OHA8 gives finite freeness
over \(A\), hence over \(S\) by transporting a basis.
Finally
\(\operatorname{Ann}_B(M_{K_0})=\ker\mu=J\): an
element is in either side exactly when it acts as zero
on every vector.  Thus for every prime ideal \(\mathfrak p\)
the source's support equivalence is literally
\(J\subset\mathfrak p\Longleftrightarrow J\subset\mathfrak p\).
These verifications give every \(\mathrm{TWBase}\) field,
including zero rank for a zero module.  They are the
general-rank construction underlying JSP8, with the
same ideal rather than a radical or enlargement of it.
\end{proof}
% END HLM23

% BEGIN HLM24
\paragraph{HLM24. Construct the actual package for arbitrary auxiliary data.}
\label{claim:HLM24}
Keep \(b\) fixed as in HLM23.  Now take any of the
\(V,h,q,v,w,B_q\) allowed in the literal
\(\mathrm{LevelAlt}\) quantifiers: \(q\in\mathbb N\),
the \(v_j\) are distinct and outside \(s.S\),
\(V.S'=s.S\cup\operatorname{range}v\), the indicated
valuation subrings lie over both \(w_j\) and \(v_j\),
and all the stated quotient-map and split-decomposition
conditions hold.  There are an original coefficient ring
\(R=B'\in\mathcal C_O\), pair \(p\) on the prescribed
\(r_N\), map \(\beta:B'\to_OB\), original union
\(\mathcal M\), commuting actions \(\alpha,\rho\),
and \(\lambda:A\to_OB'\), with the following properties:
all compact-level fixed modules are finite, every normal
compact-level quotient has the projectivity of HLM14,
the maximal-compact fixed module is \(b.M\) through the
\(\beta\)-equivariant map of HLM22, and every nonzero
vector has all the points in HLM13.
\begin{proof}
For an integral element of \(L^+\), membership in the
place defined by \(B_j\) can be tested after either
field embedding.  The two given place-over assertions
therefore imply \(\operatorname{under}(w_j)=v_j\).
They give exactly the original \(Q=\mathrm{qUp}(V,w)\),
with index set \(\mathrm{Fin}(q)\), distinct underlying
places, outside \(s.S\), and the stated range equality.
In particular \(V_0.S'\subset V.S'\).  All original
split-decomposition conditions remain among the hypotheses;
the level record's split condition supplies the sections
used above.

HLM15 pulls the actual pair \(B.u\) back to a pair
\(p'\) for \(r^{\mathrm{up}}\) on \(B\).  This
is the required nonemptiness witness for KX8, not an
assumed existence of an automorphic auxiliary point.
The global and local finite-cohomology inputs in the
original \(\mathrm{HypothesesH}'\) supply KX6 at the
actual finite set \(V.S'\); the original local
reciprocity and duality inputs supply its KX3--KX4
local interfaces.  KX7--KX8 therefore construct the
actual universal ordinary pair ring \(B'\) for
\(r^{\mathrm{up}}\).  Apply its universal property
to \(p'\) to obtain \(\beta\), with
\[
 \beta_*B'.u.D=B.u.D\circ p_0,
 \qquad \beta_*B'.u.\psi=B.u.\psi.
 \tag{HLM.21}
\]
The pair \(p=\mathrm{pairN}(B'.u)\) on \(r_N\)
is supplied by HLM15.

HLM16 constructs the actual new unitary data using
exactly \(\tau_0,\theta,a\circ j,A,\kappa\).
Put \(\lambda=\Lambda^{\mathrm{up}}.\mathrm{toPair}(B'.u)\).
Its inertia identity says
\[
       \lambda(\kappa_{v,i}(\gamma))=B'.u.\psi_{v,i}(\gamma)
                    \quad(\gamma\in I_v).
 \tag{HLM.22}
\]
FNL16 supplies its original frame axioms, AFS5 all
dense weights, UPA17 the actual automorphic parts with
their original literature inputs, NTL17 every neat
level, and HLM17 section faithfulness.  HLM9 supplies
an \(\mathrm{IsJ}\) at every one of these levels for
this \(B'\), including the zero cases.  The unique-map
hypothesis of HLM14 is the universal property of \(B'\):
two maps with the same pair correspond to the same pair
and hence coincide.  Apply HLM14 at the admitted zero
weight, using the fixed \(\iota\).  It produces all
the stated union, action, finiteness, projectivity and
point data.  HLM20--HLM22 give the specified
maximal-compact identification, then HLM23 identifies
its target with the already fixed \(b.M\).
This proves every raw output of \texttt{Base\allowbreak Object\allowbreak Hom\allowbreak Ext.\allowbreak s3}
with actual suppliers rather than a new all-level
\(\mathrm{IsJ}\) or automorphic-set assumption.
\end{proof}
% END HLM24

% BEGIN HLM25
\paragraph{HLM25. Every \texttt{Lev\allowbreak Mod} field and full classicality.}
\label{claim:HLM25}
The data of HLM24 give an original
\(\mathcal A:\mathrm{LevMod}(b,h,w)\) satisfying
\(\mathcal A.\mathrm{ClassicalAlt}\).
\begin{proof}
For a local field \(k\) and subgroup \(C\subset k_0^\times\),
the original \(\operatorname{iwC}(k,n,C)\) consists
of integral invertible matrices with upper triangular
reduction and each diagonal residue in \(C\).  It
contains the first principal congruence subgroup, hence
is open; it is a closed subgroup of the compact maximal
compact group, hence compact.  The diagonal entries of
the product of two invertible upper triangular matrices
multiply coordinatewise.  Thus reduction followed by
the diagonal modulo \(C\) is a homomorphism from the
Iwahori group, whose kernel is this subgroup.  It is
therefore normal in the full Iwahori.  For the original
\(C=\operatorname{ellC}(\ell,k_0^\times)\), consisting
of elements of order prime to \(\ell\), take finite
products to obtain the exact normal inclusion
\(U_{\mathrm{Iw},1}\triangleleft U_{\mathrm{Iw}}\).
Their quotient is finite because an open subgroup of
a compact group has finite index.  This verifies the
level assertions also for \(q=0\).

Give \(\mathcal M\) its \(R=B'\)-module structure
through \(\alpha\).  Commutation in (HLM.10) makes
each \(\rho(g)\) \(R\)-linear, without changing its
underlying action on functions.  Give \(R\) its
\(S\)-algebra structure by \(\lambda\circ e\), and
give \(\mathcal M\) its \(S\)-action through \(e\)
and the old \(A\)-action.  Then
\(\alpha(\lambda(e(a)))m=e(a)m\) proves the scalar
tower identity; the \(O\)-tower follows because both
maps are \(O\)-linear.  Equations (HLM.19) and
(HLM.22) give the exact required coordinate equation
\[
  X_k\longmapsto c_k p.\psi_{v(k),i(k)}(\gamma_k)-1.
 \tag{HLM.23}
\]

All fixed subgroups are unchanged by viewing the same
action as \(R\)-linear.  HLM24 gives finiteness over
\(A\) at \(U_{\mathrm{Iw},1}\); transporting
generators along \(e:S\simeq A\) gives finiteness
over \(S\).  Its quotient representation is projective
over \(A[U_{\mathrm{Iw}}/U_{\mathrm{Iw},1}]\).
The coefficient isomorphism induced by \(e\) identifies
this group algebra with the corresponding \(S\)-group
algebra, so it transports a projective splitting and
the quotient action to the exact projectivity field.
This is a ring isomorphism transport, not descent along
a nonflat map.  Its action on every quotient coset is
the original \(\rho(u)\) by HLM5.  The map \(e_M\)
is the HLM24 maximal-compact equivalence followed by
the fixed base identification; HLM22 gives its
\(\beta\)-scalar identity.

For clarity the complete non-instance field dictionary is
\[
\begin{array}{c|l}
\text{field}&\text{constructed value or verification}\\ \hline
r_N,\ \mathrm{compat}&\text{HLM15's original residual object and compatibility}\\
R,\ p& B',\ \mathrm{pairN}(B'.u)\text{ from HLM24}\\
\mathrm{segmentWeightExponent}&\text{equation (HLM.23)}\\
\beta,\ \beta\_\allowbreak \mathrm{pair}&\text{equation (HLM.21), with unchanged local characters}\\
P,\rho&\mathcal M\text{ with the same right-translation action, now }R\text{-linear}\\
\mathrm{fin},\ \mathrm{proj}&\text{the preceding fixed-space and group-algebra arguments}\\
e_M,\ e_M\_\allowbreak \mathrm{smul}&\text{the original-function identification of HLM22}.
\end{array}
\tag{HLM.24}
\]
The ring, algebra, module and tower instance fields were
constructed explicitly above.  Finally suppose a vector
maps to zero under every \(\mathrm{LevPtAlt}\) at the
fixed \(\iota\).  If it were nonzero, HLM13 would
give a point with a nonzero image.  Its determinant,
continuity, scalar, group, local-isotypic and every-bridge
fields are exactly those of \(\mathrm{LevPtAlt}\),
because the underlying actions and pair have not changed.
This is a contradiction.  Thus \(\mathrm{ClassicalAlt}\)
holds.  The argument is the full
\texttt{level\allowbreak UR2Mod\allowbreak Pack\_\allowbreak alt} in the source file named
\texttt{Nonempty\allowbreak Unitary\allowbreak Hecke\allowbreak Parts}; it does not conflate
that packing theorem with the separate automorphic
nonvanishing theorem used in the ordinary return argument.
\end{proof}
% END HLM25

% BEGIN HLM26
\paragraph{HLM26. The whole literal all-level \texttt{Level\allowbreak Alt} supplier.}
\label{claim:HLM26}
Assume the original \(\mathrm{HypothesesH}'\).
For every original coefficient prime \(\ell\), CM
number field \(L\), algebraic closure \(\Omega\),
Hida setting \(s\), rank \(n\geq1\), original level
\(V_0\) and \(Q_0\), coefficient object \(c\),
residual object \(r\), fixed
\(\iota:\overline{\mathbb Q}_\ell\simeq\mathbb C\),
universal ring \(B\), weight object \(\Lambda\),
\(\mathrm{UnitaryHeckeData}\ \mathscr R\), and old
auxiliary level \(K_0\), suppose \(V_0.S'=s.S\).
For every ideal \(J\subset B\) with
\(\mathscr R.\mathrm{IsJ}(B,K_0,J)\), the original
\(\mathrm{Pa.LevelAlt}(B,J)\) holds.
\begin{proof}
Choose \(b\) by HLM23, before introducing any auxiliary
data.  For every \(V,h,q,v,w,B_q\) and every list of
proofs satisfying exactly the hypotheses printed in
HLM24, construct its package by HLM24 and its
\(\mathrm{LevMod}\) and \(\mathrm{ClassicalAlt}\)
proof by HLM25.  These are the existential witness
and the universal continuation in the definition of
\(\mathrm{LevelAlt}\), with
\(D=n[L^+:\mathbb Q]\) and no auxiliary variables
in the fixed base.  The coefficient map, pair,
residual object, type, weight and maximal-compact
identification are exactly the ones checked above.
All primes and all ranks in the original statement
remain allowed, including \(\ell=3\); none of the
projectivity arguments divides by a group order or
by \(n!\).

If \(q=0\), the product group is trivial, the
selected place list and all its local tests are empty,
and HLM17--HLM22 identify the sole fixed module with
the old one.  If the old module is zero, HLM23 uses
\(J=B\) and a zero free base, exactly as its
\(\mathrm{IsJ}\) branch requires.  This does not
assert that every auxiliary union is zero; the proof
of HLM25 still separates every vector of any such
union.  Only when a particular union is zero is its
classicality assertion vacuous.  Thus neither edge
case has been removed from the quantifiers.

The conclusion is the complete original
\texttt{level\_\allowbreak of\_\allowbreak is\allowbreak J\_\allowbreak unitary\allowbreak Hecke\allowbreak Data}, retaining
the local, automorphic and cohomological foundations
in \(\mathrm{HypothesesH}'\) and in its inherited
proofs.  The later conversion of this family to
\(\mathrm{TWAut}\) and \(\mathrm{HasAut}\) is a
separate consumer; no such conclusion is inferred
merely from the present family.
\end{proof}
% END HLM26


\section{From actual level modules to the full Taylor--Wiles automorphy datum}
\label{sec:hida-taylor-wiles-v45}

This chapter proves the finite consumers between the original
\(\mathrm{LevelAlt}\) and \(\mathrm{HasAut}\).  Its only new
representation-theoretic premises are precisely the three fields
of the original \(\mathrm{LocPAlt}\): Bernstein--Zelevinsky
finiteness, the geometric lemma, and the tame local theorem.
The full \(\mathrm{IsRecT}\) correspondence laws, its existence,
and its generic and spherical good-parameter consequences have
the scope already reconstructed in GUA.  The actual level module
is supplied by HLM, with its vector-separating classical maps,
full polynomial-law comparisons, and one fixed base module.

Here \(\Lambda=O[[Z_1,\ldots,Z_D]]\), where the empty set of
additional variables is understood and \(D=n[L^+:\mathbb Q]\).
It is a complete noetherian local ring.  All modules described as
finite below are finite over this ring unless another ring is
specified.  The fixed base is \(b\), its module is \(M_0=b.M\),
and its coefficient map is \(\lambda:\Lambda\to B\).  In
particular, for every prime \(\mathfrak p\) of \(B\),
\[
 \operatorname{Ann}_B(M_0)\subseteq\mathfrak p
       \quad\Longleftrightarrow\quad J\subseteq\mathfrak p .
 \tag{HTW.1}
\]
This equivalence is retained throughout; none of the intermediate
localizations changes the chosen ideal \(J\).

At \(q\) auxiliary places write \(Q_j>1\) for the residue
cardinality, \(f_j\) for geometric Frobenius, and \(t_j\) for
the original tame generator.  Indices \(i\) run from \(0\) to
\(n-1\).  For commuting variables \(U_{ji}\) put
\[
 F(U)=\prod_j\prod_i\prod_{k\ne i}
             (U_{ji}-U_{jk})(U_{ji}-Q_jU_{jk}).
 \tag{HTW.2}
\]
This is the literal source polynomial \(\mathrm{fEl}\), with
ordered pairs, not its square-free part.  In the group algebra
define
\[
 \begin{split}
 C_j&=\prod_i([f_j]-U_{ji}),\\
 C_{ji}'&=\left(\prod_{k\ne i}([f_j]-U_{jk})\right)
                                      ([t_j]-1-Y_{ji}).
 \end{split}
 \tag{HTW.3}
\]
The order displayed in the second expression is retained.

\subsection{Finite tame level and the two local relations}

% BEGIN HTW1
\paragraph{HTW1. The tame quotient and its exact norm sequence.}
\label{claim:HTW1}
Let \(I_j\) be the original Iwahori subgroup and \(I_j^1\)
its subgroup with diagonal reductions in the prime-to-\(\ell\)
subgroup of the residue-field units.  If
\(a_j=v_\ell(Q_j-1)\), then
\[
 \Delta=\prod_j I_j/I_j^1
       \simeq\prod_{j,i}\mathbb Z/\ell^{a_j}\mathbb Z .
 \tag{HTW.4}
\]
The original \(\delta_j\) gives the generator in each diagonal
coordinate.  For any projective \(\Lambda[\Delta]\)-module \(M\),
the norm \(N_\Delta=\sum_{\gamma\in\Delta}\gamma\) induces
an isomorphism \(M_\Delta\simeq M^\Delta\).  Its kernel on \(M\)
is \(\sum_{j,i}(\gamma_{ji}-1)M\).
\begin{proof}
Reduction of the diagonal entries is surjective onto the
product of residue-field unit groups; its kernel contains the
unipotent factors and the principal congruence group.  Taking
the specified prime-to-\(\ell\) diagonal subgroup leaves the
cyclic \(\ell\)-primary quotient in each coordinate.  A finite
field's multiplicative group is cyclic, and a generator of
its quotient has order \(\ell^{a_j}\).  The defining
\(\mathrm{IsGenMod}\) property of \(\delta_j\) says exactly
that its class is such a generator.  Different places and
different diagonal coordinates commute in this quotient.

For the regular module \(\Lambda[\Delta]\), a fixed vector
has all coefficients equal and is uniquely a scalar multiple
of \(N_\Delta\).  The norm of \(\sum c_\gamma[\gamma]\)
is \((\sum c_\gamma)N_\Delta\).  Its kernel is the
augmentation ideal: subtract the coefficient sum at \(1\)
to express it as \(\sum c_\gamma([\gamma]-1)\).
Writing every \(\gamma\) as a product of the indicated
generators and using \(ab-1=a(b-1)+(a-1)\) proves the
generator description.  These assertions hold for free
modules coordinatewise and pass to a direct summand.
This proves the projective case without division by
\(|\Delta|\).  Finally \(Q_j\equiv1\bmod\ell^N\) implies
\(a_j\ge N\), also when \(N=0\).
\end{proof}
% END HTW1

% BEGIN HTW2
\paragraph{HTW2. The actual commuting torus operators and trace.}
\label{claim:HTW2}
For an original level module \(A\), put
\(M_1=A.P^{\prod I_j^1}\) and \(W=A.P^{\prod I_j}\).
There are commuting invertible torus actions \(\Theta_1,\Theta_w\)
with the original dominant double-coset values.  Define
\[
 U_{ji}=Q_j^i\Theta_1(e_i(x_j)),\qquad
 1+Y_{ji}=\Theta_1(e_i(\delta_j)),
 \tag{HTW.5}
\]
and use the same \(U\)-formula on \(W\).  The actual
double-coset trace \(M_1\to W\) is the norm in HTW1.
It is surjective, its kernel is \(\sum Y_{ji}M_1\), and
it intertwines every \(U_{ji}\).
\begin{proof}
All \(Q_j\) are units in \(O\), since the places avoid
\(\ell\).  The positive diagonal double cosets multiply
according to addition of their ordered valuation vectors.
This follows by writing an Iwahori element in upper
unipotent, diagonal-unit, and lower unipotent order:
the summation indices for a product are the disjoint
successive residue quotients of each root subgroup.
At the tame subgroup the simple-reflection relation
retains its diagonal-unit terms.  For a nonzero residue
class \(c\), choose a unit lift \(a_c\) and let \(d_c\)
have entries \(a_c,-a_c^{-1}\) in the two adjacent
coordinates and \(1\) elsewhere.  Write \(D_c\) for
its normalizer action.  The elementary two-by-two
matrix factorization of \(s(1+a_cE_{i,i+1})s\) gives
\[
 T_s^2=Q_j+T_s\sum_{c\ne0}D_c,\qquad
       T_sD_c=D_c^{\,s}T_s .
\]
Consequently \(Q_j^{-1}(T_s-\sum D_c)\) is a right
inverse and \(Q_j^{-1}(T_s-\sum D_c^{\,s})\) is a
left inverse.  They agree, so \(T_s\) is a unit.
Only on the ordinary Iwahori fixed space do all
\(D_c\) act as \(1\), giving
\(T_s^2=(Q_j-1)T_s+Q_j\).

Reduced permutation products are therefore units at
either level.  For a fundamental diagonal translation
with uniformizer in the first \(m\) entries, rotating
those entries to the end makes its product with the
rotation permutation normalize \(I_C\).  Indeed
conjugation changes the valuation bound on matrix
entry \((i,k)\) by the difference of these two
zero-or-one exponents, exactly compensating the
change of whether the entry is below the diagonal;
it permutes the prescribed diagonal residue subgroup.
The translation double coset is thus a product of
a permutation unit and a normalizer unit.
An ordered integral valuation vector is a sum of
these fundamental vectors and a constant vector.
Positive multiplication, and the normalizer actions
of diagonal units and central scalars, prove that
every dominant diagonal operator is invertible.
Every diagonal element is a quotient of two
dominant ones: add a sufficiently strictly ordered
valuation vector to make it and its translate dominant.
The positive multiplication rule therefore extends
uniquely to the torus.  Applying this independently at
each place gives commuting actions and (HTW.5).

The right cosets in \(I_j/I_j^1\) have diagonal-unit
representatives.  Thus the double coset at the identity
is exactly their sum, with no volume factor.  For a
dominant diagonal \(t\), the two expressions obtained
by first summing these units and then its root-subgroup
representatives, or in the other order, enumerate the
same cosets.  The trace intertwines \(\Theta(t)\).
For an arbitrary \(t\), multiply by a dominant \(b\)
such that \(bt\) is dominant and cancel the invertible
\(\Theta(b)\).  Hence it intertwines all torus operators,
including (HTW.5).  The level-module projectivity field
and HTW1 give the asserted kernel and surjectivity.
\end{proof}
% END HTW2

% BEGIN HTW3
\paragraph{HTW3. The uniform local nilpotence exponent is \(n!\).}
\label{claim:HTW3}
Let \(\sigma\) be smooth irreducible over \(\mathbb C\).
For the original subgroup \(I_C\), its fixed space has
dimension at most \(n!\).  A polynomial in the commuting
torus operators which vanishes at every joint generalized
character acts nilpotently with exponent \(n!\).
The same exponent works in an arbitrary isotypic sum of
copies of \(\sigma\).
\begin{proof}
Bernstein--Zelevinsky finiteness makes the fixed space finite
dimensional.  Its map to the Borel Jacquet quotient is
injective: if a fixed vector is a finite sum of unipotent
differences, a sufficiently dominant diagonal contracts
the finitely many unipotents into the fixing subgroup.
The translated vector is then a sum of differences
for elements of that subgroup.  Every root-coset
translate remains in this difference span, so the
positive double-coset operator sends the vector to
both that span and the fixed space.  Normalized
compact averaging kills the former and fixes the
latter; it is a finite average on the finitely
many smooth vectors involved.  Thus this operator
kills the vector.  HTW2 makes it invertible, so
the vector was zero.

If the fixed space is nonzero, commuting invertible torus
operators have a common eigenvector.  To see this, choose
a minimal nonzero invariant subspace; the eigenspace of
any one operator inside it is invariant under all the
others, so minimality makes every operator scalar there.
The image in the Jacquet module is nonzero.
The retained finite-length theorem for that module,
and the fact that smooth irreducible torus modules
are characters, give a finite-dimensional torus
module.  The latter fact follows by an open
stabilizer and the weak Nullstellensatz for the
remaining finitely generated abelian group, as
in HAW7.  Choose a nonzero generalized-character
block.  Its quotient by the maximal ideal of the
commutative image algebra is nonzero by Nakayama;
a nonzero functional on that quotient gives a
character quotient of the entire Jacquet module.
It need not detect the initial eigenvector.
Jacquet reciprocity applied to this quotient
gives a nonzero map from \(\sigma\) into the
corresponding principal series, which is injective
by irreducibility.  The geometric lemma gives the Jacquet
module of that principal series a filtration with one
one-dimensional character for each permutation.  Exactness
of Jacquet coinvariants gives the required bound \(n!\).
This uses the original geometric lemma on this principal
series, rather than any bound on an isotypic multiplicity.

A commuting family on a finite-dimensional complex space
can be simultaneously triangularized, by the same common
eigenvector argument and induction on the quotient.  Its
diagonal characters occur among the joint generalized
characters: equivalently decompose the finite-dimensional
commutative image algebra into its local Artinian factors.
The polynomial in the assertion therefore has zero
diagonal.  A strictly upper triangular matrix on a space
of dimension \(d\le n!\) has \(n!\)-th power zero.

Finally every fixed vector in an isotypic sum is a finite
sum of fixed vectors from actual copies of \(\sigma\).
Indeed start with a finite expression in the images of
intertwiners and average over the compact fixing subgroup;
smoothness makes each average a finite sum.  Intertwiners
commute with the torus action on dominant elements by the
finite double-coset formula, and on all elements by
HTW2's cancellation.  The same polynomial power kills
every summand.  Its exponent does not increase with the
number of copies or the number of auxiliary places.
\end{proof}
% END HTW3

% BEGIN HTW4
\paragraph{HTW4. Tame-kernel triviality from the actual Weil--Deligne comparison.}
\label{claim:HTW4}
Away from \(\ell\), an actual \(\mathrm{IsWDOf}\) comparison
gives equality of characteristic polynomials on the Weil
group.  An element of the kernel of the pro-\(\ell\) tame
character whose matrix has characteristic polynomial
\((X-1)^n\) acts as the identity.
\begin{proof}
On a sufficiently small open inertia subgroup, the
monodromy formula is
\[
                  \rho(\tau)=\exp(t_\ell(\tau)N).
\]
On the
specified tame kernel the parameter \(t_\ell\) is zero.
Some positive power of the given inertia element lies in
that open subgroup, because inertia is compact and the
quotient is finite.  Its matrix therefore has finite
order \(d\).  Its minimal polynomial divides both
\(X^d-1\) and \((X-1)^n\).  The former is separable in
characteristic zero, so the minimal polynomial divides
\(X-1\), and the matrix is \(1\).

For the first assertion, filter by the kernels of
successive powers of \(N\).  The raw Weil operator
preserves this filtration because it conjugates \(N\)
to a nonzero scalar multiple of \(N\); its exponential
correction is the identity on each successive quotient.
The corrected and uncorrected operators therefore
have identical characteristic polynomials.  The
Frobenius semisimplification also preserves these
polynomials, as in the full \(\mathrm{IsWDOf}\)
reading reconstructed in GUA.  Thus the result is the
Weil--Deligne characteristic polynomial at every Weil
element, with the same geometric Frobenius convention.
Coefficient transport uses only the field isomorphism
\(\iota\), never its continuity.
\end{proof}
% END HTW4

% BEGIN HTW5
\paragraph{HTW5. The nonresonant tame matrix identity.}
\label{claim:HTW5}
Suppose a representation \(R\) of a group has, for every \(g\),
\[
 \det(X-R(g))=\prod_i(X-c_i(g)),
\]
where the \(c_i\) are characters.  Put \(P=R(f)\), \(T=R(t)\),
\(\alpha_i=c_i(f)\), \(\beta_i=c_i(t)\).  Assume
\(P^{-1}TP=T^Q\), \(Q\ge2\), and
\(\alpha_i\ne\alpha_k,\ Q\alpha_k\) for \(i\ne k\).
Then
\[
 \prod_{k\ne i}(P-\alpha_k)\,(T-\beta_i)=0 .
 \tag{HTW.6}
\]
\begin{proof}
The character relation gives \(\beta_i^{Q-1}=1\).
Thus \(T^{Q-1}\) is unipotent by its characteristic
polynomial.  Let \(L=\log(T^{Q-1})\), using the finite
nilpotent logarithm over \(\mathbb C\).  The tame relation
gives \(P^{-1}LP=QL\).  The distinct eigenvalues of \(P\)
give an eigenbasis; in it each entry of \(L\) satisfies
\((\alpha_k/\alpha_i-Q)L_{ik}=0\).  All these coefficients
are nonzero: for different indices this is the stated
ordered-pair condition, and for equal indices it is
\(1-Q\).  Consequently \(L=0\), hence \(T^{Q-1}=1\).
The tame relation now gives \(PT=TP\).

In that eigenbasis \(T\) is diagonal, say with entries
\(d_i\).  For \(a=0,\ldots,n-1\), the trace identity at
\(f^at\) says
\(\sum_i\alpha_i^a d_i=\sum_i\alpha_i^a\beta_i\).
The Vandermonde determinant is nonzero, so \(d_i=\beta_i\).
The product in (HTW.6) is zero on every eigenline: its
first factor kills all lines except the \(i\)-th, and its
last factor kills that line.  This proves the ordered
matrix identity.  Rank zero is the unique zero-dimensional
matrix identity and needs no eigenbasis.
\end{proof}
% END HTW5

% BEGIN HTW6
\paragraph{HTW6. The full original local-data consumer.}
\label{claim:HTW6}
The original \(\mathrm{LocPAlt}\) and the full local comparison
at every classical point produce the literal
\(\mathrm{LocalLevelConditions}\), with the given geometric
Frobenius and tame generators.
\begin{proof}
Fix a prolongation and one permissible \(\mathrm{IsRecT}\)
correspondence, whose existence is among the retained
original laws.  The tame local theorem and the original
Artin map send \(f\) to a uniformizer \(x\), send \(t\)
to a unit \(\delta\) generating the indicated residue
quotient, and put \(t^{-Q}f^{-1}tf\) in the pro-\(\ell\)
tame kernel.  This choice is made before the rank,
classical point, irreducible component, or its torus
character is chosen.  After these choices fix an
arbitrary residue-unit subgroup \(C\supseteq C_\ell\),
an arbitrary admissible \(\Theta\) on the corresponding
\(I_C\)-fixed space, and any occurring joint
generalized character \(\chi\).

A joint generalized torus character \(\chi\) on the
component gives, by Jacquet reciprocity, smooth inducing
characters \(\psi_i\), trivial on the prescribed unit
subgroup, with
\[
 \psi_i(z)=Q^{i\,\operatorname{ord}(z)}
     \chi(e_i(z))\,|z|^{(n-1)/2}.
 \tag{HTW.7}
\]
This follows by evaluating the half-modulus character
of the upper Borel at the single diagonal coordinate:
its exponent is \((n-1)/2-i\).
Set \(\alpha_i=Q^i\chi(e_i(x))\).
If all the factors in (HTW.2) at this place are nonzero,
no two inducing characters are related by the modulus
character: evaluating any such relation at \(x\) would
give \(\alpha_k=Q\alpha_i\).  Thus the toral support is
unlinked.  GUA2's full \(\mathrm{PlainGd}\) law applies
to \(\mathrm{algNorm}(\sigma)\).  A spherical component
also has a good parameter by GUA12.
These are exactly the two permitted alternatives in
the local-data contract.  Only after proving this
good-parameter guard do we use the point's conditional
\(\mathrm{WDT}\) comparison to obtain the actual
\(\mathrm{IsWDOf}\) witness.

On the Weil group the corresponding characters after
the algebraic normalization are
\[
 c_i(g)=Q^{i\,\operatorname{ord}(a(g))}
                         \chi(e_i(a(g))).
 \tag{HTW.8}
\]
The full local comparison and HTW4 give the characteristic
polynomial \(\prod_i(X-c_i(g))\) for the actual matrix.
At \(f\) these values are \(\alpha_i\); at \(t\) they
are \(\chi(e_i(\delta))\).  At \(t^{-Q}f^{-1}tf\)
all the characters equal \(1\), by their triviality
on the specified unit subgroup.  HTW4 kills the actual
matrix of this element.  HTW5 gives (HTW.6).
In the spherical alternative only the characteristic
polynomial assertion is needed; in the nonresonant
alternative both assertions hold.

The argument uses the classical point's local comparison
for every permissible correspondence and every
prolongation.  Therefore it applies to the selected one
without an extra supplier.  Taking the already supplied
vector-separating classical family completes precisely
\(\mathrm{LocalLevelConditions}\).
\end{proof}
% END HTW6

% BEGIN HTW7
\paragraph{HTW7. Descent of polynomial nilpotence through classical separation.}
\label{claim:HTW7}
Let \(H\in R[U_i,Y_i]\) concern one auxiliary place.
Suppose every classical specialization and every occurring
joint generalized character makes it zero after the
substitution
\[
 U_i=Q^i\chi(e_i(x)),\qquad
 Y_i=\chi(e_i(\delta))-1 .
 \tag{HTW.9}
\]
Then \(H^{n!}\) acts as zero on \(M_1\).
\begin{proof}
Each classical map is additive, coefficient-semilinear,
and equivariant for the actual group action.  Finite
double-coset summation therefore intertwines every
dominant torus operator.  HTW2 extends this to all torus
elements, and polynomial induction extends it to \(H\).
Apply HTW3 to the classical target, which is an isotypic
sum at the chosen place.  Formula (HTW.9) is exactly
the eigenvalue substitution in (HTW.5); hence the
specialization of \(H^{n!}m\) is zero for every point.
The \(\mathrm{ClassicalAlt}\) field separates actual
module vectors, so \(H^{n!}m=0\).  The conclusion is
integral over the unchanged \(R\).
\end{proof}
% END HTW7

% BEGIN HTW8
\paragraph{HTW8. Both exact trace relations with the same exponent.}
\label{claim:HTW8}
For every \(j,i,g\), on \(M_1\) one has
\[
 \big(F(U)D.\operatorname{tr}([g]C_j)\big)^{n!}=0,\qquad
 \big(F(U)D.\operatorname{tr}([g]C'_{ji})\big)^{n!}=0 .
 \tag{HTW.10}
\]
\begin{proof}
First use only the factor \(F_j\) at place \(j\).
At a joint character either \(F_j=0\), making the
required scalar zero, or HTW6 gives the characteristic
polynomial of \(P\) and the matrix identity (HTW.6).
Cayley--Hamilton then kills \(C_j(P)\); (HTW.6) kills
the ordered \(C'_{ji}(P,T)\).  Left multiplication
by the arbitrary matrix \(R(g)\) and taking the trace
preserves zero.  The point's equality of full determinant
laws specializes their trace polynomial to exactly this
matrix trace, over the polynomial scalar algebra.
HTW7 proves both \(n!\)-th power relations for \(F_j\).

The remaining factors of \(F\) commute with these
operators, so multiplying the zero power by their
\(n!\)-th powers gives (HTW.10).  There is one
exponent \(n!\), rather than an exponent multiplied by
the number of places.  In rank zero the degree-zero
determinant has zero trace and both assertions hold
directly.  No division by \(n!\) occurs.
\end{proof}
% END HTW8

% BEGIN HTW9
\paragraph{HTW9. The symmetric operator identities on the original base module.}
\label{claim:HTW9}
Under the inclusion \(M_0\to W\), the elementary symmetric
polynomials in \(U_{j0},\ldots,U_{j,n-1}\) act by
\((-1)^kP_j{}_{\,n-k}\), where
\[
 P_j(X)=D_0.\chi([h(f_j)])\in B[X].
 \tag{HTW.11}
\]
\begin{proof}
The original level-module equivalence identifies \(M_0\)
with the vectors fixed by all maximal integral compact
subgroups.  A nonzero image of such a vector at a
classical point lies in a spherical component.
HTW6 gives its Frobenius polynomial with roots
\(Q_j^i\chi(e_i(x_j))\).  The elementary symmetric
functions in those roots are the coefficients in
(HTW.11) with sign \((-1)^k\).

The torus action on the spherical vector obeys the same
identity by the normalized Satake calculation already
given in BCS15 and UPA5--UPA7; equivalently it follows
by evaluating the elementary symmetric polynomial in
the Iwahori character decomposition.  The Hecke and
determinant scalar actions agree at the same classical
point by its full determinant-law field.  Vector
separation descends the equality to the actual fixed
module.  This holds for \(0\le k\le n\), with monic
degree \(n\) polynomials, and for the original
coefficient map \(\beta:R\to B\).
\end{proof}
% END HTW9

\subsection{Integral comparison maps and the exact correction factor}

% BEGIN HTW10
\paragraph{HTW10. An explicit integral divided-difference dual family.}
\label{claim:HTW10}
In \(O[X_0,\ldots,X_{n-1}]\) set
\[
 d=\prod_{i<k}(X_k-X_i),\quad
 N_Q=\prod_{i<k}(X_k-QX_i),\quad
 \partial f=d^{-1}\sum_{w\in S_n}\operatorname{sgn}(w)w(f).
 \tag{HTW.12}
\]
The last expression is a polynomial.  There are explicit
finite families \(a_T,b_T\), defined over \(\mathbb Z\),
such that
\[
 \sum_T a_Tw(b_T)=
     \begin{cases}d&w=1,\\0&w\ne1,\end{cases}
 \qquad
 \sum_T a_T\partial(b_Tf)=f .
 \tag{HTW.13}
\]
The same expansion holds with \(a,b\) interchanged.
\begin{proof}
An alternating polynomial vanishes on \(X_i=X_k\);
successive division by the monic differences proves
divisibility by \(d\), first over the integral universal
polynomial ring and then after coefficient substitution.
Let \(\mathcal P=\{(i,k):i<k\}\), and index by its subsets.
Put
\[
 a_T=\prod_{(i,k)\in T}X_k,\qquad
 b_T=\prod_{(i,k)\in\mathcal P\setminus T}(-X_i).
\]
Expanding a product gives
\(\sum_Ta_Tw(b_T)=\prod_{i<k}(X_k-X_{w(i)})\).
For \(w=1\) this is \(d\).
For \(w\ne1\), some \(i<w(i)\); otherwise summing
\(w(i)\le i\) forces equality everywhere.  The pair
\((i,w(i))\) makes the product zero.
Multiply the second identity in (HTW.13) by \(d\),
expand the alternating sum, and use the first identity:
only \(w=1\) survives and the result is \(df\).
Cancellation is valid in the integral universal ring,
so the identity is integral.  Applying \(w\) to the
identity for \(w^{-1}\) proves the interchanged version.
This construction inverts neither \(n!\) nor \(d\)
in the coefficient ring of any module.
\end{proof}
% END HTW10

% BEGIN HTW11
\paragraph{HTW11. The actual Iwahori trace identities.}
\label{claim:HTW11}
Let \(E_j:W\to W\) be the finite trace from Iwahori
fixed vectors to the maximal compact fixed vectors
at place \(j\), followed by inclusion.  It commutes
with symmetric polynomials at \(j\), and
\[
 \sum_T a_T(U_j)E_jb_T(U_j)=N_{Q_j}(U_j).
 \tag{HTW.14}
\]
If \(v\) is maximal-compact fixed at \(j\), then
\[
 E_j z(U_j)v=\partial(N_{Q_j}z)(U_j)v .
 \tag{HTW.15}
\]
Traces and operators at different places commute and
preserve the other maximal-compact fixed spaces.
\begin{proof}
Here are the finite Hecke calculations behind the
identities.  Bruhat decomposition of the reduction of
the maximal compact gives one Iwahori double coset
for each \(w\in S_n\); its trace is \(E=\sum_wT_w\).
For an adjacent transposition \(s=(i,i+1)\), enumerating
the root-subgroup representatives gives
\[
 \begin{split}
 T_s^2&=(Q-1)T_s+Q,\\
 T_sX_i&=X_{i+1}T_s+(Q-1)X_i,\\
 T_sX_{i+1}&=X_iT_s-(Q-1)X_i .
 \end{split}
 \tag{HTW.16}
\]
It commutes with the other \(X_k\).  In the first
enumeration one term returns to the identity and the
remaining nonzero residue classes give \(T_s\);
inserting the diagonal translation before or after
these representatives gives the last two formulas.
The scalar normalization \(U_i=Q^i\Theta(e_i(x))\)
is the one making them (HTW.16).  Reduced products
give \(T_w\); a length-increasing product enumerates
each corresponding Bruhat cell once.

Over the fraction field in the \(X_i\), these relations
are represented by
\[
 T_sf=Q\,s(f)+(Q-1)X_i
                        \frac{f-s(f)}{X_i-X_{i+1}}.
 \tag{HTW.17}
\]
The numerator is divisible by the displayed difference.
This representation is faithful on the Bernstein span:
the highest permutation term of \(T_w\) has nonzero
coefficient, and the distinct permutations are linearly
independent as automorphisms of the rational function
field.  The latter follows by the shortest-relation
argument: translate a supposed relation by multiplication
by a variable on which two automorphisms differ and
subtract.  The double-coset basis realizes this same
span in the actual Hecke algebra; its distinct
translation and Bruhat terms have disjoint leading
double-coset support.

Summing (HTW.17) over reduced Bruhat products gives
\[
 Ef=\sum_w\frac{w(N_Qf)}{w(d)}
                    =\partial(N_Qf).
 \tag{HTW.18}
\]
For completeness, its coefficient at the longest
permutation is the product of the reversed adjacent
factors \((X_i-QX_k)/(X_i-X_k)\).  The relations
\(T_sE=QE\) make the output symmetric, so permuting
this coefficient determines all the other coefficients.
This proves (HTW.18) and also verifies the normalization
at \(n=1\), where both products are \(1\).

Now HTW10 applied inside (HTW.18) gives
\(\sum_Ta_TEb_T=N_Q\) as operators on the faithful
model, hence in the Hecke algebra.  This proves
(HTW.14) on every actual module.  On a vector satisfying
\(T_sv=Qv\), the same calculation with (HTW.16)
gives \(Ezv=\partial(N_Qz)v\), proving (HTW.15).
All identities are polynomial identities over the
integral coefficient ring with \(Q\) inverted;
their verification in its fraction field therefore
does not invert any integer in the actual module.
Different places act in different factors of the
product group.  Finite sums and equivariance give the
remaining commutations and preservation assertions.
\end{proof}
% END HTW11

% BEGIN HTW12
\paragraph{HTW12. The two maps before completion.}
\label{claim:HTW12}
Let \(S\) be the simultaneous ordered splitting algebra
of the monic polynomials \(P_j\) over \(B\), with roots
\(u_{ji}\).  There are polynomial-semilinear maps
\[
 \phi:S\otimes_BM_0\longrightarrow W,\qquad
 \psi:W\longrightarrow S\otimes_BM_0
\]
whose two compositions are multiplication by
\(f_0=\prod_jN_{Q_j}(U_j)\), evaluated on the respective
sides.
\begin{proof}
Write \(\mathrm{inc}:M_0\hookrightarrow W\) for the
original inclusion through the level equivalence.
The product of all \(E_j\) lands in its image; define
\(\mathrm{tr}\) by
\(\mathrm{inc}\,\mathrm{tr}=\prod_jE_j\).
It is \(R\)-semilinear through \(\beta\).  HTW9 makes
the formula
\(\phi(z(u)\otimes m)=z(U)\mathrm{inc}(m)\)
descend through the splitting relations.

For a choice \(\mathbf T=(T_j)_j\) put
\(a_{\mathbf T}=\prod_ja_{T_j}(U_j)\) and similarly
for \(b_{\mathbf T}\).  Define the reverse map explicitly:
\[
 \psi(v)=\sum_{\mathbf T}a_{\mathbf T}(u)\otimes
                                     \mathrm{tr}(b_{\mathbf T}(U)v).
 \tag{HTW.19}
\]
The divided-difference pairing is linear over symmetric
polynomials.  HTW10's two dual expansions show
\(\sum a_T\otimes b_T\) is balanced with respect to
multiplication by any variable over the symmetric
coefficient ring: pair its two sides against arbitrary
polynomials by \(\partial\); the two dual expansions
recover equality.  Equivalently the expansions exhibit
this pairing as a perfect finite-projective pairing.
HTW11's formula, with HTW9 for symmetric scalars, gives
the same balancing for (HTW.19).  Thus \(\psi\) is
semilinear in every \(U_{ji}\), as well as in \(R\).

For \(\phi\psi\), substitute (HTW.19), move operators
at distinct places past one another, and use
(HTW.14) at each place.  The result is \(f_0(U)\).
For \(\psi\phi\), use (HTW.15) successively at every
place; (HTW.13) then gives \(f_0(u)\).
The finite product argument also covers \(q=0\),
when the splitting algebra is \(B\), the trace and
inclusion are inverse, and \(f_0=1\).
\end{proof}
% END HTW12

% BEGIN HTW13
\paragraph{HTW13. Enlargement to the literal correction polynomial.}
\label{claim:HTW13}
The maps in HTW12 can be chosen with both compositions
equal to \(F\) in (HTW.2).
\begin{proof}
Each factor \(X_k-Q_jX_i\), \(i<k\), in
\(N_{Q_j}\) is a factor of \(F\), namely the second
factor of its ordered pair \((k,i)\).
Thus \(F=f_0c\) for an integral polynomial \(c\).
Keep \(\phi\) and replace \(\psi\) by
\(\psi\circ c(U)\).  Polynomial semilinearity and
commutation give
\(\phi\psi=F(U)\), \(\psi\phi=F(u)\).
No genericity, localization, or cancellation of \(F\)
is used.  This is precisely the original
\(\mathrm{TwoMaps}\) statement.
\end{proof}
% END HTW13

% BEGIN HTW14
\paragraph{HTW14. The complete original operator-module record.}
\label{claim:HTW14}
An actual \(\mathrm{LevMod}\) with
\(\mathrm{LocalLevelConditions}\) and
\(Q_j\equiv1\bmod\ell^N\) yields the original
\(\mathrm{OpMod}\), together with both fields of
\(\mathrm{OpRel}\) and its two-map field, with exponent \(n!\).
\begin{proof}
Keep the actual coefficient ring \(R\), residual datum,
pair \(p\), coefficient map \(\beta\), and the original
weight-variable identities of \(\mathrm{LevMod}\).
Choose \(M_1,W\), the two polynomial actions, and trace
as in HTW2.  The finite-level module field makes \(M_1\)
finite over \(\Lambda\).  The projectivity field,
transported through the concrete quotient (HTW.4),
gives the required \(\Lambda[\Delta]\)-projectivity.
Formula (HTW.5) gives the exact generator relation
\(\gamma_{ji}=1+Y_{ji}\).  HTW1 supplies \(a_j\ge N\),
trace surjectivity, and its exact kernel.
HTW2 supplies the \(U\)-intertwining.
These are all fields of \(\mathrm{OpMod}\).
HTW13 gives \(\mathrm{OpRel.two}\); HTW8 gives its
two universally quantified trace relations, including
every \(g\) in the original group.  Thus no operator
field remains an additional assumption.
\end{proof}
% END HTW14

\subsection{The completed coefficient ring and free diamond module}

% BEGIN HTW15
\paragraph{HTW15. A canonical simultaneous topologically nilpotent summand.}
\label{claim:HTW15}
On a finite module over complete noetherian local \(\Lambda\),
a finite commuting family \(T_s\) has a direct summand
on which all \(T_s\) are topologically nilpotent.
Its projector commutes with every endomorphism commuting
with the family.
\begin{proof}
Let \(C\) be the commutative \(\Lambda\)-subalgebra generated
by the operators in the finite module
\(\operatorname{End}_\Lambda(M)\).  It is finite over
\(\Lambda\): each generator is integral by the determinant
trick on a finite module, and there are finitely many
generators.  It is complete.  The Artinian algebra
\(C/\mathfrak m_\Lambda C\) is a product of local
Artinian algebras.  Sum the primitive factor idempotents
in which every \(T_s\) is in the maximal ideal.
Idempotents lift uniquely through the complete
\(\mathfrak m_\Lambda\)-adic ring: lift successively
through nilpotent quotients, where the equation
\(e^2-e=0\) has invertible derivative \(2e-1\), and
take the inverse limit.

On the selected summand sufficiently high powers of
each \(T_s\) lie in \(\mathfrak m_\Lambda C\), hence
their powers tend to zero.  On any omitted factor
one \(T_s\) is a unit; if its powers on a vector tend
to zero, multiplication by its inverse and separatedness
of the finite module force that vector to be zero.
The selected range is therefore exactly the
simultaneous topologically nilpotent submodule.
The projector belongs to \(C\), so any endomorphism
centralizing the \(T_s\) centralizes it.  Inclusion
and projection are split \(\Lambda\)-linear maps.
\end{proof}
% END HTW15

% BEGIN HTW16
\paragraph{HTW16. Power-series actions and their semilinear maps.}
\label{claim:HTW16}
On the summand in HTW15 the commuting operators extend
to an action of \(R[[X_s]]\), with \(X_s\) acting as \(T_s\),
when the original operators and the projector commute
with \(R\).  A map between two such finite
\(\Lambda\)-modules that intertwines coefficients and
variables intertwines every power series.
\begin{proof}
Modulo \(\mathfrak m_\Lambda^k\), some power of every
\(T_s\) is zero, uniformly on the finite module.
Since there are finitely many variables, all monomials
of sufficiently large total degree act as zero.
Evaluate a power series by its finite truncation in
this quotient.  These evaluations are compatible as
\(k\) varies, so completeness gives a unique action.
Addition and multiplication can be checked in each
finite truncation, where they are polynomial identities.

A coefficient-semilinear additive map is
\(\Lambda\)-linear when the two coefficient maps agree
on \(\Lambda\), and hence preserves
\(\mathfrak m_\Lambda^k\)-multiples.  In every quotient
its variable-intertwining property extends to the
finite polynomial truncation.  Passing to the inverse
limit proves the full power-series assertion.
Thus no unproved continuity of an arbitrary additive
map is being used.
\end{proof}
% END HTW16

% BEGIN HTW17
\paragraph{HTW17. Freeness over the exact diamond quotient.}
\label{claim:HTW17}
Apply HTW15 to \(M_1\) with the operators
\(U_{ji}-u^0_{ji}\) and \(Y_{ji}\), where \(u^0_{ji}\in O\)
are any chosen lifts.  Let its range be \(M\).
Then for some finite \(r\),
\[
 M\simeq\left(
  \Lambda[[Y_{ji}]]/
       ((1+Y_{ji})^{\ell^{a_j}}-1)_{j,i}
                   \right)^r
 \tag{HTW.20}
\]
as a module for the displayed power-series action.
\begin{proof}
The projector commutes with \(\Delta\), since
\(\gamma_{ji}=1+Y_{ji}\) and all polynomial operators
commute.  Therefore \(M\) is a direct summand of
the projective \(\Lambda[\Delta]\)-module \(M_1\).
It is finite.  The group algebra \(\Lambda[\Delta]\)
is local: modulo \(\mathfrak m_\Lambda\) its augmentation
ideal is nilpotent, because the residue characteristic
is \(\ell\) and each generator has \(\ell\)-power order.
Its only simple quotient is the residue field.
A finite projective module over a local ring is free:
lift a residue-field basis, obtain a surjection from
a finite free module by Nakayama, split it by
projectivity, and use Nakayama again to kill the
complement.

Finally \(Y_{ji}\mapsto\gamma_{ji}-1\) identifies
the ring in (HTW.20) with \(\Lambda[\Delta]\).
Polynomial division by the monic relations gives
the basis \(\prod Y_{ji}^{b_{ji}}\), \(0\le b_{ji}<\ell^{a_j}\).
The relations are distinguished, so the same division
holds for power series over the complete local ring,
by successive reduction modulo powers of its maximal
ideal.  Both sides are finite complete free
\(\Lambda\)-modules with this basis, and the generator
map identifies them.  It respects the action on \(M\)
by HTW16 and the exact generator relation.
\end{proof}
% END HTW17

% BEGIN HTW18
\paragraph{HTW18. The augmentation quotient is the correct primary summand.}
\label{claim:HTW18}
The trace identifies
\[
 X=M/(Y_{ji})M
       \simeq W_{u^0},
 \tag{HTW.21}
\]
where the right side is the simultaneous topologically
nilpotent summand for \(U_{ji}-u^0_{ji}\) on \(W\).
The isomorphism and its split inclusion into \(W\)
intertwine all polynomial \(U\)-operators.
\begin{proof}
HTW1 gives \(W=M_1/(Y_{ji})M_1\).  The projector
in HTW15 commutes with every \(Y_{ji}\), so taking
this quotient splits as the quotient of its range
and the quotient of its complementary range.
On \(W\) all \(Y_{ji}\) are zero.  The image of the
projector is exactly the summand where the remaining
\(U_{ji}-u^0_{ji}\) are topologically nilpotent:
this follows either from its local-factor construction
or from uniqueness of the decomposition in HTW15.
Thus the trace induces the stated isomorphism and
the splitting.  Commutation with the \(U\)-operators
gives all the intertwining statements.  In particular
one has both an inclusion \(j:X\to W\) and a
projection \(\kappa:W\to X\), with \(\kappa j=1\).
\end{proof}
% END HTW18

% BEGIN HTW19
\paragraph{HTW19. Completion of the ordered splitting algebra.}
\label{claim:HTW19}
The source ring \(\mathrm{Sp}(P,\beta u^0)\) is the
factor of the finite ordered splitting algebra \(S\)
on which every \(u_{ji}-\beta(u^0_{ji})\) is
topologically nilpotent.  The maps of HTW13 restrict
to
\[
 \mathrm{Sp}(P,\beta u^0)\otimes_BM_0
          \ \rightleftarrows\ W_{u^0},
 \tag{HTW.22}
\]
with both compositions still equal to the literal \(F\).
\begin{proof}
A monic polynomial can be split successively by
adjoining one root, dividing by its monic linear
factor, and repeating.  Thus \(S\) is finite over
complete local \(B\).  Its quotient modulo
\(\mathfrak m_B\) is Artinian.  Select the factors
where all the ordered root differences have
residue zero, and lift their summed idempotent as
in HTW15.  If no such factor exists the selected
ring is the zero ring, which is allowed in this
description.

The defining power-series quotient for
\(\mathrm{Sp}\) is precisely this factor.  In one
direction, its topologically nilpotent root
differences give an evaluation of the power series
and satisfy every coefficient of the splitting
relations.  In the other, those relations give an
ordered splitting of each \(P_j\).  The universal
maps are inverse on coefficients and root
differences and hence, by completeness, on the
whole rings.  The factor is a finite direct summand
of \(S\).

Tensoring its idempotent with \(M_0\) preserves this
direct-summand description, without any flatness
assumption on \(M_0\).  The two maps of HTW13
intertwine every root difference, so they carry
the chosen summands into one another.  Restriction
of their already proved composition identities
gives (HTW.22).  These modules are finite over
\(\Lambda\): the left side is a quotient of a
finite direct sum of copies of \(M_0\), and the
right side is a summand of the finite quotient \(W\).
Their root-difference actions are topologically
nilpotent also in the \(\Lambda\)-adic topology.
If \(M_0=0\), the left module is zero and this
assertion is immediate there; the right summand
already has its \(\Lambda\)-adic nilpotence by
definition.  Otherwise the image of \(B\) in the endomorphisms
of \(M_0\) is a finite local \(\Lambda\)-algebra:
it is a \(\Lambda\)-submodule of a finite
endomorphism module and a quotient of local \(B\).
Its maximal ideal is the radical of
\(\mathfrak m_\Lambda\) times that image.
Consequently powers of \(\mathfrak m_B\)
act through arbitrarily high powers of
\(\mathfrak m_\Lambda\).  The same argument
on the finite splitting-algebra tensor module
converts nilpotence of each root difference
modulo \(\mathfrak m_B\) into the required
\(\Lambda\)-adic topological nilpotence.
\end{proof}
% END HTW19

% BEGIN HTW20
\paragraph{HTW20. The completed ring and the exact coefficient-map identity.}
\label{claim:HTW20}
From the \(\mathrm{OpMod}\) of HTW14 define
\[
 R'=R[[X_{ji},Y_{ji}]],\qquad U'_{ji}=u^0_{ji}+X_{ji}.
 \tag{HTW.23}
\]
Extend its pair by the constant coefficient inclusion.
There are original maps
\[
 s:\Lambda[[Y_{ji}]]\longrightarrow R',
 \quad
 \theta:R'\longrightarrow\mathrm{Sp}(P,\beta u^0)
\]
with the original weight and pair identities, and
\[
 \theta(U'_{ji})=u_{ji},\qquad
 \theta(s(z))=
       B\longrightarrow\mathrm{Sp}\bigl(\lambda(z|_{Y=0})\bigr).
 \tag{HTW.24}
\]
\begin{proof}
First \(\beta\) agrees with \(\lambda\) on
\(\Lambda\).  On a weight variable both values are
the prescribed scalar times the corresponding
ordinary character minus \(1\), by the
\(\mathrm{OpMod}\) weight and pair fields.
They agree on \(O\).  These are local maps of
complete coefficient rings, hence adically
continuous, so their agreement on polynomials
gives agreement on power series.

Define \(s\) by the original weight map on
coefficients and by the right-hand \(Y\)-variables.
Define \(\theta\) on \(R\) through \(\beta\),
on \(X_{ji}\) by \(u_{ji}-\beta(u^0_{ji})\),
and on every \(Y_{ji}\) by zero.  These variable
images are topologically nilpotent, so the
complete power-series maps exist.  Formula
(HTW.24) holds on coefficients and variables,
hence on all power series.

Constant base change preserves every field of
the full pair, and \(\beta.\mathrm{pair}\) gives
the full determinant and ordinary-character
identity after \(\theta\).  This is equality of
polynomial laws over every scalar algebra.
The prescribed weight equality for \(s\) follows
from that for \(R\), and
\(U'_{ji}-u^0_{ji}=X_{ji}\in\mathfrak m_{R'}\).
The ring \(R'\) is still an original complete
noetherian local coefficient ring with the same
residue field.
\end{proof}
% END HTW20

% BEGIN HTW21
\paragraph{HTW21. All fields of TWRing and TWMaps.}
\label{claim:HTW21}
HTW14 produces the literal \(\mathrm{TWRing}\)
and \(\mathrm{TWMaps}\) with exponent \(n!\), for
any specified \(u^0\), with the same \(b,J,N\).
\begin{proof}
Use the ring, pair, map \(s\), and roots in HTW20.
Let \(R'\) act on the module \(M\) in HTW17:
the left \(X\)-variable acts by \(U-u^0\) and
the right \(Y\)-variable by the original \(Y\).
HTW16 constructs this action.  The isomorphism
(HTW.20), with the same \(a_j\ge N\), is exactly
the required free diamond-quotient model for
\(\mathrm{TWRing}\).

Compose the two maps of (HTW.22) with the inverse
isomorphisms in (HTW.21).  Polynomial
semilinearity holds by construction.  HTW16
extends it to every \(t\in R'\), so these are
the exact \(\phi,\psi\) fields of \(\mathrm{TWMaps}\).
Their compositions are \(F(U')\) and its
\(\theta\)-image, without removing any factor.
HTW20 supplies \(\theta.\mathrm{pair}\),
\(\theta(U')\), and the coefficient-map identity.

Finally polynomial evaluation
\(U\mapsto u^0+X,\ Y\mapsto Y\) carries the
trace polynomials in (HTW.10) to the corresponding
trace polynomials of the constantly extended
pair.  This follows from full polynomial-law
naturality, including its trace coefficient.
Inclusion \(M\hookrightarrow M_1\) is injective
and intertwines this evaluation.  Thus both
\(n!\)-th power relations hold on every \(m\in M\),
for every \(j,i,g\).  This checks the last two
fields, completing every field of both records.
\end{proof}
% END HTW21

% BEGIN HTW22
\paragraph{HTW22. Transport through the original group isomorphism.}
\label{claim:HTW22}
The \(\mathrm{OpMod}\) and all its relations transport
through the original injective map \(e:G\to G'\)
of global data, when its map to the fixed base
commutes.  The coefficient rings, modules, roots,
base \(b\), ideal \(J\), and exponent do not change.
\begin{proof}
Such a global-data map is already surjective by
its defining fields.  With the stated injectivity,
it is an isomorphism; a continuous bijection
between these compact Hausdorff groups has a
continuous inverse.  Use that inverse \(k\) to
pull the residual datum and full determinant law
to \(G'\).  On the local groups the structural
maps intertwine the prescribed characters,
polarization, ordinary relations and type kernels;
thus the pulled object is the original pair.
For the type kernels specifically, a point lies
in the transported inertia image exactly when
its inverse image lies in the old one, and each
residual-character kernel is preserved by the
same character equality.  This checks their
intersection, not just the ambient inertia group.

The base compatibility says \(h\circ k=h'\);
hence the old coefficient map still sends the
pair to the fixed base pair.  All weight
characters remain the same.  In the trace
relations replace \(g'\) by \(k(g')\);
the group-algebra map sends \([e(f_j)]\) to
\([f_j]\) and similarly for \(t_j\).
Polynomial-law pullback commutes with coefficient
base change and trace, so both relations are
exactly the old ones.  All module and two-map
fields are unchanged.  This proves the literal
\(\mathrm{opTransport}\) consumer.
\end{proof}
% END HTW22

\subsection{One base and one exponent for the entire auxiliary family}

% BEGIN HTW23
\paragraph{HTW23. The full original HasAutOfAlt implication.}
\label{claim:HTW23}
Under the original \(\mathrm{LocPAlt}\), if
\(V_0.S'=s.S\) and \(\mathrm{LevelAlt}(B,J)\),
then \(\mathrm{HasAut}(B,J)\).  The same one
base \(b\), one residual-root lift \(u^0\), and
one exponent \(n!\) work for every auxiliary
integer \(N\).
\begin{proof}
Choose \(b\) once from \(\mathrm{LevelAlt}\),
before any auxiliary level.  Now fix an arbitrary
original Taylor--Wiles group family \(T\) and its
place realization \(Q\).  For each \((j,i)\)
choose \(u^0_{ji}\in O\) lifting the residual
character value \(\overline\chi_i(T.sg_j)\);
the residue map is surjective.  This choice
precedes \(N\).

For each \(N\), use the given \(Q.V_N\),
places \(Q.v_N,Q.w_N\), and prolongations \(Q.B_N\).
Their distinctness, avoidance, exact set of
auxiliary places, splitting, and compatibility
with the base map are fields of the same
original realization.  Thus the single
\(\mathrm{LevelAlt}\) statement supplies the
actual \(\mathrm{LevMod}\) and its separating
classical family at that level.  Avoidance of
\(s.S\), which contains the places above \(\ell\),
shows the auxiliary \(w_j\) are away from
\(\ell\), by restriction of their prolongations.

Apply HTW6 to the given \(Q.fr_N,Q.\tau_N\).
Its cardinality equality identifies the actual
\(Q_j\) with \(T.Nq_N\), and its comparison uses
the full \(\mathrm{LevPtAlt.lg}\) of the same
classical points.  The congruence field of \(T\)
then gives the hypothesis of HTW14.
Transport its operator module along the actual
\(Q.e_N\) using HTW22.  The Frobenius and tame
compatibility fields of \(Q\) identify the
transported elements with \(T.ph_N,T.ta_N\).
Apply HTW21 with the already chosen \(u^0\).

Choose these resulting rings and maps for every
\(N\).  Together with \(b,u^0\), their residual
equalities and the fixed exponent \(n!\), they
are exactly \(\mathrm{TWAut}(T)\).  Since \(T,Q\)
were arbitrary, this is the original
\(\mathrm{HasAut}\).  Empty auxiliary families
are included by the empty products already
checked above.  The abstract implication does
not impose a new lower bound on \(n\); its
zero-rank matrix and trace assertions were
included explicitly.
\end{proof}
% END HTW23

% BEGIN HTW24
\paragraph{HTW24. Actual same-ideal automorphy from the Hida level construction.}
\label{claim:HTW24}
For every original \(n\ge1\) datum covered by HLM26,
its actual \(\mathrm{IsJ}(B,K,J)\), and the original
local published hypotheses, the full
\(\mathrm{HasAut}(B,J)\) holds with exactly that ideal.
\begin{proof}
HLM26 proves the literal all-level
\(\mathrm{LevelAlt}\) producer from the supplied
unitary Hecke datum and \(\mathrm{IsJ}\).  In
particular it constructs one \(b\) satisfying
(HTW.1) before all \(q,V,h,v,w,B_q\), rather
than producing an unrelated base at each level.
Apply HTW23 to this output and to
\(V_0.S'=s.S\).  Every choice in that proof
keeps \(b\) and hence the same support ideal.
The resulting \(\mathrm{HasAut}\) is the full
original record with both trace relations,
the exact coefficient map and correction
polynomial, for every original auxiliary
family.  Nothing excludes \(\ell=3\).

This completes the level-to-automorphy producer.
It does not itself install both rank-four
\(\mathrm{Link}\) objects or identify the
entire selected \(\mathrm{COUNT\_\allowbreak II\_\allowbreak SUPPORT}\)
contract with this one endpoint.  Those wider
contracts retain their other original fields.
\end{proof}
% END HTW24

\subsection{Scope of this completion}

The three published fields of \(\mathrm{LocPAlt}\)
remain explicitly retained at their original scope.
The proof above supplies the intervening finite
Hecke, matrix, determinant-law and complete-ring
arguments and constructs the full output record.
It preserves the fixed coefficient DVR, the
given complex comparison, the auxiliary-place
realization, the complete ordinary pair, the
same \(J\), the exact ordered correction factor,
and all integer \(N\).  In particular the
nilpotence exponent is uniform because it
comes from one local Jacquet dimension bound,
while the existence of the free module comes
from projectivity over the finite diamond
group algebra, not from dividing by its order.


\section{One Hida ideal and the complete selected Count-II support input}
\label{sec:hida-installation-v45}

This chapter assembles the original objects, after their producers
have been reconstructed.  It has two different scopes.  For every
original positive rank it constructs the whole
\(\mathrm{HidaUAlt}\): one ideal with all six Hida properties
and the entire same-ideal Taylor--Wiles automorphy system.  For
each of the two rank-four Links used in Count II it also supplies
the original \(\mathrm{P1core}\) on that ideal.  The latter uses
the completed one-point support proof.  The general assertion
\(\mathrm{PfromAut}\), with arbitrary chains of
\(\mathrm{NiceE}\) points and its \(\mathrm{P3core}\) conclusion,
is a separate antecedent in the all-rank modularity implication
below.  It is not inferred from the one-point argument.

Retain the original \(H'=\mathrm{HypothesesH}'\), including its
named automorphic, local-correspondence and cohomological
foundations.  In particular the original definite-inner-form
base change, cuspidal Galois attachment and printed local
classification inputs retain their previous status.  The
cohomological package \(\mathrm{CohP}\) means all eight original
fields: local finiteness, local Tate duality, local Euler
characteristic, global finiteness, global Euler characteristic,
Greenberg--Wiles, normalized Poitou--Tate, and the original tame
local theorem.  The full \(\mathrm{LocCFT}\) and the actual
Artin normalizations are retained, with their already reconstructed
finite consumers.  No general support or automorphy assertion is
added to these foundations.

For the all-rank statements fix the original
\(\ell,L,\Omega,s,V,c,n,r,B,\Lambda,\iota\), where
\(n\geq1\), \(O=c.O\), \(V.S'=s.S\), and
\(\iota:\overline{\mathbb Q}_\ell\simeq\mathbb C\) is the
given field isomorphism.  All coefficient maps, residual labels,
decomposition groups and local character orders are the original
ones.  There is no continuity requirement on \(\iota\).
The original setting already has \(\ell\ne2\); the prime three
remains included throughout.

% BEGIN HIA1
\paragraph{HIA1. All six Hida fields on the same given ideal.}
\label{claim:HIA1}
Let \(R\) be any of the original unitary Hecke data, let \(K\)
be its given auxiliary level, and suppose
\(R.\mathrm{IsJ}(B,K,J)\).  Then, under the retained original
inputs and \(V.S'=s.S\), the literal record
\[
                         B.\mathrm{HidaPropsAlt}(\iota,J,\Lambda)
                         \tag{HIA.1}
\]
holds for this very ideal \(J\).
\begin{proof}
Use the following six proved endpoints with the same original
arguments and the same ideal:
\[
\begin{array}{c|c|l}
\text{field}&\text{proved endpoint}&\text{literal content}\\ \hline
\mathrm{ring}&\mathrm{JRP15}&\mathrm{RingProps}\ \Lambda\ J\\
\mathrm{dense}&\mathrm{UPA18}&\mathrm{Dense}\ J\\
\mathrm{pt}&\mathrm{HPT18}&\mathrm{PointTheoremBaseObject}\ \iota\ J\\
\mathrm{end1}&\mathrm{HED16}&\mathrm{ArithmeticPointDecomposition}\ \iota\ J\\
\mathrm{end2}&\mathrm{HAW19}&\mathrm{IrrArithPointsAutomorphicAway}\ \iota\ J\\
\mathrm{twist}&\mathrm{HTS8}&\mathrm{TwistOk}\ J.
\end{array}
\tag{HIA.2}
\]
The original aliases \(\mathrm{arithmeticPointPurityDense}\)
and \(\mathrm{arithmeticPointClassicalityDense}\) are exactly
the fourth and fifth contents of this table.  No implication
from a weaker purity or classicality predicate is being used.

There are several substantive guards inside these fields.
The point theorem quantifies over a supplied representation
and proves \(J\subseteq\ker x\) under its original conditions
\(2\leq n\leq4\), continuity, automorphy, full auxiliary
inertia triviality, and the prescribed ordered flags and
nonlinkage.  It is an implication with those rank guards even
when the ambient rank is outside that interval.  The first
end field retains every positive summand rank and positive
string length, their exact total rank, the full group
characteristic-polynomial identity, and its strict cofinite
root-ratio assertion over all original places under the original
condition that every string length is one.  The second
retains \(\mathrm{AutCsd0}\) without a continuity antecedent;
only its \(\mathrm{WitOff}\) conclusion is conditional on
continuity.  Its one real window center is chosen before all
places for each attached cusp.  The twist field is the
original finite \(\ell\)-power-order twist assertion, on this same ideal.

The ring field retains the original \(r.\mathrm{V0ok}\)
antecedent in its auxiliary-unramified clause; that antecedent
is not imposed on the entire six-field record.  HPT18 and
HAW19 likewise add no such hypothesis to their endpoints.
Every endpoint in (HIA.2) accepts the given
\(\mathrm{IsJ}\), rather than selecting an unrelated ideal.
Thus the six proofs are the six fields of (HIA.1).

This also includes the legitimate zero-module branch
\(J=B\).  Fields whose antecedents contain
\(J\subseteq\ker x\) then have no such point.  The
\(\mathrm{pt}\) field does not have that antecedent and is
not declared vacuous: its full conclusion is precisely the
one supplied by HPT18 for this \(\mathrm{IsJ}\) branch.
\end{proof}
% END HIA1

% BEGIN HIA2
\paragraph{HIA2. The full all-rank Hida and automorphy supplier.}
\label{claim:HIA2}
Under the original \(H'\), the entire \(\mathrm{HidaUAlt}\)
holds.  In particular, for every original datum in its definition,
including \(V.S'=s.S\), \(\mathrm{noMu}\),
\(\mathrm{Large}(L,c)\), every \(n\geq1\), and every
actual \(r,B,\Lambda,\iota\), there is one ideal \(J\subseteq B\)
such that
\[
 B.\mathrm{HidaPropsAlt}(\iota,J,\Lambda)
                   \quad\hbox{and}\quad \mathrm{HasAut}(B,J).
 \tag{HIA.3}
\]
No further coefficient extension is required.
\begin{proof}
Take \(Q=\mathrm{QData.empty}(V,S'=S)\).  ITA13 gives
the actual eight-field unitary Hecke datum \(R\) over the
specified \(O\) and \(\Lambda\).  There is a unique tuple
of auxiliary levels indexed by its empty set; use it as \(K\).
Put
\[
 H=\mathrm{hset}(Q,R.\tau_0),\qquad
 M=H.\mathrm{MK}(R.\theta,R.a,K),\qquad
 \mathbb T=H.\mathrm{TK}(R.\theta,R.a,K).
 \tag{HIA.4}
\]
HLM9 supplies the literal \(\mathrm{IsJ}\) on these objects
for every rank in the statement.  If \(M=0\), it takes
\(J=B\).  Otherwise it constructs the continuous original
determinant and local characters on \(\mathbb T\), proves
their full pair relations and their prescribed weight map,
and obtains a surjection \(f:B\to\mathbb T\).  It takes
\(J=\ker f\).  Since \(\mathbb T\) is its actual subalgebra
of \(\operatorname{End}_{\Lambda}(M)\), this is also
\(\operatorname{Ann}_B(M)\).  Thus the support ideal has been
chosen before any later arithmetic point or auxiliary
Taylor--Wiles family.  HIA1 gives all six fields on it.

JRP1 verifies \(\mathrm{IsCoef}(\ell,O)\): the given
coefficient ring is a complete discrete valuation ring with
finite residue field of characteristic \(\ell\).
The original local package used in converting levels to
automorphy is exactly
\[
 \mathrm{LocPAlt}=
 (\text{Jacquet finite length/admissibility},
   \text{geometric lemma},\text{tame local theorem}).
 \tag{HIA.5}
\]
Its first two inputs are the original literature fields
\(H'.\mathrm{litN.loc.bz}\) and
\(H'.\mathrm{litN.loc.geo}\).  Its third is the original
tame-local consequence of the retained local Tate duality
field.  These are the same named foundations as in HTW24;
none is inferred merely from the six Hida fields.

HLM26 constructs the full \(\mathrm{LevelAlt}(B,J)\) with
one base \(b\) before every auxiliary datum.  HTW24 converts
that very family, with (HIA.5), into the whole
\(\mathrm{HasAut}(B,J)\).  In expanded form this says
\[
 \begin{gathered}
 \forall\mathcal T:\mathrm{TWGrp}(\mathrm{globOf}(r,V)),
 \quad\mathrm{TWPl}(r,\mathcal T)
 \ \Longrightarrow\\
 \mathrm{Nonempty}\!\left(
 \mathrm{TWAut}(\mathrm{pairOf}(B.u),J,\mathrm{Empty},
                   n[L^+:\mathbb Q],\mathcal T)\right).
 \end{gathered}
 \tag{HIA.6}
\]
In each system the residual eigenvalue lifts are chosen before
all its levels; the exponent is the original \(n!\), uniform
in the level, and both original trace-polynomial relations
and the coefficient-map correction are retained.  No
division by \(n!\) occurs.  The zero-module branch is the
one already covered by the full level and automorphy proofs.

The choices above use arbitrary supplied \(B\) and
\(\Lambda\); universality is not restricted to one preferred
representative.  They take place over the original \(O\),
and therefore work for every coefficient object in the
definition, including every sufficiently large later object.
Together with (HIA.1), (HIA.6) proves exactly (HIA.3) and
the full quantifier order of \(\mathrm{HidaUAlt}\).
\end{proof}
% END HIA2

% BEGIN HIA3
\paragraph{HIA3. The cyclotomic condition and the separate modularity implication.}
\label{claim:HIA3}
For every original level \(V\), its \(\mathrm{noMu}\)
condition implies the actual residual cyclotomic condition
\[
 \exists g\in V.\mathrm{frame}.G_0,
                 \overline\varepsilon(g)\ne1.
 \tag{HIA.7}
\]
Independently of the selected rank-four support proof below,
the original all-rank implication is
\[
 \mathrm{CohP}+\mathrm{LocCFTAll}+\mathrm{PfromAut}
              +\mathrm{HidaUAlt}
                       \ \Longrightarrow\ \mathrm{ModUAlt}.
 \tag{HIA.8}
\]
Here \(\mathrm{PfromAut}\) is an explicit antecedent of
(HIA.8), not a conclusion of this chapter.
\begin{proof}
Choose a primitive \(\ell\)th root \(\zeta\) in the
given algebraic closure.  If every element of \(G_L\)
fixed \(\zeta\), the fixed-field theorem would put it
in \(L\), contrary to \(\mathrm{noMu}\).  Thus some
\(\sigma\in G_L\) moves it.  The action on these roots
is multiplication by the reduction of the cyclotomic
character, so \(\overline\varepsilon(\sigma)\ne1\).
By the defining cyclotomic compatibility of the original
level projection, the element \(g=V.\mathrm{pr}(\sigma)\)
lies in the original \(G_0\) and has that same reduction.
This proves (HIA.7).  It does not construct a new level
quotient or replace its cyclotomic character.

To prove (HIA.8), fix any input in the literal definition
of \(\mathrm{ModUAlt}\).  The \(\mathrm{HidaUAlt}\)
antecedent supplies one \(J\) satisfying (HIA.3); HIA2
has now supplied that antecedent under the original \(H'\).
Use JRP1 for \(\mathrm{IsCoef}\), and use
\(\mathrm{LocCFTAll}\) on this \(s\).
The given \(n\geq1\), \(V.S'=s.S\), (HIA.7), the full
\(\mathrm{CohP}\), and \(\mathrm{HasAut}(B,J)\) are
precisely the antecedents of the original
\(\mathrm{PfromAut}\).  That hypothesis returns
\[
                       B.\mathrm{P1core}(J)
                          \quad\hbox{and}\quad
                       B.\mathrm{P3core}(J).
 \tag{HIA.9}
\]
Appending these two fields to \(\mathrm{HidaPropsAlt}\)
is exactly \(B.\mathrm{ModularAlt}(\iota,J,\Lambda)\).
This proves (HIA.8) at all ranks and on the same ideal.

For clarity about its boundary, \(\mathrm{P3core}\)
quantifies over every finite chain of \(\mathrm{NiceE}\)
points with the original thick connecting primes.  The
one-point \(\mathrm{NiceI}\) argument that follows has
no such connecting primes and supplies only its stated
\(\mathrm{P1core}\) conclusion.  It does not discharge
the \(\mathrm{PfromAut}\) antecedent of (HIA.8).
\end{proof}
% END HIA3

% BEGIN HIA4
\paragraph{HIA4. The selected rank-four support conclusion on that ideal.}
\label{claim:HIA4}
Fix either original rank-four Link, its actual \(B,\Lambda\),
and the ideal \(J\) obtained in HIA2.  Retain the full
original \(\mathrm{CohP}\) and local reciprocity hypotheses.
Then the entire original one-point predicate
\(B.\mathrm{P1core}(J)\) holds:
\[
 \begin{gathered}
 x.\mathrm{NiceI},\quad B.\mathrm{IsPt}(x,f),\quad
 J\subseteq\ker f,\quad
 \mathfrak q\in\operatorname{Spec}B,\quad
 \mathfrak q\subseteq\ker f\\
                       \Longrightarrow\quad J\subseteq\mathfrak q.
 \end{gathered}
 \tag{HIA.10}
\]
All the maps, primes and point fields in (HIA.10) are the
original ones.  This conclusion requires no general
\(\mathrm{PfromAut}\) or \(\mathrm{P3core}\).
\begin{proof}
We verify the completed inputs to JSP11 in their actual order.
JRP1 supplies the coefficient condition.  KX3 supplies the
whole \(\mathrm{LocCFT}\) on the chosen decomposition
groups; CIP12 supplies the exact additional Artin matching
and the already reconstructed normalization consumer.
The Link has \(V.S'=s.S\) and \(\mathrm{noMu}\), so the
unconditional first assertion of HIA3 supplies (HIA.7).
The required \(\mathrm{CohP}\) is the retained eight-field
package, and HIA2 supplies \(\mathrm{HasAut}(B,J)\)
before \(x,f,\mathfrak q\) are introduced.

Fix arbitrary \(x,f,\mathfrak q\) as in (HIA.10).
CIP10 and CIP11 normalize this \(\mathrm{NiceI}\) point
on the actual updated inertia frame.  They give the original
\(a=4\), \(N3\), and \(N4(D_0)\), where \(D_0\)
is the original auxiliary subgroup, together with an
injective local map of point rings and the full determinant
and ordered-character identities.  The first part of CAP13
supplies the required \(\Phi\)-finiteness before a finite
list is selected.  CCH15 then supplies that list, its second
point extension and its common Kap with \(\kappa.g=0\).
CCH13 preserves the original \(N4(D_0)\), the complete
determinant law and the ordered characters.  If \(j\)
denotes the composite point-ring map, it is injective, has
the original compatible coefficient structure, and gives
\[
                             g=jf,\qquad\ker g=\ker f.
 \tag{HIA.11}
\]
Only these pointwise finite extensions are made; the outer
coefficient object and its earlier threshold stay fixed.

Now use the second part of CAP13 to realize precisely the
selected list by the actual \(\mathcal T\) and
\(\mathrm{TWPl}\).  The residual-kernel identity of this
record makes every residual character one on its subgroup
\(H\).  The full determinant identity therefore gives
\(\operatorname{cp}(\overline{\rho(h)})=(X-1)^4\)
for \(h\in H\).  JSP11's Cayley--Hamilton and
characteristic-\(\ell\) argument applies on this same
point: \((\rho(h)-1)^4\in tM_4(k[[t]])\), whence
\(\rho(h)^{\ell^{4m}}-1\in t^mM_4(k[[t]])\).
Taking the open normal cores and then every \(m\)
proves the exact \(\mathrm{PL}(\ell,H)\)-kernel
condition.  Thus the later bounds receive the full pro-\(\ell\)
condition, not only a residual statement.

CFO20 and CFO21 now give the required first-order field
on this very point and Kap, with parameter
\[
                            4[L^+:\mathbb Q]+4\mathcal T.q.
 \tag{HIA.12}
\]
Its constant is chosen after the point and list, before
the truncation and eventual level.  CCH16 gives the
obstruction field on the same Kap with constant zero;
CCH17 checks their coupled installation.  No independent
Kap or alternative point is used for the two bounds.
Together with the original eigenvalue and detector data
these construct the literal \(\mathrm{Pt}\) required
by the support consumer.

JSP11 supplies the actual \(\mathrm{Adapted}\),
\(\mathrm{GenBy}\), stable-discriminant and kernel
comparisons.  Its use of \(\mathrm{HasAut}\) is now
fed by HIA2 at this \(\mathcal T/\mathrm{TWPl}\), with
the same \(J\), the empty parameter type and
dimension \(4[L^+:\mathbb Q]\).  SCI13 supplies the
full selected \(\mathrm{Raw}\) conclusion for exactly
\(h=0,a=n=4,\kappa.g=0\), including its actual
quotient, coefficient normalization and prime contraction.
Apply JSP11 with these supplied inputs and initial stable
prime \(\ker g\).  Its conclusion is
\(J\subseteq\mathfrak q\), after (HIA.11) identifies
the original point kernel.  This proves (HIA.10) for
every original test.

All the hypotheses formerly left in the selected route
have now been supplied: normalization, choice, actual
auxiliary primes, the coupled first- and second-order
bounds, selected Raw, and same-ideal HasAut.  The arithmetic
and cohomological hypotheses printed above are unchanged.
The proof does not invert \(24\), exclude \(\ell=3\),
or assert the support theorem for arbitrary chains.
\end{proof}
% END HIA4

% BEGIN HIA5
\paragraph{HIA5. Two excluded places give one global window.}
\label{claim:HIA5}
Let \(n\geq1\), let \(u\ne u'\) be finite places of the
maximal real subfield, and use the same \(\iota\) and
the same representation \(\rho\).  Then the exact
original implication is
\[
 \begin{split}
 &\mathrm{AutOffN}(u,\iota,\rho)\ \wedge\
       \mathrm{AutOffN}(u',\iota,\rho)\\
 &\hspace{12mm}\Longrightarrow\
 \mathrm{AutCsd0}(\iota,\rho)\ \wedge\
       \bigl(\mathrm{Continuous}(\rho)\Longrightarrow
                                  \mathrm{WitA}(\iota,\rho)\bigr).
 \end{split}
 \tag{HIA.13}
\]
In particular its global window has one real center before
all eligible finite places; the two input centers need not agree.
\begin{proof}
Keep the \(\mathrm{AutCsd0}\) part of the first input.
Only for the remaining conditional assertion assume
\(\rho\) continuous.  The two inputs give
\(\mathrm{WitOff}(u,\iota,\rho)\) and
\(\mathrm{WitOff}(u',\iota,\rho)\).  Either gives the
first part of \(\mathrm{WitA}\): at every place above
\(\ell\), \(\mathrm{Unl1}\) implies
\(\mathrm{PotUnl}\), with all original prolongations
and positive powers.

Fix any regular algebraic cusp \(\pi\) attached to
this \(\rho\).  HAW6 supplies a single real \(t\),
before any place is considered, such that toral unlinked
support at a place implies \(\mathrm{WinAt}(\pi,t,v)\).
The first and second \(\mathrm{WitOff}\) also supply
centers \(a\) and \(b\) for their respective excluded-place
conditions.  Take an arbitrary eligible place \(v\), meaning
\(v\nmid\ell\) or \(\mathrm{PotUnl}(\rho,v)\).
If its underlying finite place in the maximal real subfield
is different from \(u\), use the first window; otherwise
that place is different from \(u'\),
so use the second.  In either case the resulting window
gives toral support and, for every support character,
\[
                    q_v^{c-1/2}<|\xi(\varpi_v)|<q_v^{c+1/2}
                     \quad(c=a\text{ or }b).
 \tag{HIA.14}
\]
If two such characters had unramified quotient taking the
value \(q_v\) on \(\varpi_v\), their absolute values
would differ by the factor \(q_v\).  The smaller strict
lower bound in (HIA.14) would then put the larger one
strictly above its upper bound.  Thus the support is
unlinked in the exact original sense.  Repeated equal
characters cause no difficulty.  Apply the already chosen
\(t\) from HAW6 to get \(\mathrm{WinAt}(\pi,t,v)\).
This works for every eligible \(v\) without a new choice
of \(t\).  It proves the second part of \(\mathrm{WitA}\)
for every attached cusp, and hence (HIA.13).
\end{proof}
% END HIA5

% BEGIN HIA6
\paragraph{HIA6. Both original Links above one common coefficient threshold.}
\label{claim:HIA6}
Retain the original Fit/Deep tower, the original named
foundations, and a fixed \(\iota\).  The entire selected
\(\mathrm{COUNT\_\allowbreak II\_\allowbreak SUPPORT}\) input is now supplied.
More precisely, there is one original coefficient threshold
\(c_1\), chosen before every later \(c'\), Space,
universal ring and weight object, such that both original
rank-four Links have the following property at every
allowed \(c'\): for every actual \(B,\Lambda\) there
is one \(J\subseteq B\) carrying all six Hida fields
and the full \(\mathrm{P1core}\).  The Links retain
their common \(\kappa,\beta,t\), distinct auxiliary
primes, and every original local and residual field.
\begin{proof}
Fix the tower \(D,T\) and \(\iota\) first.  QX4 constructs
its characters before choosing an auxiliary prime, and
chooses \(c_1\) containing their required finite coefficient
values and all the original embeddings of \(\mathcal O_L\).
It follows that for every \(c'\) containing both the original
\(c\) and \(c_1\), the literal Links
\(\mathcal L_0,\mathcal L_1\) exist over \(c'\), with
the same \(\kappa,\beta,t\) and with distinct auxiliary
rational primes.  This coefficient choice precedes all
Spaces and all their universal data.

Each Link includes exactly the hypotheses needed here:
\(V.S'=s.S\), \(\mathrm{noMu}\),
\(\mathrm{Large}(L,c')\), the actual residual rank-four
object, and \(r.\mathrm{V0ok}\), as well as the original
local-place identification and character equalities.
For either Link and arbitrary supplied \(B,\Lambda\),
apply HIA2 over that same \(c'.O\).  It gives a single
ideal with all six fields and full HasAut.  HIA4 adds
\(\mathrm{P1core}\) to that exact ideal.  The cohomological
input is the retained full \(\mathrm{CohP}\); the local
class-field and cyclotomic conditions have been checked
in HIA4.  The type, ordinary and arithmetic-point
constructions introduce no coefficient enlargement.
The normalization and support extensions occur only inside
the later \(\mathrm{P1core}\) point quantifier.  Hence no
new external coefficient threshold is needed beyond QX4's.

The resulting quantifier order is therefore
\[
 \begin{gathered}
 D,T,\iota\ ;\quad \exists c_1\ ;\quad
 \forall c'\supset c,c_1\ ;\quad
 \exists\mathcal L_0,\mathcal L_1\ ;\\
 \forall i\in\{0,1\}\ \forall B_i,\Lambda_i\ ;\quad
 \exists J_i\ ;\quad
 B_i.\mathrm{HidaPropsAlt}(\iota,J_i,\Lambda_i)
                    \ \wedge\ B_i.\mathrm{P1core}(J_i).
 \end{gathered}
 \tag{HIA.15}
\]
It remains valid for every later original Space.  Each
\(J_i\) is chosen before the point, prime, detector,
truncation and level tests in its properties.  The two
ideals belong to their two respective original universal
rings; no equality between ideals in different rings is
asserted.  Within each ring all properties have one witness.
This is exactly the same-ideal and two-Link contract isolated
by JSP12, now with its previously retained producers supplied.

To exhibit its original consumer, take an arbitrary original
Space \(X\) above this threshold.  QX5 constructs the
actual nonempty ordinary pairs, so KX9 supplies actual
universal rings, KX5 supplies their weight objects, and
QX5 gives the maps \(\phi_i:B_i\to X.B\) for \(a=0\).
Use (HIA.15) and set
\[
                    I=\phi_0(J_0)X.B+\phi_1(J_1)X.B.
 \tag{HIA.16}
\]
QX10 requires precisely the six Hida fields and selected
\(\mathrm{P1core}\) now proved, rather than an arbitrary-chain
\(\mathrm{P3core}\).  It proves all five original fields
\(\mathrm{fin},\mathrm{p1},\mathrm{seed},\mathrm{pur},\mathrm{aut}\)
of \(\mathrm{PotentialModularityTower}\) on this same
\(X,I\), and thus the entire \(\mathrm{PMExAlt}\).
The already proved YX2 implication then gives the original
\(\mathrm{CountII}\) under its unchanged background inputs.

The stronger end2 fields also preserve the original two-place
gluing step explicitly.  Whenever the source consumer gives
\(\mathrm{AutOffN}\) on both Links, QX5's comparison and
their common \(\kappa\) identify its input as the same
twisted representation.  Distinct auxiliary rational primes
give distinct auxiliary finite places of the maximal real subfield:
equality of the latter
would make their absolute norms, and hence their degree-one
rational primes, equal.  HIA5 therefore gives the full
global window with one center.  This observation neither
identifies unrelated representations nor requires the two
original window centers to coincide.

All these conclusions retain \(\ell=3\), the fixed
\(\iota\), full determinant laws and pro-\(\ell\) conditions,
and the original cohomological and literature boundaries.
They close this selected support producer.  They do not
prove the general \(\mathrm{PfromAut}\) antecedent in
HIA3 or remove the other independent obligations in the
complete RH and GRH reconstruction.
\end{proof}
% END HIA6

\endgroup
\end{document}
